European option contracts specify an expiration date, and if the option is to be exercised at all, the exercise must occur on the expiration date. An option whose owner can choose to exercise at any time up to and including the expiration date is called \textit{American}.

Whereas a European option can be exercised only at a fixed date, an American option can be exercised any time up to its expiration. The value of an American option is the value achieved by exercising optimally. Finding this value entails finding the optimal exercise rule by solving an optimal stopping problem and computing the expected discounted payoff of the option under this rule. The embedded optimisation problem makes this a difficult problem for simulation.

Because an American option can be exercised at any time prior to its expiration, it can never be worth less than the payoff associated with immediate exercise. This is called the \textit{intrinsic value} of the option.

To price an American option, just as with a European option, we could imagine selling the option in exchange for some initial capital and then consider how to use this capital to hedge the short position in the option. In this case, we would need to be ready to pay off the option at all times prior to the expiration date because we do not know when it will be exercised. We could determine when, from our point of view, is the worst time for the owner to exercise the option. From the owner's point of view, this would be the \textit{optimal exercise time}, and we shall call it that. We could then compute the initial capital we need in order to be hedged against exercise at the optimal exercise time. Finally, we could show how to invest this capital so that we are hedged even if the owner exercises at a nonoptimal time. In the subsequent sections, we do all these things but begin the analysis at a different point than for European options. We define the price of American options using a risk-neutral pricing formula and then show that this price is the smallest initial capital that permits construction of the hedge just described.

\subsection{Perpetual American Put}

The simplest interesting American option is the \textit{perpetual American put}. It is interesting because the optimal exercise policy is not obvious, and it is simple because this policy can be determined explicitly. Although this is not a traded option, we begin our discussion with it in order to present in a simple context the ideas behind the subsequent analysis of more realistic options.
The underlying asset (except where the asset pays dividends, which will be explicitly mentioned) has the price process \(S(t)\) given by
\begin{equation}
	\label{eq:GBMRiskNeutral}
	dS(t) = rS(t) \, \di t + \sigma S(t) \, \di \widetilde{W}(t),
\end{equation}
where the interest rate \(r\) and the volatility \(\sigma\) are strictly positive constants and \(\widetilde{W}(t)\) is a Brownian motion under the risk-neutral probability measure \(\widetilde{\Prob}\). The perpetual American put pays \(K - S(t)\) if it is exercised at time \(t\). This is its \textit{intrinsic value}.

\begin{defn}{Perpetual American Put}{PerpetualAmericanPut}
	Let \(\mathcal{T}\) be the set of all stopping times. The price of the perpetual American put is defined to be
	\begin{equation}
		\label{eq:PerpetualAmericanPutPrice}
		v_*(x) = \max_{\tau \in \mathcal{T}} \widetilde{\E} \left[ e^{-r\tau} (K - S(\tau)) \right],
	\end{equation}
	where \(x = S(0)\) in \Cref{eq:PerpetualAmericanPutPrice} is the initial stock price. In the event that \(\tau = \infty\), we interpret \(e^{-r\tau}(K - S(\tau))\) to be zero.
\end{defn}

The idea behind \Cref{defn:PerpetualAmericanPut} is that the owner of the perpetual American put can choose an exercise time \(\tau\), subject only to the condition that she may not look ahead to determine when to exercise. The mathematical formulation of this ``not look ahead'' restriction is that \(\tau\) must be a stopping time. The price of the option at time zero is the risk-neutral expected payoff of the option, discounted from the exercise time back to time zero. If the option is never exercised, its payoff is zero. This explains the term under the expectation on the right-hand side of \Cref{eq:PerpetualAmericanPutPrice}. The owner of the option should choose the exercise strategy that maximizes this expected payoff, discounted back to time zero, and thus we define the price of the option to be the maximum over \(\tau \in \mathcal{T}\) of the discounted expected payoffs.

The owner of the perpetual American put can exercise at any time. In particular, there is no expiration date after which the put can no longer be exercised. This makes every date like every other date; the time remaining to expiration is always the same (i.e., infinity). Because every date is like every other date, it is reasonable to expect that the optimal exercise policy depends only on the value of \(S(t)\) and not on the time variable \(t\). The owner of the put should exercise as soon as \(S(t)\) falls far enough below \(K\). In other words, it is reasonable to expect that the optimal exercise policy is of the form
\begin{quote} Exercise the put as soon as \(S(t)\) falls to the level \(L_*\).
\end{quote}

We have two questions to answer:
\begin{enumerate}[label=\textup{(\roman*)}]
\item What is the value of \(L_*\), and how do we know it corresponds to optimal exercise?
	\item What is the value of the put?
\end{enumerate}

For the perpetual American put, we can base the answers to these questions on explicit computations.

\subsubsection{Price under Arbitrary Exercise}

\begin{thm}{Laplace Transform For First Passage Time Of Drifted Brownian Motion}{LaplaceTransformForFirstPassageTimeOfDriftedBrownianMotion}
	Let \(\widetilde{W}(t)\) be a Brownian motion under a probability measure \(\widetilde{\Prob}\), let \(\mu\) be a real number, and let \(m\) be a positive number. Define \(X(t) = \mu t + \widetilde{W}(t)\), and set
	\[
		\tau_m = \min\{t \ge 0 ; X(t) = m\},
	\]
	If \(X(t)\) never reaches the level \(m\), then we interpret \(\tau_m\) to be \(\infty\). Then
	\begin{equation} \label{eq:LaplaceTransformFormula}
		\widetilde{\E}\e^{-\lambda \tau_m} = \e^{-m(-\mu+\sqrt{\mu^2+2\lambda})} \text{ for all } \lambda > 0,
	\end{equation}
	where we interpret \(\e^{-\lambda \tau_m}\) to be zero if \(\tau_m = \infty\).
\end{thm}

\begin{proof}

	Define \(\sigma = -\mu + \sqrt{\mu^2 + 2\lambda}\) so that \(\sigma > 0\) and
	\begin{align*}
		\sigma \mu + \frac{1}{2}\sigma^2 & = -\mu^2 + \mu\sqrt{\mu^2 + 2\lambda} + \frac{1}{2}\left(-\mu + \sqrt{\mu^2 + 2\lambda}\right)^2                   \\
		                                 & = -\mu^2 + \mu\sqrt{\mu^2 + 2\lambda} + \frac{1}{2}\mu^2 - \mu\sqrt{\mu^2 + 2\lambda} + \frac{1}{2}\mu^2 + \lambda \\
		                                 & = \lambda.
	\end{align*}
	Then
	\[
		\e^{\sigma X(t)-\lambda t} = \e^{\sigma \mu t + \sigma \widetilde{W}(t) - \sigma \mu t - \frac{1}{2}\sigma^2 t} = \e^{\sigma \widetilde{W}(t) - \frac{1}{2}\sigma^2 t}.
	\]
	This process is an exponential martingale under \(\widetilde{\Prob}\) (its differential has a \(d\widetilde{W}(t)\) term and no \(dt\) term). According to \Cref{thm:DoobsStoppingTheoremForContinuousMartingales}, the stopped martingale
	\[
		M(t) = \e^{\sigma \widetilde{W}(t \wedge \tau_m) - \frac{1}{2}\sigma^2 (t \wedge \tau_m)}
	\]
	is also a martingale. Therefore, for each positive integer \(n\),
	\begin{align}
		1 = M(0) & = \widetilde{\E}M(n) \nonumber                                                                                                                                                            \\
		         & = \widetilde{\E}\left[\e^{\sigma X(n \wedge \tau_m) - \lambda(n \wedge \tau_m)}\right] \nonumber                                                                                          \\
		         & = \widetilde{\E}\left[\e^{\sigma m - \lambda \tau_m}\I_{\{\tau_m \le n\}}\right] + \widetilde{\E}\left[\e^{\sigma X(n) - \lambda n}\I_{\{\tau_m > n\}}\right]. \label{eq:MartingaleSplit}
	\end{align}
	The nonnegative random variables \(\e^{\sigma m - \lambda \tau_m}\I_{\{\tau_m \le n\}}\) increase with \(n\), and their limit is \(\e^{\sigma m - \lambda \tau_m}\I_{\{\tau_m < \infty\}}\). In other words,
	\[
		0 \le \e^{\sigma m - \lambda \tau_m}\I_{\{\tau_m \le 1\}} \le \e^{\sigma m - \lambda \tau_m}\I_{\{\tau_m \le 2\}} \le \dots \text{ almost surely,}
	\]
	and
	\[
		\lim_{n \to \infty} \e^{\sigma m - \lambda \tau_m}\I_{\{\tau_m \le n\}} = \e^{\sigma m - \lambda \tau_m}\I_{\{\tau_m < \infty\}} \text{ almost surely.}
	\]
	The Monotone Convergence Theorem~\Cref{thm:PropertiesOfAnExpectation} implies
	\begin{equation} \label{eq:MonotoneConvergenceLimit}
		\lim_{n \to \infty} \widetilde{\E}\left[\e^{\sigma m - \lambda \tau_m}\I_{\{\tau_m \le n\}}\right] = \widetilde{\E}\left[\e^{\sigma m - \lambda \tau_m}\I_{\{\tau_m < \infty\}}\right].
	\end{equation}
	On the other hand, the random variable \(\e^{\sigma X(n) - \lambda n}\I_{\{\tau_m > n\}}\) satisfies
	\[
		0 \le \e^{\sigma X(n) - \lambda n}\I_{\{\tau_m > n\}} \le \e^{\sigma m - \lambda n} \le \e^{\sigma m} \text{ almost surely}
	\]
	because \(X(n) \le m\) for \(n < \tau_m\) and \(\sigma\) is positive. Because \(\lambda\) is positive, we have
	\[
		\lim_{n \to \infty} \e^{\sigma X(n) - \lambda n}\I_{\{\tau_m > n\}} \le \lim_{n \to \infty} \e^{\sigma m - \lambda n} = 0.
	\]
	According to the Dominated Convergence Theorem~\Cref{thm:DominatedAndBoundedConvergence}
	\begin{equation} \label{eq:DominatedConvergenceLimit}
		\lim_{n \to \infty} \widetilde{\E}\left[\e^{\sigma X(n) - \lambda n}\I_{\{\tau_m > n\}}\right] = 0.
	\end{equation}
	Taking the limit in \Cref{eq:MartingaleSplit} and using \Cref{eq:MonotoneConvergenceLimit,eq:DominatedConvergenceLimit}, we obtain
	\[
		1 = \widetilde{\E}\left[\e^{\sigma m - \lambda \tau_m}\I_{\{\tau_m < \infty\}}\right]
	\]
	or, equivalently,
	\begin{equation} \label{eq:FinalLaplaceResult}
		\widetilde{\E}\left[\e^{-\lambda \tau_m}\I_{\{\tau_m < \infty\}}\right] = \e^{-\sigma m} = \e^{-m(-\mu + \sqrt{\mu^2 + 2\lambda})} \text{ for all } \lambda > 0.
	\end{equation}
	This is \Cref{eq:LaplaceTransformFormula} when we interpret \(\e^{-\lambda \tau_m}\) to be zero if \(\tau_m = \infty\).
\end{proof}

\begin{xtip}{Limit Case for Lambda}{LimitCaseLambda}
	We used the strict positivity of \(\lambda\) to derive (8.3.7), but now that we have it, we can take the limit as \(\lambda \downarrow 0\). The random variables \(e^{-\lambda \tau_m} \I_{\set{\tau_m < \infty}}\) are nonnegative and increase to \(\I_{\set{\tau_m < \infty}}\) as \(\lambda \downarrow 0\), and the Monotone Convergence Theorem, \Cref{thm:PropertiesOfAnExpectation}, allows us to conclude that
	\begin{equation*}
		\widetilde{\Prob}\set{\tau_m < \infty} = \widetilde{\E}\I_{\set{\tau_m < \infty}} = \lim_{\lambda \downarrow 0} e^{-m(-\mu + \sqrt{\mu^2 + 2\lambda})} = e^{m\mu - m\abs{\mu}}.
	\end{equation*}
	If \(\mu \ge 0\), the drift in \(X(t)\) is zero or upward, toward level \(m\), and \(\widetilde{\Prob}\set{\tau_m < \infty} = 1\); the level \(X(t)\) is reached with probability one. On the other hand, if \(\mu < 0\), the drift in \(X(t)\) is downward, away from level \(m\), and \(\widetilde{\Prob}\set{\tau_m < \infty} = e^{-2m\abs{\mu}} < 1\); there is a positive probability of never reaching \(m\).
\end{xtip}

The solution to \Cref{eq:GBMRiskNeutral} is
\begin{equation} \label{eq:StockPriceSolution}
	S(t) = S(0) \exp \set{ \sigma \widetilde{W}(t) + \left( r - \frac{1}{2}\sigma^2 \right) t }.
\end{equation}
Suppose the owner of the perpetual American put sets a positive level \(L < K\) and resolves to exercise the put the first time the stock price falls to \(L\). If the initial stock price is at or below \(L\), she exercises immediately (at time zero). The value of the put in this case is \(v_L(S(0)) = K - S(0)\). If the initial stock price is above \(L\), she exercises at the stopping time
\begin{equation} \label{eq:StoppingTime}
	\tau_L = \min\set{t \ge 0; S(t) = L},
\end{equation}
where \(\tau_L\) is set equal to \(\infty\) if the stock price never reaches the level \(L\). At the time of exercise, the put pays \(K - S(\tau_L) = K - L\). Discounting this back to time zero and taking the risk-neutral expected value, we compute the value of the put under this exercise strategy to be
\begin{equation} \label{eq:RiskNeutralValue}
	v_L(S(0)) = (K - L) \widetilde{\E}e^{-r\tau_L} \text{ for all } S(0) \ge L.
\end{equation}
On those paths where \(\tau_L = \infty\), we interpret \(e^{-r\tau_L}\) to be zero. (Recall our assumption at the beginning of this section that \(r\) is strictly positive.) Although not explicitly indicated by the notation, the distribution of \(\tau_L\) depends on the initial stock price \(S(0)\), so the right-hand side~\eqref{eq:RiskNeutralValue} is a function of \(S(0)\).

\begin{lem}{Value Function Formula}{ValueFunctionFormula}
	The function \(v_L(x)\) is given by the formula
	\begin{equation} \label{eq:ValueFunctionPiecewise}
		v_L(x) = \begin{cases}
			K - x,                                                     & 0 \le x \le L, \\
			(K - L) \left( \frac{x}{L} \right)^{-\frac{2r}{\sigma^2}}, & x \ge L.
		\end{cases}
	\end{equation}
\end{lem}

\begin{proof}
	We only need to establish the second line of \Cref{eq:ValueFunctionPiecewise}. If \(x = L\), then \(\tau_L = 0\) and \Cref{eq:RiskNeutralValue} implies \(v_L(x) = K - L\).

	We consider the case \(S(0) = x > L\). The stopping time \(\tau_L\) is the first time \(S(t)\) reaches the level \(L\). But \(S(t) = L\) if and only if
	\[
		S(t) = x \exp \left\{ \sigma \widetilde{W}(t) + \left( r - \frac{1}{2}\sigma^2 \right) t \right\}
	\]
	implies
	\[
		-\widetilde{W}(t) - \frac{1}{\sigma}\left( r - \frac{1}{2}\sigma^2 \right) t = \frac{1}{\sigma} \log \frac{x}{L}.
	\]
	We now apply \Cref{thm:LaplaceTransformForFirstPassageTimeOfDriftedBrownianMotion} with \(X(t)\) in that theorem replaced by \(-\widetilde{W}(t) - \frac{1}{\sigma}\left( r - \frac{1}{2}\sigma^2 \right) t\) (the processes \(\widetilde{W}(t)\) and \(-\widetilde{W}(t)\) are both Brownian motions under \(\widetilde{\Prob}\)), with \(\lambda\) replaced by \(r\), with \(\mu\) replaced by \(-\frac{1}{\sigma}\left( r - \frac{1}{2}\sigma^2 \right)\), and with \(m\) replaced by \(\frac{1}{\sigma} \log \frac{x}{L}\), which is positive. With these replacements, \(\tau_m\) in \Cref{thm:LaplaceTransformForFirstPassageTimeOfDriftedBrownianMotion} is \(\tau_L\) and
	\begin{align*}
		\mu^2 + 2\lambda & = \frac{1}{\sigma^2} \left( r^2 - r\sigma^2 + \frac{1}{4}\sigma^4 \right) + 2r \\
		                 & = \frac{1}{\sigma^2} \left( r^2 + r\sigma^2 + \frac{1}{4}\sigma^4 \right)      \\
		                 & = \frac{1}{\sigma^2} \left( r + \frac{1}{2}\sigma^2 \right)^2.
	\end{align*}
	Therefore,
	\[
		-\mu + \sqrt{\mu^2 + 2\lambda} = \frac{1}{\sigma}\left( r - \frac{1}{2}\sigma^2 \right) + \frac{1}{\sigma}\left( r + \frac{1}{2}\sigma^2 \right) = \frac{2r}{\sigma}.
	\]
	\Cref{eq:LaplaceTransformFormula} implies
	\[
		\widetilde{\E} \left[ \e^{-r\tau_L} \right] = \exp \left\{ -\frac{1}{\sigma} \log \frac{x}{L} \cdot \frac{2r}{\sigma} \right\} = \left( \frac{x}{L} \right)^{-\frac{2r}{\sigma^2}}.
	\]
	The second line in \Cref{eq:ValueFunctionPiecewise} follows.
\end{proof}

\subsubsection{Price under Optimal Exercise}

\Cref{fig:ExerciseBoundaryValue} shows the function \(v_L(x)\) for three different values of \(L\). The function \(v_{L_1}(x)\) in that figure actually lies below the intrinsic value \(K - x\) for \(x\) between \(L_1\) and \(L_2\). If the initial stock price is between \(L_1\) and \(L_2\), then the strategy of exercising the first time the stock price falls to \(L_1\) is obviously a poor one; it would be better to exercise at time zero and receive the intrinsic value. The function \(v_{L_2}(x)\) agrees with the intrinsic value for \(0 \le x \le L_2\) and follows the indicated curve for \(x \ge L_2\). The function \(v_{L_*}(x)\) agrees with the intrinsic value for \(0 \le x \le L_*\) and follows the indicated curve for \(x \ge L_*\). For \(x \ge L_*\), the function \(v_{L_*}(x)\) is strictly larger than the function \(v_{L_2}(x)\), and hence the strategy of exercising the first time the stock price falls to \(L_*\) is better than exercising the first time the stock price falls to \(L_2\).


\begin{Pcode}{PerpetualPut.py}
	# Model Parameters
	K, alpha = 10.0, 0.82
	L_star = (alpha / (1 + alpha)) * K  # Optimal exercise boundary
	x = np.linspace(0.1, 15.0, 500)     # State space grid

	# Valuation
	v = np.where(x > L_star,
	(K - L_star) * (x / L_star) ** -alpha,
	K - x)
\end{Pcode}

\begin{figure}[H]
	\centering
	%% Creator: Matplotlib, PGF backend
%%
%% To include the figure in your LaTeX document, write
%%   \input{<filename>.pgf}
%%
%% Make sure the required packages are loaded in your preamble
%%   \usepackage{pgf}
%%
%% Also ensure that all the required font packages are loaded; for instance,
%% the lmodern package is sometimes necessary when using math font.
%%   \usepackage{lmodern}
%%
%% Figures using additional raster images can only be included by \input if
%% they are in the same directory as the main LaTeX file. For loading figures
%% from other directories you can use the `import` package
%%   \usepackage{import}
%%
%% and then include the figures with
%%   \import{<path to file>}{<filename>.pgf}
%%
%% Matplotlib used the following preamble
%%   \def\mathdefault#1{#1}
%%   \everymath=\expandafter{\the\everymath\displaystyle}
%%   \IfFileExists{scrextend.sty}{
%%     \usepackage[fontsize=10.000000pt]{scrextend}
%%   }{
%%     \renewcommand{\normalsize}{\fontsize{10.000000}{12.000000}\selectfont}
%%     \normalsize
%%   }
%%   \usepackage{amsmath}
%%   \usepackage{amssymb}
%%   \usepackage{amsfonts}
%%   \usepackage{libertine}
%%   \usepackage[libertine]{newtxmath}
%%   \makeatletter\@ifpackageloaded{underscore}{}{\usepackage[strings]{underscore}}\makeatother
%%
\begingroup%
\makeatletter%
\begin{pgfpicture}%
\pgfpathrectangle{\pgfpointorigin}{\pgfqpoint{6.249663in}{4.030868in}}%
\pgfusepath{use as bounding box, clip}%
\begin{pgfscope}%
\pgfsetbuttcap%
\pgfsetmiterjoin%
\definecolor{currentfill}{rgb}{1.000000,1.000000,1.000000}%
\pgfsetfillcolor{currentfill}%
\pgfsetlinewidth{0.000000pt}%
\definecolor{currentstroke}{rgb}{1.000000,1.000000,1.000000}%
\pgfsetstrokecolor{currentstroke}%
\pgfsetdash{}{0pt}%
\pgfpathmoveto{\pgfqpoint{0.000000in}{0.000000in}}%
\pgfpathlineto{\pgfqpoint{6.249663in}{0.000000in}}%
\pgfpathlineto{\pgfqpoint{6.249663in}{4.030868in}}%
\pgfpathlineto{\pgfqpoint{0.000000in}{4.030868in}}%
\pgfpathlineto{\pgfqpoint{0.000000in}{0.000000in}}%
\pgfpathclose%
\pgfusepath{fill}%
\end{pgfscope}%
\begin{pgfscope}%
\pgfsetbuttcap%
\pgfsetmiterjoin%
\definecolor{currentfill}{rgb}{1.000000,1.000000,1.000000}%
\pgfsetfillcolor{currentfill}%
\pgfsetlinewidth{0.000000pt}%
\definecolor{currentstroke}{rgb}{0.000000,0.000000,0.000000}%
\pgfsetstrokecolor{currentstroke}%
\pgfsetstrokeopacity{0.000000}%
\pgfsetdash{}{0pt}%
\pgfpathmoveto{\pgfqpoint{0.396477in}{0.232212in}}%
\pgfpathlineto{\pgfqpoint{6.171885in}{0.232212in}}%
\pgfpathlineto{\pgfqpoint{6.171885in}{3.904501in}}%
\pgfpathlineto{\pgfqpoint{0.396477in}{3.904501in}}%
\pgfpathlineto{\pgfqpoint{0.396477in}{0.232212in}}%
\pgfpathclose%
\pgfusepath{fill}%
\end{pgfscope}%
\begin{pgfscope}%
\pgfsetbuttcap%
\pgfsetroundjoin%
\definecolor{currentfill}{rgb}{0.000000,0.000000,0.000000}%
\pgfsetfillcolor{currentfill}%
\pgfsetlinewidth{0.501875pt}%
\definecolor{currentstroke}{rgb}{0.000000,0.000000,0.000000}%
\pgfsetstrokecolor{currentstroke}%
\pgfsetdash{}{0pt}%
\pgfsys@defobject{currentmarker}{\pgfqpoint{0.000000in}{0.000000in}}{\pgfqpoint{0.000000in}{0.041667in}}{%
\pgfpathmoveto{\pgfqpoint{0.000000in}{0.000000in}}%
\pgfpathlineto{\pgfqpoint{0.000000in}{0.041667in}}%
\pgfusepath{stroke,fill}%
}%
\begin{pgfscope}%
\pgfsys@transformshift{1.359045in}{0.232212in}%
\pgfsys@useobject{currentmarker}{}%
\end{pgfscope}%
\end{pgfscope}%
\begin{pgfscope}%
\pgfsetbuttcap%
\pgfsetroundjoin%
\definecolor{currentfill}{rgb}{0.000000,0.000000,0.000000}%
\pgfsetfillcolor{currentfill}%
\pgfsetlinewidth{0.501875pt}%
\definecolor{currentstroke}{rgb}{0.000000,0.000000,0.000000}%
\pgfsetstrokecolor{currentstroke}%
\pgfsetdash{}{0pt}%
\pgfsys@defobject{currentmarker}{\pgfqpoint{0.000000in}{-0.041667in}}{\pgfqpoint{0.000000in}{0.000000in}}{%
\pgfpathmoveto{\pgfqpoint{0.000000in}{0.000000in}}%
\pgfpathlineto{\pgfqpoint{0.000000in}{-0.041667in}}%
\pgfusepath{stroke,fill}%
}%
\begin{pgfscope}%
\pgfsys@transformshift{1.359045in}{3.904501in}%
\pgfsys@useobject{currentmarker}{}%
\end{pgfscope}%
\end{pgfscope}%
\begin{pgfscope}%
\definecolor{textcolor}{rgb}{0.000000,0.000000,0.000000}%
\pgfsetstrokecolor{textcolor}%
\pgfsetfillcolor{textcolor}%
\pgftext[x=1.359045in,y=0.183600in,,top]{\color{textcolor}{\rmfamily\fontsize{10.000000}{12.000000}\selectfont\catcode`\^=\active\def^{\ifmmode\sp\else\^{}\fi}\catcode`\%=\active\def%{\%}$L_1$}}%
\end{pgfscope}%
\begin{pgfscope}%
\pgfsetbuttcap%
\pgfsetroundjoin%
\definecolor{currentfill}{rgb}{0.000000,0.000000,0.000000}%
\pgfsetfillcolor{currentfill}%
\pgfsetlinewidth{0.501875pt}%
\definecolor{currentstroke}{rgb}{0.000000,0.000000,0.000000}%
\pgfsetstrokecolor{currentstroke}%
\pgfsetdash{}{0pt}%
\pgfsys@defobject{currentmarker}{\pgfqpoint{0.000000in}{0.000000in}}{\pgfqpoint{0.000000in}{0.041667in}}{%
\pgfpathmoveto{\pgfqpoint{0.000000in}{0.000000in}}%
\pgfpathlineto{\pgfqpoint{0.000000in}{0.041667in}}%
\pgfusepath{stroke,fill}%
}%
\begin{pgfscope}%
\pgfsys@transformshift{2.131215in}{0.232212in}%
\pgfsys@useobject{currentmarker}{}%
\end{pgfscope}%
\end{pgfscope}%
\begin{pgfscope}%
\pgfsetbuttcap%
\pgfsetroundjoin%
\definecolor{currentfill}{rgb}{0.000000,0.000000,0.000000}%
\pgfsetfillcolor{currentfill}%
\pgfsetlinewidth{0.501875pt}%
\definecolor{currentstroke}{rgb}{0.000000,0.000000,0.000000}%
\pgfsetstrokecolor{currentstroke}%
\pgfsetdash{}{0pt}%
\pgfsys@defobject{currentmarker}{\pgfqpoint{0.000000in}{-0.041667in}}{\pgfqpoint{0.000000in}{0.000000in}}{%
\pgfpathmoveto{\pgfqpoint{0.000000in}{0.000000in}}%
\pgfpathlineto{\pgfqpoint{0.000000in}{-0.041667in}}%
\pgfusepath{stroke,fill}%
}%
\begin{pgfscope}%
\pgfsys@transformshift{2.131215in}{3.904501in}%
\pgfsys@useobject{currentmarker}{}%
\end{pgfscope}%
\end{pgfscope}%
\begin{pgfscope}%
\definecolor{textcolor}{rgb}{0.000000,0.000000,0.000000}%
\pgfsetstrokecolor{textcolor}%
\pgfsetfillcolor{textcolor}%
\pgftext[x=2.131215in,y=0.183600in,,top]{\color{textcolor}{\rmfamily\fontsize{10.000000}{12.000000}\selectfont\catcode`\^=\active\def^{\ifmmode\sp\else\^{}\fi}\catcode`\%=\active\def%{\%}$L_*$}}%
\end{pgfscope}%
\begin{pgfscope}%
\pgfsetbuttcap%
\pgfsetroundjoin%
\definecolor{currentfill}{rgb}{0.000000,0.000000,0.000000}%
\pgfsetfillcolor{currentfill}%
\pgfsetlinewidth{0.501875pt}%
\definecolor{currentstroke}{rgb}{0.000000,0.000000,0.000000}%
\pgfsetstrokecolor{currentstroke}%
\pgfsetdash{}{0pt}%
\pgfsys@defobject{currentmarker}{\pgfqpoint{0.000000in}{0.000000in}}{\pgfqpoint{0.000000in}{0.041667in}}{%
\pgfpathmoveto{\pgfqpoint{0.000000in}{0.000000in}}%
\pgfpathlineto{\pgfqpoint{0.000000in}{0.041667in}}%
\pgfusepath{stroke,fill}%
}%
\begin{pgfscope}%
\pgfsys@transformshift{3.091667in}{0.232212in}%
\pgfsys@useobject{currentmarker}{}%
\end{pgfscope}%
\end{pgfscope}%
\begin{pgfscope}%
\pgfsetbuttcap%
\pgfsetroundjoin%
\definecolor{currentfill}{rgb}{0.000000,0.000000,0.000000}%
\pgfsetfillcolor{currentfill}%
\pgfsetlinewidth{0.501875pt}%
\definecolor{currentstroke}{rgb}{0.000000,0.000000,0.000000}%
\pgfsetstrokecolor{currentstroke}%
\pgfsetdash{}{0pt}%
\pgfsys@defobject{currentmarker}{\pgfqpoint{0.000000in}{-0.041667in}}{\pgfqpoint{0.000000in}{0.000000in}}{%
\pgfpathmoveto{\pgfqpoint{0.000000in}{0.000000in}}%
\pgfpathlineto{\pgfqpoint{0.000000in}{-0.041667in}}%
\pgfusepath{stroke,fill}%
}%
\begin{pgfscope}%
\pgfsys@transformshift{3.091667in}{3.904501in}%
\pgfsys@useobject{currentmarker}{}%
\end{pgfscope}%
\end{pgfscope}%
\begin{pgfscope}%
\definecolor{textcolor}{rgb}{0.000000,0.000000,0.000000}%
\pgfsetstrokecolor{textcolor}%
\pgfsetfillcolor{textcolor}%
\pgftext[x=3.091667in,y=0.183600in,,top]{\color{textcolor}{\rmfamily\fontsize{10.000000}{12.000000}\selectfont\catcode`\^=\active\def^{\ifmmode\sp\else\^{}\fi}\catcode`\%=\active\def%{\%}$L_2$}}%
\end{pgfscope}%
\begin{pgfscope}%
\pgfsetbuttcap%
\pgfsetroundjoin%
\definecolor{currentfill}{rgb}{0.000000,0.000000,0.000000}%
\pgfsetfillcolor{currentfill}%
\pgfsetlinewidth{0.501875pt}%
\definecolor{currentstroke}{rgb}{0.000000,0.000000,0.000000}%
\pgfsetstrokecolor{currentstroke}%
\pgfsetdash{}{0pt}%
\pgfsys@defobject{currentmarker}{\pgfqpoint{0.000000in}{0.000000in}}{\pgfqpoint{0.000000in}{0.041667in}}{%
\pgfpathmoveto{\pgfqpoint{0.000000in}{0.000000in}}%
\pgfpathlineto{\pgfqpoint{0.000000in}{0.041667in}}%
\pgfusepath{stroke,fill}%
}%
\begin{pgfscope}%
\pgfsys@transformshift{4.246749in}{0.232212in}%
\pgfsys@useobject{currentmarker}{}%
\end{pgfscope}%
\end{pgfscope}%
\begin{pgfscope}%
\pgfsetbuttcap%
\pgfsetroundjoin%
\definecolor{currentfill}{rgb}{0.000000,0.000000,0.000000}%
\pgfsetfillcolor{currentfill}%
\pgfsetlinewidth{0.501875pt}%
\definecolor{currentstroke}{rgb}{0.000000,0.000000,0.000000}%
\pgfsetstrokecolor{currentstroke}%
\pgfsetdash{}{0pt}%
\pgfsys@defobject{currentmarker}{\pgfqpoint{0.000000in}{-0.041667in}}{\pgfqpoint{0.000000in}{0.000000in}}{%
\pgfpathmoveto{\pgfqpoint{0.000000in}{0.000000in}}%
\pgfpathlineto{\pgfqpoint{0.000000in}{-0.041667in}}%
\pgfusepath{stroke,fill}%
}%
\begin{pgfscope}%
\pgfsys@transformshift{4.246749in}{3.904501in}%
\pgfsys@useobject{currentmarker}{}%
\end{pgfscope}%
\end{pgfscope}%
\begin{pgfscope}%
\definecolor{textcolor}{rgb}{0.000000,0.000000,0.000000}%
\pgfsetstrokecolor{textcolor}%
\pgfsetfillcolor{textcolor}%
\pgftext[x=4.246749in,y=0.183600in,,top]{\color{textcolor}{\rmfamily\fontsize{10.000000}{12.000000}\selectfont\catcode`\^=\active\def^{\ifmmode\sp\else\^{}\fi}\catcode`\%=\active\def%{\%}$K$}}%
\end{pgfscope}%
\begin{pgfscope}%
\pgfsetbuttcap%
\pgfsetroundjoin%
\definecolor{currentfill}{rgb}{0.000000,0.000000,0.000000}%
\pgfsetfillcolor{currentfill}%
\pgfsetlinewidth{0.501875pt}%
\definecolor{currentstroke}{rgb}{0.000000,0.000000,0.000000}%
\pgfsetstrokecolor{currentstroke}%
\pgfsetdash{}{0pt}%
\pgfsys@defobject{currentmarker}{\pgfqpoint{0.000000in}{0.000000in}}{\pgfqpoint{0.000000in}{0.020833in}}{%
\pgfpathmoveto{\pgfqpoint{0.000000in}{0.000000in}}%
\pgfpathlineto{\pgfqpoint{0.000000in}{0.020833in}}%
\pgfusepath{stroke,fill}%
}%
\begin{pgfscope}%
\pgfsys@transformshift{0.586875in}{0.232212in}%
\pgfsys@useobject{currentmarker}{}%
\end{pgfscope}%
\end{pgfscope}%
\begin{pgfscope}%
\pgfsetbuttcap%
\pgfsetroundjoin%
\definecolor{currentfill}{rgb}{0.000000,0.000000,0.000000}%
\pgfsetfillcolor{currentfill}%
\pgfsetlinewidth{0.501875pt}%
\definecolor{currentstroke}{rgb}{0.000000,0.000000,0.000000}%
\pgfsetstrokecolor{currentstroke}%
\pgfsetdash{}{0pt}%
\pgfsys@defobject{currentmarker}{\pgfqpoint{0.000000in}{-0.020833in}}{\pgfqpoint{0.000000in}{0.000000in}}{%
\pgfpathmoveto{\pgfqpoint{0.000000in}{0.000000in}}%
\pgfpathlineto{\pgfqpoint{0.000000in}{-0.020833in}}%
\pgfusepath{stroke,fill}%
}%
\begin{pgfscope}%
\pgfsys@transformshift{0.586875in}{3.904501in}%
\pgfsys@useobject{currentmarker}{}%
\end{pgfscope}%
\end{pgfscope}%
\begin{pgfscope}%
\pgfsetbuttcap%
\pgfsetroundjoin%
\definecolor{currentfill}{rgb}{0.000000,0.000000,0.000000}%
\pgfsetfillcolor{currentfill}%
\pgfsetlinewidth{0.501875pt}%
\definecolor{currentstroke}{rgb}{0.000000,0.000000,0.000000}%
\pgfsetstrokecolor{currentstroke}%
\pgfsetdash{}{0pt}%
\pgfsys@defobject{currentmarker}{\pgfqpoint{0.000000in}{0.000000in}}{\pgfqpoint{0.000000in}{0.020833in}}{%
\pgfpathmoveto{\pgfqpoint{0.000000in}{0.000000in}}%
\pgfpathlineto{\pgfqpoint{0.000000in}{0.020833in}}%
\pgfusepath{stroke,fill}%
}%
\begin{pgfscope}%
\pgfsys@transformshift{0.779917in}{0.232212in}%
\pgfsys@useobject{currentmarker}{}%
\end{pgfscope}%
\end{pgfscope}%
\begin{pgfscope}%
\pgfsetbuttcap%
\pgfsetroundjoin%
\definecolor{currentfill}{rgb}{0.000000,0.000000,0.000000}%
\pgfsetfillcolor{currentfill}%
\pgfsetlinewidth{0.501875pt}%
\definecolor{currentstroke}{rgb}{0.000000,0.000000,0.000000}%
\pgfsetstrokecolor{currentstroke}%
\pgfsetdash{}{0pt}%
\pgfsys@defobject{currentmarker}{\pgfqpoint{0.000000in}{-0.020833in}}{\pgfqpoint{0.000000in}{0.000000in}}{%
\pgfpathmoveto{\pgfqpoint{0.000000in}{0.000000in}}%
\pgfpathlineto{\pgfqpoint{0.000000in}{-0.020833in}}%
\pgfusepath{stroke,fill}%
}%
\begin{pgfscope}%
\pgfsys@transformshift{0.779917in}{3.904501in}%
\pgfsys@useobject{currentmarker}{}%
\end{pgfscope}%
\end{pgfscope}%
\begin{pgfscope}%
\pgfsetbuttcap%
\pgfsetroundjoin%
\definecolor{currentfill}{rgb}{0.000000,0.000000,0.000000}%
\pgfsetfillcolor{currentfill}%
\pgfsetlinewidth{0.501875pt}%
\definecolor{currentstroke}{rgb}{0.000000,0.000000,0.000000}%
\pgfsetstrokecolor{currentstroke}%
\pgfsetdash{}{0pt}%
\pgfsys@defobject{currentmarker}{\pgfqpoint{0.000000in}{0.000000in}}{\pgfqpoint{0.000000in}{0.020833in}}{%
\pgfpathmoveto{\pgfqpoint{0.000000in}{0.000000in}}%
\pgfpathlineto{\pgfqpoint{0.000000in}{0.020833in}}%
\pgfusepath{stroke,fill}%
}%
\begin{pgfscope}%
\pgfsys@transformshift{0.972960in}{0.232212in}%
\pgfsys@useobject{currentmarker}{}%
\end{pgfscope}%
\end{pgfscope}%
\begin{pgfscope}%
\pgfsetbuttcap%
\pgfsetroundjoin%
\definecolor{currentfill}{rgb}{0.000000,0.000000,0.000000}%
\pgfsetfillcolor{currentfill}%
\pgfsetlinewidth{0.501875pt}%
\definecolor{currentstroke}{rgb}{0.000000,0.000000,0.000000}%
\pgfsetstrokecolor{currentstroke}%
\pgfsetdash{}{0pt}%
\pgfsys@defobject{currentmarker}{\pgfqpoint{0.000000in}{-0.020833in}}{\pgfqpoint{0.000000in}{0.000000in}}{%
\pgfpathmoveto{\pgfqpoint{0.000000in}{0.000000in}}%
\pgfpathlineto{\pgfqpoint{0.000000in}{-0.020833in}}%
\pgfusepath{stroke,fill}%
}%
\begin{pgfscope}%
\pgfsys@transformshift{0.972960in}{3.904501in}%
\pgfsys@useobject{currentmarker}{}%
\end{pgfscope}%
\end{pgfscope}%
\begin{pgfscope}%
\pgfsetbuttcap%
\pgfsetroundjoin%
\definecolor{currentfill}{rgb}{0.000000,0.000000,0.000000}%
\pgfsetfillcolor{currentfill}%
\pgfsetlinewidth{0.501875pt}%
\definecolor{currentstroke}{rgb}{0.000000,0.000000,0.000000}%
\pgfsetstrokecolor{currentstroke}%
\pgfsetdash{}{0pt}%
\pgfsys@defobject{currentmarker}{\pgfqpoint{0.000000in}{0.000000in}}{\pgfqpoint{0.000000in}{0.020833in}}{%
\pgfpathmoveto{\pgfqpoint{0.000000in}{0.000000in}}%
\pgfpathlineto{\pgfqpoint{0.000000in}{0.020833in}}%
\pgfusepath{stroke,fill}%
}%
\begin{pgfscope}%
\pgfsys@transformshift{1.166002in}{0.232212in}%
\pgfsys@useobject{currentmarker}{}%
\end{pgfscope}%
\end{pgfscope}%
\begin{pgfscope}%
\pgfsetbuttcap%
\pgfsetroundjoin%
\definecolor{currentfill}{rgb}{0.000000,0.000000,0.000000}%
\pgfsetfillcolor{currentfill}%
\pgfsetlinewidth{0.501875pt}%
\definecolor{currentstroke}{rgb}{0.000000,0.000000,0.000000}%
\pgfsetstrokecolor{currentstroke}%
\pgfsetdash{}{0pt}%
\pgfsys@defobject{currentmarker}{\pgfqpoint{0.000000in}{-0.020833in}}{\pgfqpoint{0.000000in}{0.000000in}}{%
\pgfpathmoveto{\pgfqpoint{0.000000in}{0.000000in}}%
\pgfpathlineto{\pgfqpoint{0.000000in}{-0.020833in}}%
\pgfusepath{stroke,fill}%
}%
\begin{pgfscope}%
\pgfsys@transformshift{1.166002in}{3.904501in}%
\pgfsys@useobject{currentmarker}{}%
\end{pgfscope}%
\end{pgfscope}%
\begin{pgfscope}%
\pgfsetbuttcap%
\pgfsetroundjoin%
\definecolor{currentfill}{rgb}{0.000000,0.000000,0.000000}%
\pgfsetfillcolor{currentfill}%
\pgfsetlinewidth{0.501875pt}%
\definecolor{currentstroke}{rgb}{0.000000,0.000000,0.000000}%
\pgfsetstrokecolor{currentstroke}%
\pgfsetdash{}{0pt}%
\pgfsys@defobject{currentmarker}{\pgfqpoint{0.000000in}{0.000000in}}{\pgfqpoint{0.000000in}{0.020833in}}{%
\pgfpathmoveto{\pgfqpoint{0.000000in}{0.000000in}}%
\pgfpathlineto{\pgfqpoint{0.000000in}{0.020833in}}%
\pgfusepath{stroke,fill}%
}%
\begin{pgfscope}%
\pgfsys@transformshift{1.552087in}{0.232212in}%
\pgfsys@useobject{currentmarker}{}%
\end{pgfscope}%
\end{pgfscope}%
\begin{pgfscope}%
\pgfsetbuttcap%
\pgfsetroundjoin%
\definecolor{currentfill}{rgb}{0.000000,0.000000,0.000000}%
\pgfsetfillcolor{currentfill}%
\pgfsetlinewidth{0.501875pt}%
\definecolor{currentstroke}{rgb}{0.000000,0.000000,0.000000}%
\pgfsetstrokecolor{currentstroke}%
\pgfsetdash{}{0pt}%
\pgfsys@defobject{currentmarker}{\pgfqpoint{0.000000in}{-0.020833in}}{\pgfqpoint{0.000000in}{0.000000in}}{%
\pgfpathmoveto{\pgfqpoint{0.000000in}{0.000000in}}%
\pgfpathlineto{\pgfqpoint{0.000000in}{-0.020833in}}%
\pgfusepath{stroke,fill}%
}%
\begin{pgfscope}%
\pgfsys@transformshift{1.552087in}{3.904501in}%
\pgfsys@useobject{currentmarker}{}%
\end{pgfscope}%
\end{pgfscope}%
\begin{pgfscope}%
\pgfsetbuttcap%
\pgfsetroundjoin%
\definecolor{currentfill}{rgb}{0.000000,0.000000,0.000000}%
\pgfsetfillcolor{currentfill}%
\pgfsetlinewidth{0.501875pt}%
\definecolor{currentstroke}{rgb}{0.000000,0.000000,0.000000}%
\pgfsetstrokecolor{currentstroke}%
\pgfsetdash{}{0pt}%
\pgfsys@defobject{currentmarker}{\pgfqpoint{0.000000in}{0.000000in}}{\pgfqpoint{0.000000in}{0.020833in}}{%
\pgfpathmoveto{\pgfqpoint{0.000000in}{0.000000in}}%
\pgfpathlineto{\pgfqpoint{0.000000in}{0.020833in}}%
\pgfusepath{stroke,fill}%
}%
\begin{pgfscope}%
\pgfsys@transformshift{1.745130in}{0.232212in}%
\pgfsys@useobject{currentmarker}{}%
\end{pgfscope}%
\end{pgfscope}%
\begin{pgfscope}%
\pgfsetbuttcap%
\pgfsetroundjoin%
\definecolor{currentfill}{rgb}{0.000000,0.000000,0.000000}%
\pgfsetfillcolor{currentfill}%
\pgfsetlinewidth{0.501875pt}%
\definecolor{currentstroke}{rgb}{0.000000,0.000000,0.000000}%
\pgfsetstrokecolor{currentstroke}%
\pgfsetdash{}{0pt}%
\pgfsys@defobject{currentmarker}{\pgfqpoint{0.000000in}{-0.020833in}}{\pgfqpoint{0.000000in}{0.000000in}}{%
\pgfpathmoveto{\pgfqpoint{0.000000in}{0.000000in}}%
\pgfpathlineto{\pgfqpoint{0.000000in}{-0.020833in}}%
\pgfusepath{stroke,fill}%
}%
\begin{pgfscope}%
\pgfsys@transformshift{1.745130in}{3.904501in}%
\pgfsys@useobject{currentmarker}{}%
\end{pgfscope}%
\end{pgfscope}%
\begin{pgfscope}%
\pgfsetbuttcap%
\pgfsetroundjoin%
\definecolor{currentfill}{rgb}{0.000000,0.000000,0.000000}%
\pgfsetfillcolor{currentfill}%
\pgfsetlinewidth{0.501875pt}%
\definecolor{currentstroke}{rgb}{0.000000,0.000000,0.000000}%
\pgfsetstrokecolor{currentstroke}%
\pgfsetdash{}{0pt}%
\pgfsys@defobject{currentmarker}{\pgfqpoint{0.000000in}{0.000000in}}{\pgfqpoint{0.000000in}{0.020833in}}{%
\pgfpathmoveto{\pgfqpoint{0.000000in}{0.000000in}}%
\pgfpathlineto{\pgfqpoint{0.000000in}{0.020833in}}%
\pgfusepath{stroke,fill}%
}%
\begin{pgfscope}%
\pgfsys@transformshift{1.938172in}{0.232212in}%
\pgfsys@useobject{currentmarker}{}%
\end{pgfscope}%
\end{pgfscope}%
\begin{pgfscope}%
\pgfsetbuttcap%
\pgfsetroundjoin%
\definecolor{currentfill}{rgb}{0.000000,0.000000,0.000000}%
\pgfsetfillcolor{currentfill}%
\pgfsetlinewidth{0.501875pt}%
\definecolor{currentstroke}{rgb}{0.000000,0.000000,0.000000}%
\pgfsetstrokecolor{currentstroke}%
\pgfsetdash{}{0pt}%
\pgfsys@defobject{currentmarker}{\pgfqpoint{0.000000in}{-0.020833in}}{\pgfqpoint{0.000000in}{0.000000in}}{%
\pgfpathmoveto{\pgfqpoint{0.000000in}{0.000000in}}%
\pgfpathlineto{\pgfqpoint{0.000000in}{-0.020833in}}%
\pgfusepath{stroke,fill}%
}%
\begin{pgfscope}%
\pgfsys@transformshift{1.938172in}{3.904501in}%
\pgfsys@useobject{currentmarker}{}%
\end{pgfscope}%
\end{pgfscope}%
\begin{pgfscope}%
\pgfsetbuttcap%
\pgfsetroundjoin%
\definecolor{currentfill}{rgb}{0.000000,0.000000,0.000000}%
\pgfsetfillcolor{currentfill}%
\pgfsetlinewidth{0.501875pt}%
\definecolor{currentstroke}{rgb}{0.000000,0.000000,0.000000}%
\pgfsetstrokecolor{currentstroke}%
\pgfsetdash{}{0pt}%
\pgfsys@defobject{currentmarker}{\pgfqpoint{0.000000in}{0.000000in}}{\pgfqpoint{0.000000in}{0.020833in}}{%
\pgfpathmoveto{\pgfqpoint{0.000000in}{0.000000in}}%
\pgfpathlineto{\pgfqpoint{0.000000in}{0.020833in}}%
\pgfusepath{stroke,fill}%
}%
\begin{pgfscope}%
\pgfsys@transformshift{2.324257in}{0.232212in}%
\pgfsys@useobject{currentmarker}{}%
\end{pgfscope}%
\end{pgfscope}%
\begin{pgfscope}%
\pgfsetbuttcap%
\pgfsetroundjoin%
\definecolor{currentfill}{rgb}{0.000000,0.000000,0.000000}%
\pgfsetfillcolor{currentfill}%
\pgfsetlinewidth{0.501875pt}%
\definecolor{currentstroke}{rgb}{0.000000,0.000000,0.000000}%
\pgfsetstrokecolor{currentstroke}%
\pgfsetdash{}{0pt}%
\pgfsys@defobject{currentmarker}{\pgfqpoint{0.000000in}{-0.020833in}}{\pgfqpoint{0.000000in}{0.000000in}}{%
\pgfpathmoveto{\pgfqpoint{0.000000in}{0.000000in}}%
\pgfpathlineto{\pgfqpoint{0.000000in}{-0.020833in}}%
\pgfusepath{stroke,fill}%
}%
\begin{pgfscope}%
\pgfsys@transformshift{2.324257in}{3.904501in}%
\pgfsys@useobject{currentmarker}{}%
\end{pgfscope}%
\end{pgfscope}%
\begin{pgfscope}%
\pgfsetbuttcap%
\pgfsetroundjoin%
\definecolor{currentfill}{rgb}{0.000000,0.000000,0.000000}%
\pgfsetfillcolor{currentfill}%
\pgfsetlinewidth{0.501875pt}%
\definecolor{currentstroke}{rgb}{0.000000,0.000000,0.000000}%
\pgfsetstrokecolor{currentstroke}%
\pgfsetdash{}{0pt}%
\pgfsys@defobject{currentmarker}{\pgfqpoint{0.000000in}{0.000000in}}{\pgfqpoint{0.000000in}{0.020833in}}{%
\pgfpathmoveto{\pgfqpoint{0.000000in}{0.000000in}}%
\pgfpathlineto{\pgfqpoint{0.000000in}{0.020833in}}%
\pgfusepath{stroke,fill}%
}%
\begin{pgfscope}%
\pgfsys@transformshift{2.517300in}{0.232212in}%
\pgfsys@useobject{currentmarker}{}%
\end{pgfscope}%
\end{pgfscope}%
\begin{pgfscope}%
\pgfsetbuttcap%
\pgfsetroundjoin%
\definecolor{currentfill}{rgb}{0.000000,0.000000,0.000000}%
\pgfsetfillcolor{currentfill}%
\pgfsetlinewidth{0.501875pt}%
\definecolor{currentstroke}{rgb}{0.000000,0.000000,0.000000}%
\pgfsetstrokecolor{currentstroke}%
\pgfsetdash{}{0pt}%
\pgfsys@defobject{currentmarker}{\pgfqpoint{0.000000in}{-0.020833in}}{\pgfqpoint{0.000000in}{0.000000in}}{%
\pgfpathmoveto{\pgfqpoint{0.000000in}{0.000000in}}%
\pgfpathlineto{\pgfqpoint{0.000000in}{-0.020833in}}%
\pgfusepath{stroke,fill}%
}%
\begin{pgfscope}%
\pgfsys@transformshift{2.517300in}{3.904501in}%
\pgfsys@useobject{currentmarker}{}%
\end{pgfscope}%
\end{pgfscope}%
\begin{pgfscope}%
\pgfsetbuttcap%
\pgfsetroundjoin%
\definecolor{currentfill}{rgb}{0.000000,0.000000,0.000000}%
\pgfsetfillcolor{currentfill}%
\pgfsetlinewidth{0.501875pt}%
\definecolor{currentstroke}{rgb}{0.000000,0.000000,0.000000}%
\pgfsetstrokecolor{currentstroke}%
\pgfsetdash{}{0pt}%
\pgfsys@defobject{currentmarker}{\pgfqpoint{0.000000in}{0.000000in}}{\pgfqpoint{0.000000in}{0.020833in}}{%
\pgfpathmoveto{\pgfqpoint{0.000000in}{0.000000in}}%
\pgfpathlineto{\pgfqpoint{0.000000in}{0.020833in}}%
\pgfusepath{stroke,fill}%
}%
\begin{pgfscope}%
\pgfsys@transformshift{2.710342in}{0.232212in}%
\pgfsys@useobject{currentmarker}{}%
\end{pgfscope}%
\end{pgfscope}%
\begin{pgfscope}%
\pgfsetbuttcap%
\pgfsetroundjoin%
\definecolor{currentfill}{rgb}{0.000000,0.000000,0.000000}%
\pgfsetfillcolor{currentfill}%
\pgfsetlinewidth{0.501875pt}%
\definecolor{currentstroke}{rgb}{0.000000,0.000000,0.000000}%
\pgfsetstrokecolor{currentstroke}%
\pgfsetdash{}{0pt}%
\pgfsys@defobject{currentmarker}{\pgfqpoint{0.000000in}{-0.020833in}}{\pgfqpoint{0.000000in}{0.000000in}}{%
\pgfpathmoveto{\pgfqpoint{0.000000in}{0.000000in}}%
\pgfpathlineto{\pgfqpoint{0.000000in}{-0.020833in}}%
\pgfusepath{stroke,fill}%
}%
\begin{pgfscope}%
\pgfsys@transformshift{2.710342in}{3.904501in}%
\pgfsys@useobject{currentmarker}{}%
\end{pgfscope}%
\end{pgfscope}%
\begin{pgfscope}%
\pgfsetbuttcap%
\pgfsetroundjoin%
\definecolor{currentfill}{rgb}{0.000000,0.000000,0.000000}%
\pgfsetfillcolor{currentfill}%
\pgfsetlinewidth{0.501875pt}%
\definecolor{currentstroke}{rgb}{0.000000,0.000000,0.000000}%
\pgfsetstrokecolor{currentstroke}%
\pgfsetdash{}{0pt}%
\pgfsys@defobject{currentmarker}{\pgfqpoint{0.000000in}{0.000000in}}{\pgfqpoint{0.000000in}{0.020833in}}{%
\pgfpathmoveto{\pgfqpoint{0.000000in}{0.000000in}}%
\pgfpathlineto{\pgfqpoint{0.000000in}{0.020833in}}%
\pgfusepath{stroke,fill}%
}%
\begin{pgfscope}%
\pgfsys@transformshift{2.903385in}{0.232212in}%
\pgfsys@useobject{currentmarker}{}%
\end{pgfscope}%
\end{pgfscope}%
\begin{pgfscope}%
\pgfsetbuttcap%
\pgfsetroundjoin%
\definecolor{currentfill}{rgb}{0.000000,0.000000,0.000000}%
\pgfsetfillcolor{currentfill}%
\pgfsetlinewidth{0.501875pt}%
\definecolor{currentstroke}{rgb}{0.000000,0.000000,0.000000}%
\pgfsetstrokecolor{currentstroke}%
\pgfsetdash{}{0pt}%
\pgfsys@defobject{currentmarker}{\pgfqpoint{0.000000in}{-0.020833in}}{\pgfqpoint{0.000000in}{0.000000in}}{%
\pgfpathmoveto{\pgfqpoint{0.000000in}{0.000000in}}%
\pgfpathlineto{\pgfqpoint{0.000000in}{-0.020833in}}%
\pgfusepath{stroke,fill}%
}%
\begin{pgfscope}%
\pgfsys@transformshift{2.903385in}{3.904501in}%
\pgfsys@useobject{currentmarker}{}%
\end{pgfscope}%
\end{pgfscope}%
\begin{pgfscope}%
\pgfsetbuttcap%
\pgfsetroundjoin%
\definecolor{currentfill}{rgb}{0.000000,0.000000,0.000000}%
\pgfsetfillcolor{currentfill}%
\pgfsetlinewidth{0.501875pt}%
\definecolor{currentstroke}{rgb}{0.000000,0.000000,0.000000}%
\pgfsetstrokecolor{currentstroke}%
\pgfsetdash{}{0pt}%
\pgfsys@defobject{currentmarker}{\pgfqpoint{0.000000in}{0.000000in}}{\pgfqpoint{0.000000in}{0.020833in}}{%
\pgfpathmoveto{\pgfqpoint{0.000000in}{0.000000in}}%
\pgfpathlineto{\pgfqpoint{0.000000in}{0.020833in}}%
\pgfusepath{stroke,fill}%
}%
\begin{pgfscope}%
\pgfsys@transformshift{3.096427in}{0.232212in}%
\pgfsys@useobject{currentmarker}{}%
\end{pgfscope}%
\end{pgfscope}%
\begin{pgfscope}%
\pgfsetbuttcap%
\pgfsetroundjoin%
\definecolor{currentfill}{rgb}{0.000000,0.000000,0.000000}%
\pgfsetfillcolor{currentfill}%
\pgfsetlinewidth{0.501875pt}%
\definecolor{currentstroke}{rgb}{0.000000,0.000000,0.000000}%
\pgfsetstrokecolor{currentstroke}%
\pgfsetdash{}{0pt}%
\pgfsys@defobject{currentmarker}{\pgfqpoint{0.000000in}{-0.020833in}}{\pgfqpoint{0.000000in}{0.000000in}}{%
\pgfpathmoveto{\pgfqpoint{0.000000in}{0.000000in}}%
\pgfpathlineto{\pgfqpoint{0.000000in}{-0.020833in}}%
\pgfusepath{stroke,fill}%
}%
\begin{pgfscope}%
\pgfsys@transformshift{3.096427in}{3.904501in}%
\pgfsys@useobject{currentmarker}{}%
\end{pgfscope}%
\end{pgfscope}%
\begin{pgfscope}%
\pgfsetbuttcap%
\pgfsetroundjoin%
\definecolor{currentfill}{rgb}{0.000000,0.000000,0.000000}%
\pgfsetfillcolor{currentfill}%
\pgfsetlinewidth{0.501875pt}%
\definecolor{currentstroke}{rgb}{0.000000,0.000000,0.000000}%
\pgfsetstrokecolor{currentstroke}%
\pgfsetdash{}{0pt}%
\pgfsys@defobject{currentmarker}{\pgfqpoint{0.000000in}{0.000000in}}{\pgfqpoint{0.000000in}{0.020833in}}{%
\pgfpathmoveto{\pgfqpoint{0.000000in}{0.000000in}}%
\pgfpathlineto{\pgfqpoint{0.000000in}{0.020833in}}%
\pgfusepath{stroke,fill}%
}%
\begin{pgfscope}%
\pgfsys@transformshift{3.289470in}{0.232212in}%
\pgfsys@useobject{currentmarker}{}%
\end{pgfscope}%
\end{pgfscope}%
\begin{pgfscope}%
\pgfsetbuttcap%
\pgfsetroundjoin%
\definecolor{currentfill}{rgb}{0.000000,0.000000,0.000000}%
\pgfsetfillcolor{currentfill}%
\pgfsetlinewidth{0.501875pt}%
\definecolor{currentstroke}{rgb}{0.000000,0.000000,0.000000}%
\pgfsetstrokecolor{currentstroke}%
\pgfsetdash{}{0pt}%
\pgfsys@defobject{currentmarker}{\pgfqpoint{0.000000in}{-0.020833in}}{\pgfqpoint{0.000000in}{0.000000in}}{%
\pgfpathmoveto{\pgfqpoint{0.000000in}{0.000000in}}%
\pgfpathlineto{\pgfqpoint{0.000000in}{-0.020833in}}%
\pgfusepath{stroke,fill}%
}%
\begin{pgfscope}%
\pgfsys@transformshift{3.289470in}{3.904501in}%
\pgfsys@useobject{currentmarker}{}%
\end{pgfscope}%
\end{pgfscope}%
\begin{pgfscope}%
\pgfsetbuttcap%
\pgfsetroundjoin%
\definecolor{currentfill}{rgb}{0.000000,0.000000,0.000000}%
\pgfsetfillcolor{currentfill}%
\pgfsetlinewidth{0.501875pt}%
\definecolor{currentstroke}{rgb}{0.000000,0.000000,0.000000}%
\pgfsetstrokecolor{currentstroke}%
\pgfsetdash{}{0pt}%
\pgfsys@defobject{currentmarker}{\pgfqpoint{0.000000in}{0.000000in}}{\pgfqpoint{0.000000in}{0.020833in}}{%
\pgfpathmoveto{\pgfqpoint{0.000000in}{0.000000in}}%
\pgfpathlineto{\pgfqpoint{0.000000in}{0.020833in}}%
\pgfusepath{stroke,fill}%
}%
\begin{pgfscope}%
\pgfsys@transformshift{3.482512in}{0.232212in}%
\pgfsys@useobject{currentmarker}{}%
\end{pgfscope}%
\end{pgfscope}%
\begin{pgfscope}%
\pgfsetbuttcap%
\pgfsetroundjoin%
\definecolor{currentfill}{rgb}{0.000000,0.000000,0.000000}%
\pgfsetfillcolor{currentfill}%
\pgfsetlinewidth{0.501875pt}%
\definecolor{currentstroke}{rgb}{0.000000,0.000000,0.000000}%
\pgfsetstrokecolor{currentstroke}%
\pgfsetdash{}{0pt}%
\pgfsys@defobject{currentmarker}{\pgfqpoint{0.000000in}{-0.020833in}}{\pgfqpoint{0.000000in}{0.000000in}}{%
\pgfpathmoveto{\pgfqpoint{0.000000in}{0.000000in}}%
\pgfpathlineto{\pgfqpoint{0.000000in}{-0.020833in}}%
\pgfusepath{stroke,fill}%
}%
\begin{pgfscope}%
\pgfsys@transformshift{3.482512in}{3.904501in}%
\pgfsys@useobject{currentmarker}{}%
\end{pgfscope}%
\end{pgfscope}%
\begin{pgfscope}%
\pgfsetbuttcap%
\pgfsetroundjoin%
\definecolor{currentfill}{rgb}{0.000000,0.000000,0.000000}%
\pgfsetfillcolor{currentfill}%
\pgfsetlinewidth{0.501875pt}%
\definecolor{currentstroke}{rgb}{0.000000,0.000000,0.000000}%
\pgfsetstrokecolor{currentstroke}%
\pgfsetdash{}{0pt}%
\pgfsys@defobject{currentmarker}{\pgfqpoint{0.000000in}{0.000000in}}{\pgfqpoint{0.000000in}{0.020833in}}{%
\pgfpathmoveto{\pgfqpoint{0.000000in}{0.000000in}}%
\pgfpathlineto{\pgfqpoint{0.000000in}{0.020833in}}%
\pgfusepath{stroke,fill}%
}%
\begin{pgfscope}%
\pgfsys@transformshift{3.675555in}{0.232212in}%
\pgfsys@useobject{currentmarker}{}%
\end{pgfscope}%
\end{pgfscope}%
\begin{pgfscope}%
\pgfsetbuttcap%
\pgfsetroundjoin%
\definecolor{currentfill}{rgb}{0.000000,0.000000,0.000000}%
\pgfsetfillcolor{currentfill}%
\pgfsetlinewidth{0.501875pt}%
\definecolor{currentstroke}{rgb}{0.000000,0.000000,0.000000}%
\pgfsetstrokecolor{currentstroke}%
\pgfsetdash{}{0pt}%
\pgfsys@defobject{currentmarker}{\pgfqpoint{0.000000in}{-0.020833in}}{\pgfqpoint{0.000000in}{0.000000in}}{%
\pgfpathmoveto{\pgfqpoint{0.000000in}{0.000000in}}%
\pgfpathlineto{\pgfqpoint{0.000000in}{-0.020833in}}%
\pgfusepath{stroke,fill}%
}%
\begin{pgfscope}%
\pgfsys@transformshift{3.675555in}{3.904501in}%
\pgfsys@useobject{currentmarker}{}%
\end{pgfscope}%
\end{pgfscope}%
\begin{pgfscope}%
\pgfsetbuttcap%
\pgfsetroundjoin%
\definecolor{currentfill}{rgb}{0.000000,0.000000,0.000000}%
\pgfsetfillcolor{currentfill}%
\pgfsetlinewidth{0.501875pt}%
\definecolor{currentstroke}{rgb}{0.000000,0.000000,0.000000}%
\pgfsetstrokecolor{currentstroke}%
\pgfsetdash{}{0pt}%
\pgfsys@defobject{currentmarker}{\pgfqpoint{0.000000in}{0.000000in}}{\pgfqpoint{0.000000in}{0.020833in}}{%
\pgfpathmoveto{\pgfqpoint{0.000000in}{0.000000in}}%
\pgfpathlineto{\pgfqpoint{0.000000in}{0.020833in}}%
\pgfusepath{stroke,fill}%
}%
\begin{pgfscope}%
\pgfsys@transformshift{3.868597in}{0.232212in}%
\pgfsys@useobject{currentmarker}{}%
\end{pgfscope}%
\end{pgfscope}%
\begin{pgfscope}%
\pgfsetbuttcap%
\pgfsetroundjoin%
\definecolor{currentfill}{rgb}{0.000000,0.000000,0.000000}%
\pgfsetfillcolor{currentfill}%
\pgfsetlinewidth{0.501875pt}%
\definecolor{currentstroke}{rgb}{0.000000,0.000000,0.000000}%
\pgfsetstrokecolor{currentstroke}%
\pgfsetdash{}{0pt}%
\pgfsys@defobject{currentmarker}{\pgfqpoint{0.000000in}{-0.020833in}}{\pgfqpoint{0.000000in}{0.000000in}}{%
\pgfpathmoveto{\pgfqpoint{0.000000in}{0.000000in}}%
\pgfpathlineto{\pgfqpoint{0.000000in}{-0.020833in}}%
\pgfusepath{stroke,fill}%
}%
\begin{pgfscope}%
\pgfsys@transformshift{3.868597in}{3.904501in}%
\pgfsys@useobject{currentmarker}{}%
\end{pgfscope}%
\end{pgfscope}%
\begin{pgfscope}%
\pgfsetbuttcap%
\pgfsetroundjoin%
\definecolor{currentfill}{rgb}{0.000000,0.000000,0.000000}%
\pgfsetfillcolor{currentfill}%
\pgfsetlinewidth{0.501875pt}%
\definecolor{currentstroke}{rgb}{0.000000,0.000000,0.000000}%
\pgfsetstrokecolor{currentstroke}%
\pgfsetdash{}{0pt}%
\pgfsys@defobject{currentmarker}{\pgfqpoint{0.000000in}{0.000000in}}{\pgfqpoint{0.000000in}{0.020833in}}{%
\pgfpathmoveto{\pgfqpoint{0.000000in}{0.000000in}}%
\pgfpathlineto{\pgfqpoint{0.000000in}{0.020833in}}%
\pgfusepath{stroke,fill}%
}%
\begin{pgfscope}%
\pgfsys@transformshift{4.061640in}{0.232212in}%
\pgfsys@useobject{currentmarker}{}%
\end{pgfscope}%
\end{pgfscope}%
\begin{pgfscope}%
\pgfsetbuttcap%
\pgfsetroundjoin%
\definecolor{currentfill}{rgb}{0.000000,0.000000,0.000000}%
\pgfsetfillcolor{currentfill}%
\pgfsetlinewidth{0.501875pt}%
\definecolor{currentstroke}{rgb}{0.000000,0.000000,0.000000}%
\pgfsetstrokecolor{currentstroke}%
\pgfsetdash{}{0pt}%
\pgfsys@defobject{currentmarker}{\pgfqpoint{0.000000in}{-0.020833in}}{\pgfqpoint{0.000000in}{0.000000in}}{%
\pgfpathmoveto{\pgfqpoint{0.000000in}{0.000000in}}%
\pgfpathlineto{\pgfqpoint{0.000000in}{-0.020833in}}%
\pgfusepath{stroke,fill}%
}%
\begin{pgfscope}%
\pgfsys@transformshift{4.061640in}{3.904501in}%
\pgfsys@useobject{currentmarker}{}%
\end{pgfscope}%
\end{pgfscope}%
\begin{pgfscope}%
\pgfsetbuttcap%
\pgfsetroundjoin%
\definecolor{currentfill}{rgb}{0.000000,0.000000,0.000000}%
\pgfsetfillcolor{currentfill}%
\pgfsetlinewidth{0.501875pt}%
\definecolor{currentstroke}{rgb}{0.000000,0.000000,0.000000}%
\pgfsetstrokecolor{currentstroke}%
\pgfsetdash{}{0pt}%
\pgfsys@defobject{currentmarker}{\pgfqpoint{0.000000in}{0.000000in}}{\pgfqpoint{0.000000in}{0.020833in}}{%
\pgfpathmoveto{\pgfqpoint{0.000000in}{0.000000in}}%
\pgfpathlineto{\pgfqpoint{0.000000in}{0.020833in}}%
\pgfusepath{stroke,fill}%
}%
\begin{pgfscope}%
\pgfsys@transformshift{4.254682in}{0.232212in}%
\pgfsys@useobject{currentmarker}{}%
\end{pgfscope}%
\end{pgfscope}%
\begin{pgfscope}%
\pgfsetbuttcap%
\pgfsetroundjoin%
\definecolor{currentfill}{rgb}{0.000000,0.000000,0.000000}%
\pgfsetfillcolor{currentfill}%
\pgfsetlinewidth{0.501875pt}%
\definecolor{currentstroke}{rgb}{0.000000,0.000000,0.000000}%
\pgfsetstrokecolor{currentstroke}%
\pgfsetdash{}{0pt}%
\pgfsys@defobject{currentmarker}{\pgfqpoint{0.000000in}{-0.020833in}}{\pgfqpoint{0.000000in}{0.000000in}}{%
\pgfpathmoveto{\pgfqpoint{0.000000in}{0.000000in}}%
\pgfpathlineto{\pgfqpoint{0.000000in}{-0.020833in}}%
\pgfusepath{stroke,fill}%
}%
\begin{pgfscope}%
\pgfsys@transformshift{4.254682in}{3.904501in}%
\pgfsys@useobject{currentmarker}{}%
\end{pgfscope}%
\end{pgfscope}%
\begin{pgfscope}%
\pgfsetbuttcap%
\pgfsetroundjoin%
\definecolor{currentfill}{rgb}{0.000000,0.000000,0.000000}%
\pgfsetfillcolor{currentfill}%
\pgfsetlinewidth{0.501875pt}%
\definecolor{currentstroke}{rgb}{0.000000,0.000000,0.000000}%
\pgfsetstrokecolor{currentstroke}%
\pgfsetdash{}{0pt}%
\pgfsys@defobject{currentmarker}{\pgfqpoint{0.000000in}{0.000000in}}{\pgfqpoint{0.000000in}{0.020833in}}{%
\pgfpathmoveto{\pgfqpoint{0.000000in}{0.000000in}}%
\pgfpathlineto{\pgfqpoint{0.000000in}{0.020833in}}%
\pgfusepath{stroke,fill}%
}%
\begin{pgfscope}%
\pgfsys@transformshift{4.447725in}{0.232212in}%
\pgfsys@useobject{currentmarker}{}%
\end{pgfscope}%
\end{pgfscope}%
\begin{pgfscope}%
\pgfsetbuttcap%
\pgfsetroundjoin%
\definecolor{currentfill}{rgb}{0.000000,0.000000,0.000000}%
\pgfsetfillcolor{currentfill}%
\pgfsetlinewidth{0.501875pt}%
\definecolor{currentstroke}{rgb}{0.000000,0.000000,0.000000}%
\pgfsetstrokecolor{currentstroke}%
\pgfsetdash{}{0pt}%
\pgfsys@defobject{currentmarker}{\pgfqpoint{0.000000in}{-0.020833in}}{\pgfqpoint{0.000000in}{0.000000in}}{%
\pgfpathmoveto{\pgfqpoint{0.000000in}{0.000000in}}%
\pgfpathlineto{\pgfqpoint{0.000000in}{-0.020833in}}%
\pgfusepath{stroke,fill}%
}%
\begin{pgfscope}%
\pgfsys@transformshift{4.447725in}{3.904501in}%
\pgfsys@useobject{currentmarker}{}%
\end{pgfscope}%
\end{pgfscope}%
\begin{pgfscope}%
\pgfsetbuttcap%
\pgfsetroundjoin%
\definecolor{currentfill}{rgb}{0.000000,0.000000,0.000000}%
\pgfsetfillcolor{currentfill}%
\pgfsetlinewidth{0.501875pt}%
\definecolor{currentstroke}{rgb}{0.000000,0.000000,0.000000}%
\pgfsetstrokecolor{currentstroke}%
\pgfsetdash{}{0pt}%
\pgfsys@defobject{currentmarker}{\pgfqpoint{0.000000in}{0.000000in}}{\pgfqpoint{0.000000in}{0.020833in}}{%
\pgfpathmoveto{\pgfqpoint{0.000000in}{0.000000in}}%
\pgfpathlineto{\pgfqpoint{0.000000in}{0.020833in}}%
\pgfusepath{stroke,fill}%
}%
\begin{pgfscope}%
\pgfsys@transformshift{4.640767in}{0.232212in}%
\pgfsys@useobject{currentmarker}{}%
\end{pgfscope}%
\end{pgfscope}%
\begin{pgfscope}%
\pgfsetbuttcap%
\pgfsetroundjoin%
\definecolor{currentfill}{rgb}{0.000000,0.000000,0.000000}%
\pgfsetfillcolor{currentfill}%
\pgfsetlinewidth{0.501875pt}%
\definecolor{currentstroke}{rgb}{0.000000,0.000000,0.000000}%
\pgfsetstrokecolor{currentstroke}%
\pgfsetdash{}{0pt}%
\pgfsys@defobject{currentmarker}{\pgfqpoint{0.000000in}{-0.020833in}}{\pgfqpoint{0.000000in}{0.000000in}}{%
\pgfpathmoveto{\pgfqpoint{0.000000in}{0.000000in}}%
\pgfpathlineto{\pgfqpoint{0.000000in}{-0.020833in}}%
\pgfusepath{stroke,fill}%
}%
\begin{pgfscope}%
\pgfsys@transformshift{4.640767in}{3.904501in}%
\pgfsys@useobject{currentmarker}{}%
\end{pgfscope}%
\end{pgfscope}%
\begin{pgfscope}%
\pgfsetbuttcap%
\pgfsetroundjoin%
\definecolor{currentfill}{rgb}{0.000000,0.000000,0.000000}%
\pgfsetfillcolor{currentfill}%
\pgfsetlinewidth{0.501875pt}%
\definecolor{currentstroke}{rgb}{0.000000,0.000000,0.000000}%
\pgfsetstrokecolor{currentstroke}%
\pgfsetdash{}{0pt}%
\pgfsys@defobject{currentmarker}{\pgfqpoint{0.000000in}{0.000000in}}{\pgfqpoint{0.000000in}{0.020833in}}{%
\pgfpathmoveto{\pgfqpoint{0.000000in}{0.000000in}}%
\pgfpathlineto{\pgfqpoint{0.000000in}{0.020833in}}%
\pgfusepath{stroke,fill}%
}%
\begin{pgfscope}%
\pgfsys@transformshift{4.833810in}{0.232212in}%
\pgfsys@useobject{currentmarker}{}%
\end{pgfscope}%
\end{pgfscope}%
\begin{pgfscope}%
\pgfsetbuttcap%
\pgfsetroundjoin%
\definecolor{currentfill}{rgb}{0.000000,0.000000,0.000000}%
\pgfsetfillcolor{currentfill}%
\pgfsetlinewidth{0.501875pt}%
\definecolor{currentstroke}{rgb}{0.000000,0.000000,0.000000}%
\pgfsetstrokecolor{currentstroke}%
\pgfsetdash{}{0pt}%
\pgfsys@defobject{currentmarker}{\pgfqpoint{0.000000in}{-0.020833in}}{\pgfqpoint{0.000000in}{0.000000in}}{%
\pgfpathmoveto{\pgfqpoint{0.000000in}{0.000000in}}%
\pgfpathlineto{\pgfqpoint{0.000000in}{-0.020833in}}%
\pgfusepath{stroke,fill}%
}%
\begin{pgfscope}%
\pgfsys@transformshift{4.833810in}{3.904501in}%
\pgfsys@useobject{currentmarker}{}%
\end{pgfscope}%
\end{pgfscope}%
\begin{pgfscope}%
\pgfsetbuttcap%
\pgfsetroundjoin%
\definecolor{currentfill}{rgb}{0.000000,0.000000,0.000000}%
\pgfsetfillcolor{currentfill}%
\pgfsetlinewidth{0.501875pt}%
\definecolor{currentstroke}{rgb}{0.000000,0.000000,0.000000}%
\pgfsetstrokecolor{currentstroke}%
\pgfsetdash{}{0pt}%
\pgfsys@defobject{currentmarker}{\pgfqpoint{0.000000in}{0.000000in}}{\pgfqpoint{0.000000in}{0.020833in}}{%
\pgfpathmoveto{\pgfqpoint{0.000000in}{0.000000in}}%
\pgfpathlineto{\pgfqpoint{0.000000in}{0.020833in}}%
\pgfusepath{stroke,fill}%
}%
\begin{pgfscope}%
\pgfsys@transformshift{5.026852in}{0.232212in}%
\pgfsys@useobject{currentmarker}{}%
\end{pgfscope}%
\end{pgfscope}%
\begin{pgfscope}%
\pgfsetbuttcap%
\pgfsetroundjoin%
\definecolor{currentfill}{rgb}{0.000000,0.000000,0.000000}%
\pgfsetfillcolor{currentfill}%
\pgfsetlinewidth{0.501875pt}%
\definecolor{currentstroke}{rgb}{0.000000,0.000000,0.000000}%
\pgfsetstrokecolor{currentstroke}%
\pgfsetdash{}{0pt}%
\pgfsys@defobject{currentmarker}{\pgfqpoint{0.000000in}{-0.020833in}}{\pgfqpoint{0.000000in}{0.000000in}}{%
\pgfpathmoveto{\pgfqpoint{0.000000in}{0.000000in}}%
\pgfpathlineto{\pgfqpoint{0.000000in}{-0.020833in}}%
\pgfusepath{stroke,fill}%
}%
\begin{pgfscope}%
\pgfsys@transformshift{5.026852in}{3.904501in}%
\pgfsys@useobject{currentmarker}{}%
\end{pgfscope}%
\end{pgfscope}%
\begin{pgfscope}%
\pgfsetbuttcap%
\pgfsetroundjoin%
\definecolor{currentfill}{rgb}{0.000000,0.000000,0.000000}%
\pgfsetfillcolor{currentfill}%
\pgfsetlinewidth{0.501875pt}%
\definecolor{currentstroke}{rgb}{0.000000,0.000000,0.000000}%
\pgfsetstrokecolor{currentstroke}%
\pgfsetdash{}{0pt}%
\pgfsys@defobject{currentmarker}{\pgfqpoint{0.000000in}{0.000000in}}{\pgfqpoint{0.000000in}{0.020833in}}{%
\pgfpathmoveto{\pgfqpoint{0.000000in}{0.000000in}}%
\pgfpathlineto{\pgfqpoint{0.000000in}{0.020833in}}%
\pgfusepath{stroke,fill}%
}%
\begin{pgfscope}%
\pgfsys@transformshift{5.219895in}{0.232212in}%
\pgfsys@useobject{currentmarker}{}%
\end{pgfscope}%
\end{pgfscope}%
\begin{pgfscope}%
\pgfsetbuttcap%
\pgfsetroundjoin%
\definecolor{currentfill}{rgb}{0.000000,0.000000,0.000000}%
\pgfsetfillcolor{currentfill}%
\pgfsetlinewidth{0.501875pt}%
\definecolor{currentstroke}{rgb}{0.000000,0.000000,0.000000}%
\pgfsetstrokecolor{currentstroke}%
\pgfsetdash{}{0pt}%
\pgfsys@defobject{currentmarker}{\pgfqpoint{0.000000in}{-0.020833in}}{\pgfqpoint{0.000000in}{0.000000in}}{%
\pgfpathmoveto{\pgfqpoint{0.000000in}{0.000000in}}%
\pgfpathlineto{\pgfqpoint{0.000000in}{-0.020833in}}%
\pgfusepath{stroke,fill}%
}%
\begin{pgfscope}%
\pgfsys@transformshift{5.219895in}{3.904501in}%
\pgfsys@useobject{currentmarker}{}%
\end{pgfscope}%
\end{pgfscope}%
\begin{pgfscope}%
\pgfsetbuttcap%
\pgfsetroundjoin%
\definecolor{currentfill}{rgb}{0.000000,0.000000,0.000000}%
\pgfsetfillcolor{currentfill}%
\pgfsetlinewidth{0.501875pt}%
\definecolor{currentstroke}{rgb}{0.000000,0.000000,0.000000}%
\pgfsetstrokecolor{currentstroke}%
\pgfsetdash{}{0pt}%
\pgfsys@defobject{currentmarker}{\pgfqpoint{0.000000in}{0.000000in}}{\pgfqpoint{0.000000in}{0.020833in}}{%
\pgfpathmoveto{\pgfqpoint{0.000000in}{0.000000in}}%
\pgfpathlineto{\pgfqpoint{0.000000in}{0.020833in}}%
\pgfusepath{stroke,fill}%
}%
\begin{pgfscope}%
\pgfsys@transformshift{5.412937in}{0.232212in}%
\pgfsys@useobject{currentmarker}{}%
\end{pgfscope}%
\end{pgfscope}%
\begin{pgfscope}%
\pgfsetbuttcap%
\pgfsetroundjoin%
\definecolor{currentfill}{rgb}{0.000000,0.000000,0.000000}%
\pgfsetfillcolor{currentfill}%
\pgfsetlinewidth{0.501875pt}%
\definecolor{currentstroke}{rgb}{0.000000,0.000000,0.000000}%
\pgfsetstrokecolor{currentstroke}%
\pgfsetdash{}{0pt}%
\pgfsys@defobject{currentmarker}{\pgfqpoint{0.000000in}{-0.020833in}}{\pgfqpoint{0.000000in}{0.000000in}}{%
\pgfpathmoveto{\pgfqpoint{0.000000in}{0.000000in}}%
\pgfpathlineto{\pgfqpoint{0.000000in}{-0.020833in}}%
\pgfusepath{stroke,fill}%
}%
\begin{pgfscope}%
\pgfsys@transformshift{5.412937in}{3.904501in}%
\pgfsys@useobject{currentmarker}{}%
\end{pgfscope}%
\end{pgfscope}%
\begin{pgfscope}%
\pgfsetbuttcap%
\pgfsetroundjoin%
\definecolor{currentfill}{rgb}{0.000000,0.000000,0.000000}%
\pgfsetfillcolor{currentfill}%
\pgfsetlinewidth{0.501875pt}%
\definecolor{currentstroke}{rgb}{0.000000,0.000000,0.000000}%
\pgfsetstrokecolor{currentstroke}%
\pgfsetdash{}{0pt}%
\pgfsys@defobject{currentmarker}{\pgfqpoint{0.000000in}{0.000000in}}{\pgfqpoint{0.000000in}{0.020833in}}{%
\pgfpathmoveto{\pgfqpoint{0.000000in}{0.000000in}}%
\pgfpathlineto{\pgfqpoint{0.000000in}{0.020833in}}%
\pgfusepath{stroke,fill}%
}%
\begin{pgfscope}%
\pgfsys@transformshift{5.605980in}{0.232212in}%
\pgfsys@useobject{currentmarker}{}%
\end{pgfscope}%
\end{pgfscope}%
\begin{pgfscope}%
\pgfsetbuttcap%
\pgfsetroundjoin%
\definecolor{currentfill}{rgb}{0.000000,0.000000,0.000000}%
\pgfsetfillcolor{currentfill}%
\pgfsetlinewidth{0.501875pt}%
\definecolor{currentstroke}{rgb}{0.000000,0.000000,0.000000}%
\pgfsetstrokecolor{currentstroke}%
\pgfsetdash{}{0pt}%
\pgfsys@defobject{currentmarker}{\pgfqpoint{0.000000in}{-0.020833in}}{\pgfqpoint{0.000000in}{0.000000in}}{%
\pgfpathmoveto{\pgfqpoint{0.000000in}{0.000000in}}%
\pgfpathlineto{\pgfqpoint{0.000000in}{-0.020833in}}%
\pgfusepath{stroke,fill}%
}%
\begin{pgfscope}%
\pgfsys@transformshift{5.605980in}{3.904501in}%
\pgfsys@useobject{currentmarker}{}%
\end{pgfscope}%
\end{pgfscope}%
\begin{pgfscope}%
\pgfsetbuttcap%
\pgfsetroundjoin%
\definecolor{currentfill}{rgb}{0.000000,0.000000,0.000000}%
\pgfsetfillcolor{currentfill}%
\pgfsetlinewidth{0.501875pt}%
\definecolor{currentstroke}{rgb}{0.000000,0.000000,0.000000}%
\pgfsetstrokecolor{currentstroke}%
\pgfsetdash{}{0pt}%
\pgfsys@defobject{currentmarker}{\pgfqpoint{0.000000in}{0.000000in}}{\pgfqpoint{0.000000in}{0.020833in}}{%
\pgfpathmoveto{\pgfqpoint{0.000000in}{0.000000in}}%
\pgfpathlineto{\pgfqpoint{0.000000in}{0.020833in}}%
\pgfusepath{stroke,fill}%
}%
\begin{pgfscope}%
\pgfsys@transformshift{5.799022in}{0.232212in}%
\pgfsys@useobject{currentmarker}{}%
\end{pgfscope}%
\end{pgfscope}%
\begin{pgfscope}%
\pgfsetbuttcap%
\pgfsetroundjoin%
\definecolor{currentfill}{rgb}{0.000000,0.000000,0.000000}%
\pgfsetfillcolor{currentfill}%
\pgfsetlinewidth{0.501875pt}%
\definecolor{currentstroke}{rgb}{0.000000,0.000000,0.000000}%
\pgfsetstrokecolor{currentstroke}%
\pgfsetdash{}{0pt}%
\pgfsys@defobject{currentmarker}{\pgfqpoint{0.000000in}{-0.020833in}}{\pgfqpoint{0.000000in}{0.000000in}}{%
\pgfpathmoveto{\pgfqpoint{0.000000in}{0.000000in}}%
\pgfpathlineto{\pgfqpoint{0.000000in}{-0.020833in}}%
\pgfusepath{stroke,fill}%
}%
\begin{pgfscope}%
\pgfsys@transformshift{5.799022in}{3.904501in}%
\pgfsys@useobject{currentmarker}{}%
\end{pgfscope}%
\end{pgfscope}%
\begin{pgfscope}%
\pgfsetbuttcap%
\pgfsetroundjoin%
\definecolor{currentfill}{rgb}{0.000000,0.000000,0.000000}%
\pgfsetfillcolor{currentfill}%
\pgfsetlinewidth{0.501875pt}%
\definecolor{currentstroke}{rgb}{0.000000,0.000000,0.000000}%
\pgfsetstrokecolor{currentstroke}%
\pgfsetdash{}{0pt}%
\pgfsys@defobject{currentmarker}{\pgfqpoint{0.000000in}{0.000000in}}{\pgfqpoint{0.000000in}{0.020833in}}{%
\pgfpathmoveto{\pgfqpoint{0.000000in}{0.000000in}}%
\pgfpathlineto{\pgfqpoint{0.000000in}{0.020833in}}%
\pgfusepath{stroke,fill}%
}%
\begin{pgfscope}%
\pgfsys@transformshift{5.992065in}{0.232212in}%
\pgfsys@useobject{currentmarker}{}%
\end{pgfscope}%
\end{pgfscope}%
\begin{pgfscope}%
\pgfsetbuttcap%
\pgfsetroundjoin%
\definecolor{currentfill}{rgb}{0.000000,0.000000,0.000000}%
\pgfsetfillcolor{currentfill}%
\pgfsetlinewidth{0.501875pt}%
\definecolor{currentstroke}{rgb}{0.000000,0.000000,0.000000}%
\pgfsetstrokecolor{currentstroke}%
\pgfsetdash{}{0pt}%
\pgfsys@defobject{currentmarker}{\pgfqpoint{0.000000in}{-0.020833in}}{\pgfqpoint{0.000000in}{0.000000in}}{%
\pgfpathmoveto{\pgfqpoint{0.000000in}{0.000000in}}%
\pgfpathlineto{\pgfqpoint{0.000000in}{-0.020833in}}%
\pgfusepath{stroke,fill}%
}%
\begin{pgfscope}%
\pgfsys@transformshift{5.992065in}{3.904501in}%
\pgfsys@useobject{currentmarker}{}%
\end{pgfscope}%
\end{pgfscope}%
\begin{pgfscope}%
\definecolor{textcolor}{rgb}{0.000000,0.000000,0.000000}%
\pgfsetstrokecolor{textcolor}%
\pgfsetfillcolor{textcolor}%
\pgftext[x=6.171885in,y=0.192732in,right,top]{\color{textcolor}{\rmfamily\fontsize{11.000000}{13.200000}\selectfont\catcode`\^=\active\def^{\ifmmode\sp\else\^{}\fi}\catcode`\%=\active\def%{\%}$x$}}%
\end{pgfscope}%
\begin{pgfscope}%
\pgfsetbuttcap%
\pgfsetroundjoin%
\definecolor{currentfill}{rgb}{0.000000,0.000000,0.000000}%
\pgfsetfillcolor{currentfill}%
\pgfsetlinewidth{0.501875pt}%
\definecolor{currentstroke}{rgb}{0.000000,0.000000,0.000000}%
\pgfsetstrokecolor{currentstroke}%
\pgfsetdash{}{0pt}%
\pgfsys@defobject{currentmarker}{\pgfqpoint{0.000000in}{0.000000in}}{\pgfqpoint{0.041667in}{0.000000in}}{%
\pgfpathmoveto{\pgfqpoint{0.000000in}{0.000000in}}%
\pgfpathlineto{\pgfqpoint{0.041667in}{0.000000in}}%
\pgfusepath{stroke,fill}%
}%
\begin{pgfscope}%
\pgfsys@transformshift{0.396477in}{2.855276in}%
\pgfsys@useobject{currentmarker}{}%
\end{pgfscope}%
\end{pgfscope}%
\begin{pgfscope}%
\pgfsetbuttcap%
\pgfsetroundjoin%
\definecolor{currentfill}{rgb}{0.000000,0.000000,0.000000}%
\pgfsetfillcolor{currentfill}%
\pgfsetlinewidth{0.501875pt}%
\definecolor{currentstroke}{rgb}{0.000000,0.000000,0.000000}%
\pgfsetstrokecolor{currentstroke}%
\pgfsetdash{}{0pt}%
\pgfsys@defobject{currentmarker}{\pgfqpoint{-0.041667in}{0.000000in}}{\pgfqpoint{-0.000000in}{0.000000in}}{%
\pgfpathmoveto{\pgfqpoint{-0.000000in}{0.000000in}}%
\pgfpathlineto{\pgfqpoint{-0.041667in}{0.000000in}}%
\pgfusepath{stroke,fill}%
}%
\begin{pgfscope}%
\pgfsys@transformshift{6.171885in}{2.855276in}%
\pgfsys@useobject{currentmarker}{}%
\end{pgfscope}%
\end{pgfscope}%
\begin{pgfscope}%
\definecolor{textcolor}{rgb}{0.000000,0.000000,0.000000}%
\pgfsetstrokecolor{textcolor}%
\pgfsetfillcolor{textcolor}%
\pgftext[x=0.249671in, y=2.806578in, left, base]{\color{textcolor}{\rmfamily\fontsize{10.000000}{12.000000}\selectfont\catcode`\^=\active\def^{\ifmmode\sp\else\^{}\fi}\catcode`\%=\active\def%{\%}$K$}}%
\end{pgfscope}%
\begin{pgfscope}%
\definecolor{textcolor}{rgb}{0.000000,0.000000,0.000000}%
\pgfsetstrokecolor{textcolor}%
\pgfsetfillcolor{textcolor}%
\pgftext[x=0.388560in,y=3.904501in,right,bottom]{\color{textcolor}{\rmfamily\fontsize{11.000000}{13.200000}\selectfont\catcode`\^=\active\def^{\ifmmode\sp\else\^{}\fi}\catcode`\%=\active\def%{\%}$v_L(x)$}}%
\end{pgfscope}%
\begin{pgfscope}%
\pgfpathrectangle{\pgfqpoint{0.396477in}{0.232212in}}{\pgfqpoint{5.775408in}{3.672290in}}%
\pgfusepath{clip}%
\pgfsetrectcap%
\pgfsetroundjoin%
\pgfsetlinewidth{1.204500pt}%
\definecolor{currentstroke}{rgb}{0.000000,0.000000,0.000000}%
\pgfsetstrokecolor{currentstroke}%
\pgfsetdash{}{0pt}%
\pgfpathmoveto{\pgfqpoint{0.396477in}{2.855276in}}%
\pgfpathlineto{\pgfqpoint{4.246749in}{0.232212in}}%
\pgfusepath{stroke}%
\end{pgfscope}%
\begin{pgfscope}%
\pgfpathrectangle{\pgfqpoint{0.396477in}{0.232212in}}{\pgfqpoint{5.775408in}{3.672290in}}%
\pgfusepath{clip}%
\pgfsetbuttcap%
\pgfsetroundjoin%
\pgfsetlinewidth{1.204500pt}%
\definecolor{currentstroke}{rgb}{0.000000,0.000000,0.000000}%
\pgfsetstrokecolor{currentstroke}%
\pgfsetdash{{4.440000pt}{1.920000pt}}{0.000000pt}%
\pgfpathmoveto{\pgfqpoint{0.844688in}{3.914501in}}%
\pgfpathlineto{\pgfqpoint{0.848864in}{3.886183in}}%
\pgfpathlineto{\pgfqpoint{0.860361in}{3.811757in}}%
\pgfpathlineto{\pgfqpoint{0.871858in}{3.740614in}}%
\pgfpathlineto{\pgfqpoint{0.883355in}{3.672536in}}%
\pgfpathlineto{\pgfqpoint{0.894852in}{3.607321in}}%
\pgfpathlineto{\pgfqpoint{0.906348in}{3.544789in}}%
\pgfpathlineto{\pgfqpoint{0.917845in}{3.484771in}}%
\pgfpathlineto{\pgfqpoint{0.929342in}{3.427114in}}%
\pgfpathlineto{\pgfqpoint{0.940839in}{3.371678in}}%
\pgfpathlineto{\pgfqpoint{0.952336in}{3.318333in}}%
\pgfpathlineto{\pgfqpoint{0.963832in}{3.266959in}}%
\pgfpathlineto{\pgfqpoint{0.975329in}{3.217445in}}%
\pgfpathlineto{\pgfqpoint{0.986826in}{3.169689in}}%
\pgfpathlineto{\pgfqpoint{0.998323in}{3.123596in}}%
\pgfpathlineto{\pgfqpoint{1.009820in}{3.079079in}}%
\pgfpathlineto{\pgfqpoint{1.021316in}{3.036054in}}%
\pgfpathlineto{\pgfqpoint{1.032813in}{2.994447in}}%
\pgfpathlineto{\pgfqpoint{1.044310in}{2.954186in}}%
\pgfpathlineto{\pgfqpoint{1.055807in}{2.915205in}}%
\pgfpathlineto{\pgfqpoint{1.067304in}{2.877441in}}%
\pgfpathlineto{\pgfqpoint{1.078800in}{2.840837in}}%
\pgfpathlineto{\pgfqpoint{1.090297in}{2.805339in}}%
\pgfpathlineto{\pgfqpoint{1.101794in}{2.770895in}}%
\pgfpathlineto{\pgfqpoint{1.113291in}{2.737459in}}%
\pgfpathlineto{\pgfqpoint{1.124788in}{2.704984in}}%
\pgfpathlineto{\pgfqpoint{1.136284in}{2.673429in}}%
\pgfpathlineto{\pgfqpoint{1.147781in}{2.642754in}}%
\pgfpathlineto{\pgfqpoint{1.159278in}{2.612922in}}%
\pgfpathlineto{\pgfqpoint{1.170775in}{2.583897in}}%
\pgfpathlineto{\pgfqpoint{1.182272in}{2.555646in}}%
\pgfpathlineto{\pgfqpoint{1.193769in}{2.528137in}}%
\pgfpathlineto{\pgfqpoint{1.205265in}{2.501341in}}%
\pgfpathlineto{\pgfqpoint{1.216762in}{2.475229in}}%
\pgfpathlineto{\pgfqpoint{1.228259in}{2.449775in}}%
\pgfpathlineto{\pgfqpoint{1.239756in}{2.424953in}}%
\pgfpathlineto{\pgfqpoint{1.251253in}{2.400740in}}%
\pgfpathlineto{\pgfqpoint{1.262749in}{2.377112in}}%
\pgfpathlineto{\pgfqpoint{1.274246in}{2.354048in}}%
\pgfpathlineto{\pgfqpoint{1.285743in}{2.331528in}}%
\pgfpathlineto{\pgfqpoint{1.297240in}{2.309531in}}%
\pgfpathlineto{\pgfqpoint{1.308737in}{2.288039in}}%
\pgfpathlineto{\pgfqpoint{1.320233in}{2.267035in}}%
\pgfpathlineto{\pgfqpoint{1.331730in}{2.246501in}}%
\pgfpathlineto{\pgfqpoint{1.343227in}{2.226422in}}%
\pgfpathlineto{\pgfqpoint{1.354724in}{2.206781in}}%
\pgfusepath{stroke}%
\end{pgfscope}%
\begin{pgfscope}%
\pgfpathrectangle{\pgfqpoint{0.396477in}{0.232212in}}{\pgfqpoint{5.775408in}{3.672290in}}%
\pgfusepath{clip}%
\pgfsetrectcap%
\pgfsetroundjoin%
\pgfsetlinewidth{1.204500pt}%
\definecolor{currentstroke}{rgb}{0.000000,0.000000,0.000000}%
\pgfsetstrokecolor{currentstroke}%
\pgfsetdash{}{0pt}%
\pgfpathmoveto{\pgfqpoint{1.366221in}{2.187565in}}%
\pgfpathlineto{\pgfqpoint{1.400711in}{2.132324in}}%
\pgfpathlineto{\pgfqpoint{1.435201in}{2.080431in}}%
\pgfpathlineto{\pgfqpoint{1.481189in}{2.015929in}}%
\pgfpathlineto{\pgfqpoint{1.527176in}{1.956220in}}%
\pgfpathlineto{\pgfqpoint{1.573163in}{1.900773in}}%
\pgfpathlineto{\pgfqpoint{1.619150in}{1.849134in}}%
\pgfpathlineto{\pgfqpoint{1.665137in}{1.800914in}}%
\pgfpathlineto{\pgfqpoint{1.711125in}{1.755774in}}%
\pgfpathlineto{\pgfqpoint{1.757112in}{1.713419in}}%
\pgfpathlineto{\pgfqpoint{1.803099in}{1.673591in}}%
\pgfpathlineto{\pgfqpoint{1.860583in}{1.627019in}}%
\pgfpathlineto{\pgfqpoint{1.918067in}{1.583661in}}%
\pgfpathlineto{\pgfqpoint{1.975551in}{1.543185in}}%
\pgfpathlineto{\pgfqpoint{2.033035in}{1.505304in}}%
\pgfpathlineto{\pgfqpoint{2.090519in}{1.469771in}}%
\pgfpathlineto{\pgfqpoint{2.159500in}{1.429923in}}%
\pgfpathlineto{\pgfqpoint{2.228481in}{1.392816in}}%
\pgfpathlineto{\pgfqpoint{2.297462in}{1.358168in}}%
\pgfpathlineto{\pgfqpoint{2.366443in}{1.325734in}}%
\pgfpathlineto{\pgfqpoint{2.446920in}{1.290414in}}%
\pgfpathlineto{\pgfqpoint{2.527398in}{1.257530in}}%
\pgfpathlineto{\pgfqpoint{2.619372in}{1.222611in}}%
\pgfpathlineto{\pgfqpoint{2.711347in}{1.190226in}}%
\pgfpathlineto{\pgfqpoint{2.803321in}{1.160102in}}%
\pgfpathlineto{\pgfqpoint{2.906792in}{1.128622in}}%
\pgfpathlineto{\pgfqpoint{3.010264in}{1.099418in}}%
\pgfpathlineto{\pgfqpoint{3.125232in}{1.069342in}}%
\pgfpathlineto{\pgfqpoint{3.240200in}{1.041487in}}%
\pgfpathlineto{\pgfqpoint{3.366665in}{1.013122in}}%
\pgfpathlineto{\pgfqpoint{3.504626in}{0.984583in}}%
\pgfpathlineto{\pgfqpoint{3.642588in}{0.958261in}}%
\pgfpathlineto{\pgfqpoint{3.792046in}{0.931950in}}%
\pgfpathlineto{\pgfqpoint{3.953002in}{0.905875in}}%
\pgfpathlineto{\pgfqpoint{4.125454in}{0.880220in}}%
\pgfpathlineto{\pgfqpoint{4.309403in}{0.855132in}}%
\pgfpathlineto{\pgfqpoint{4.504848in}{0.830726in}}%
\pgfpathlineto{\pgfqpoint{4.711791in}{0.807087in}}%
\pgfpathlineto{\pgfqpoint{4.941727in}{0.783129in}}%
\pgfpathlineto{\pgfqpoint{5.183160in}{0.760238in}}%
\pgfpathlineto{\pgfqpoint{5.447586in}{0.737462in}}%
\pgfpathlineto{\pgfqpoint{5.735006in}{0.715046in}}%
\pgfpathlineto{\pgfqpoint{6.045420in}{0.693180in}}%
\pgfpathlineto{\pgfqpoint{6.171885in}{0.684886in}}%
\pgfpathlineto{\pgfqpoint{6.171885in}{0.684886in}}%
\pgfusepath{stroke}%
\end{pgfscope}%
\begin{pgfscope}%
\pgfpathrectangle{\pgfqpoint{0.396477in}{0.232212in}}{\pgfqpoint{5.775408in}{3.672290in}}%
\pgfusepath{clip}%
\pgfsetbuttcap%
\pgfsetroundjoin%
\pgfsetlinewidth{1.003750pt}%
\definecolor{currentstroke}{rgb}{0.000000,0.000000,0.000000}%
\pgfsetstrokecolor{currentstroke}%
\pgfsetdash{{1.000000pt}{1.650000pt}}{0.000000pt}%
\pgfpathmoveto{\pgfqpoint{1.359045in}{0.232212in}}%
\pgfpathlineto{\pgfqpoint{1.359045in}{2.199510in}}%
\pgfusepath{stroke}%
\end{pgfscope}%
\begin{pgfscope}%
\pgfpathrectangle{\pgfqpoint{0.396477in}{0.232212in}}{\pgfqpoint{5.775408in}{3.672290in}}%
\pgfusepath{clip}%
\pgfsetbuttcap%
\pgfsetroundjoin%
\pgfsetlinewidth{1.204500pt}%
\definecolor{currentstroke}{rgb}{0.000000,0.000000,0.000000}%
\pgfsetstrokecolor{currentstroke}%
\pgfsetdash{{4.440000pt}{1.920000pt}}{0.000000pt}%
\pgfpathmoveto{\pgfqpoint{0.949143in}{3.914501in}}%
\pgfpathlineto{\pgfqpoint{0.963832in}{3.835905in}}%
\pgfpathlineto{\pgfqpoint{0.986826in}{3.720399in}}%
\pgfpathlineto{\pgfqpoint{1.009820in}{3.612802in}}%
\pgfpathlineto{\pgfqpoint{1.044310in}{3.464495in}}%
\pgfpathlineto{\pgfqpoint{1.078800in}{3.329895in}}%
\pgfpathlineto{\pgfqpoint{1.113291in}{3.207136in}}%
\pgfpathlineto{\pgfqpoint{1.147781in}{3.094676in}}%
\pgfpathlineto{\pgfqpoint{1.182272in}{2.991237in}}%
\pgfpathlineto{\pgfqpoint{1.216762in}{2.895744in}}%
\pgfpathlineto{\pgfqpoint{1.251253in}{2.807290in}}%
\pgfpathlineto{\pgfqpoint{1.285743in}{2.725102in}}%
\pgfpathlineto{\pgfqpoint{1.320233in}{2.648518in}}%
\pgfpathlineto{\pgfqpoint{1.354724in}{2.576968in}}%
\pgfpathlineto{\pgfqpoint{1.389214in}{2.509956in}}%
\pgfpathlineto{\pgfqpoint{1.423705in}{2.447052in}}%
\pgfpathlineto{\pgfqpoint{1.458195in}{2.387879in}}%
\pgfpathlineto{\pgfqpoint{1.492685in}{2.332103in}}%
\pgfpathlineto{\pgfqpoint{1.538673in}{2.262520in}}%
\pgfpathlineto{\pgfqpoint{1.584660in}{2.197856in}}%
\pgfpathlineto{\pgfqpoint{1.630647in}{2.137592in}}%
\pgfpathlineto{\pgfqpoint{1.676634in}{2.081281in}}%
\pgfpathlineto{\pgfqpoint{1.722622in}{2.028536in}}%
\pgfpathlineto{\pgfqpoint{1.768609in}{1.979018in}}%
\pgfpathlineto{\pgfqpoint{1.814596in}{1.932431in}}%
\pgfpathlineto{\pgfqpoint{1.860583in}{1.888514in}}%
\pgfpathlineto{\pgfqpoint{1.918067in}{1.837027in}}%
\pgfpathlineto{\pgfqpoint{1.975551in}{1.788963in}}%
\pgfpathlineto{\pgfqpoint{2.033035in}{1.743981in}}%
\pgfpathlineto{\pgfqpoint{2.090519in}{1.701785in}}%
\pgfpathlineto{\pgfqpoint{2.125010in}{1.677697in}}%
\pgfpathlineto{\pgfqpoint{2.125010in}{1.677697in}}%
\pgfusepath{stroke}%
\end{pgfscope}%
\begin{pgfscope}%
\pgfpathrectangle{\pgfqpoint{0.396477in}{0.232212in}}{\pgfqpoint{5.775408in}{3.672290in}}%
\pgfusepath{clip}%
\pgfsetrectcap%
\pgfsetroundjoin%
\pgfsetlinewidth{1.204500pt}%
\definecolor{currentstroke}{rgb}{0.000000,0.000000,0.000000}%
\pgfsetstrokecolor{currentstroke}%
\pgfsetdash{}{0pt}%
\pgfpathmoveto{\pgfqpoint{2.136506in}{1.669860in}}%
\pgfpathlineto{\pgfqpoint{2.193991in}{1.632051in}}%
\pgfpathlineto{\pgfqpoint{2.262971in}{1.589485in}}%
\pgfpathlineto{\pgfqpoint{2.331952in}{1.549690in}}%
\pgfpathlineto{\pgfqpoint{2.400933in}{1.512395in}}%
\pgfpathlineto{\pgfqpoint{2.469914in}{1.477365in}}%
\pgfpathlineto{\pgfqpoint{2.550391in}{1.439086in}}%
\pgfpathlineto{\pgfqpoint{2.630869in}{1.403324in}}%
\pgfpathlineto{\pgfqpoint{2.711347in}{1.369832in}}%
\pgfpathlineto{\pgfqpoint{2.803321in}{1.334060in}}%
\pgfpathlineto{\pgfqpoint{2.895296in}{1.300693in}}%
\pgfpathlineto{\pgfqpoint{2.998767in}{1.265729in}}%
\pgfpathlineto{\pgfqpoint{3.102238in}{1.233207in}}%
\pgfpathlineto{\pgfqpoint{3.205709in}{1.202872in}}%
\pgfpathlineto{\pgfqpoint{3.320677in}{1.171466in}}%
\pgfpathlineto{\pgfqpoint{3.447142in}{1.139417in}}%
\pgfpathlineto{\pgfqpoint{3.573607in}{1.109698in}}%
\pgfpathlineto{\pgfqpoint{3.711569in}{1.079640in}}%
\pgfpathlineto{\pgfqpoint{3.861027in}{1.049544in}}%
\pgfpathlineto{\pgfqpoint{4.010486in}{1.021722in}}%
\pgfpathlineto{\pgfqpoint{4.171441in}{0.994011in}}%
\pgfpathlineto{\pgfqpoint{4.343893in}{0.966612in}}%
\pgfpathlineto{\pgfqpoint{4.527842in}{0.939689in}}%
\pgfpathlineto{\pgfqpoint{4.723288in}{0.913375in}}%
\pgfpathlineto{\pgfqpoint{4.930230in}{0.887774in}}%
\pgfpathlineto{\pgfqpoint{5.148669in}{0.862960in}}%
\pgfpathlineto{\pgfqpoint{5.390102in}{0.837843in}}%
\pgfpathlineto{\pgfqpoint{5.643032in}{0.813796in}}%
\pgfpathlineto{\pgfqpoint{5.918955in}{0.789859in}}%
\pgfpathlineto{\pgfqpoint{6.171885in}{0.769752in}}%
\pgfpathlineto{\pgfqpoint{6.171885in}{0.769752in}}%
\pgfusepath{stroke}%
\end{pgfscope}%
\begin{pgfscope}%
\pgfpathrectangle{\pgfqpoint{0.396477in}{0.232212in}}{\pgfqpoint{5.775408in}{3.672290in}}%
\pgfusepath{clip}%
\pgfsetbuttcap%
\pgfsetroundjoin%
\pgfsetlinewidth{1.003750pt}%
\definecolor{currentstroke}{rgb}{0.000000,0.000000,0.000000}%
\pgfsetstrokecolor{currentstroke}%
\pgfsetdash{{1.000000pt}{1.650000pt}}{0.000000pt}%
\pgfpathmoveto{\pgfqpoint{2.131215in}{0.232212in}}%
\pgfpathlineto{\pgfqpoint{2.131215in}{1.673456in}}%
\pgfusepath{stroke}%
\end{pgfscope}%
\begin{pgfscope}%
\pgfpathrectangle{\pgfqpoint{0.396477in}{0.232212in}}{\pgfqpoint{5.775408in}{3.672290in}}%
\pgfusepath{clip}%
\pgfsetbuttcap%
\pgfsetroundjoin%
\pgfsetlinewidth{1.204500pt}%
\definecolor{currentstroke}{rgb}{0.000000,0.000000,0.000000}%
\pgfsetstrokecolor{currentstroke}%
\pgfsetdash{{4.440000pt}{1.920000pt}}{0.000000pt}%
\pgfpathmoveto{\pgfqpoint{0.807019in}{3.914501in}}%
\pgfpathlineto{\pgfqpoint{0.825871in}{3.780930in}}%
\pgfpathlineto{\pgfqpoint{0.848864in}{3.632335in}}%
\pgfpathlineto{\pgfqpoint{0.871858in}{3.496880in}}%
\pgfpathlineto{\pgfqpoint{0.894852in}{3.372847in}}%
\pgfpathlineto{\pgfqpoint{0.917845in}{3.258810in}}%
\pgfpathlineto{\pgfqpoint{0.940839in}{3.153574in}}%
\pgfpathlineto{\pgfqpoint{0.963832in}{3.056129in}}%
\pgfpathlineto{\pgfqpoint{0.986826in}{2.965617in}}%
\pgfpathlineto{\pgfqpoint{1.009820in}{2.881302in}}%
\pgfpathlineto{\pgfqpoint{1.044310in}{2.765086in}}%
\pgfpathlineto{\pgfqpoint{1.078800in}{2.659611in}}%
\pgfpathlineto{\pgfqpoint{1.113291in}{2.563415in}}%
\pgfpathlineto{\pgfqpoint{1.147781in}{2.475290in}}%
\pgfpathlineto{\pgfqpoint{1.182272in}{2.394233in}}%
\pgfpathlineto{\pgfqpoint{1.216762in}{2.319403in}}%
\pgfpathlineto{\pgfqpoint{1.251253in}{2.250089in}}%
\pgfpathlineto{\pgfqpoint{1.285743in}{2.185685in}}%
\pgfpathlineto{\pgfqpoint{1.320233in}{2.125672in}}%
\pgfpathlineto{\pgfqpoint{1.354724in}{2.069604in}}%
\pgfpathlineto{\pgfqpoint{1.389214in}{2.017093in}}%
\pgfpathlineto{\pgfqpoint{1.423705in}{1.967800in}}%
\pgfpathlineto{\pgfqpoint{1.469692in}{1.906577in}}%
\pgfpathlineto{\pgfqpoint{1.515679in}{1.849951in}}%
\pgfpathlineto{\pgfqpoint{1.561666in}{1.797406in}}%
\pgfpathlineto{\pgfqpoint{1.607653in}{1.748505in}}%
\pgfpathlineto{\pgfqpoint{1.653641in}{1.702871in}}%
\pgfpathlineto{\pgfqpoint{1.699628in}{1.660177in}}%
\pgfpathlineto{\pgfqpoint{1.745615in}{1.620140in}}%
\pgfpathlineto{\pgfqpoint{1.803099in}{1.573456in}}%
\pgfpathlineto{\pgfqpoint{1.860583in}{1.530120in}}%
\pgfpathlineto{\pgfqpoint{1.918067in}{1.489773in}}%
\pgfpathlineto{\pgfqpoint{1.975551in}{1.452109in}}%
\pgfpathlineto{\pgfqpoint{2.033035in}{1.416860in}}%
\pgfpathlineto{\pgfqpoint{2.102016in}{1.377426in}}%
\pgfpathlineto{\pgfqpoint{2.170997in}{1.340792in}}%
\pgfpathlineto{\pgfqpoint{2.239978in}{1.306661in}}%
\pgfpathlineto{\pgfqpoint{2.308959in}{1.274778in}}%
\pgfpathlineto{\pgfqpoint{2.389436in}{1.240129in}}%
\pgfpathlineto{\pgfqpoint{2.469914in}{1.207936in}}%
\pgfpathlineto{\pgfqpoint{2.550391in}{1.177939in}}%
\pgfpathlineto{\pgfqpoint{2.642366in}{1.146062in}}%
\pgfpathlineto{\pgfqpoint{2.734340in}{1.116475in}}%
\pgfpathlineto{\pgfqpoint{2.837812in}{1.085624in}}%
\pgfpathlineto{\pgfqpoint{2.941283in}{1.057065in}}%
\pgfpathlineto{\pgfqpoint{3.056251in}{1.027713in}}%
\pgfpathlineto{\pgfqpoint{3.090741in}{1.019353in}}%
\pgfpathlineto{\pgfqpoint{3.090741in}{1.019353in}}%
\pgfusepath{stroke}%
\end{pgfscope}%
\begin{pgfscope}%
\pgfpathrectangle{\pgfqpoint{0.396477in}{0.232212in}}{\pgfqpoint{5.775408in}{3.672290in}}%
\pgfusepath{clip}%
\pgfsetrectcap%
\pgfsetroundjoin%
\pgfsetlinewidth{1.204500pt}%
\definecolor{currentstroke}{rgb}{0.000000,0.000000,0.000000}%
\pgfsetstrokecolor{currentstroke}%
\pgfsetdash{}{0pt}%
\pgfpathmoveto{\pgfqpoint{3.102238in}{1.016609in}}%
\pgfpathlineto{\pgfqpoint{3.228703in}{0.987771in}}%
\pgfpathlineto{\pgfqpoint{3.355168in}{0.961185in}}%
\pgfpathlineto{\pgfqpoint{3.493129in}{0.934445in}}%
\pgfpathlineto{\pgfqpoint{3.642588in}{0.907821in}}%
\pgfpathlineto{\pgfqpoint{3.803543in}{0.881536in}}%
\pgfpathlineto{\pgfqpoint{3.975995in}{0.855770in}}%
\pgfpathlineto{\pgfqpoint{4.159944in}{0.830666in}}%
\pgfpathlineto{\pgfqpoint{4.355390in}{0.806330in}}%
\pgfpathlineto{\pgfqpoint{4.562332in}{0.782837in}}%
\pgfpathlineto{\pgfqpoint{4.780772in}{0.760238in}}%
\pgfpathlineto{\pgfqpoint{5.022204in}{0.737531in}}%
\pgfpathlineto{\pgfqpoint{5.275134in}{0.715946in}}%
\pgfpathlineto{\pgfqpoint{5.551057in}{0.694609in}}%
\pgfpathlineto{\pgfqpoint{5.849974in}{0.673721in}}%
\pgfpathlineto{\pgfqpoint{6.171885in}{0.653438in}}%
\pgfpathlineto{\pgfqpoint{6.171885in}{0.653438in}}%
\pgfusepath{stroke}%
\end{pgfscope}%
\begin{pgfscope}%
\pgfpathrectangle{\pgfqpoint{0.396477in}{0.232212in}}{\pgfqpoint{5.775408in}{3.672290in}}%
\pgfusepath{clip}%
\pgfsetbuttcap%
\pgfsetroundjoin%
\pgfsetlinewidth{1.003750pt}%
\definecolor{currentstroke}{rgb}{0.000000,0.000000,0.000000}%
\pgfsetstrokecolor{currentstroke}%
\pgfsetdash{{1.000000pt}{1.650000pt}}{0.000000pt}%
\pgfpathmoveto{\pgfqpoint{3.091667in}{0.232212in}}%
\pgfpathlineto{\pgfqpoint{3.091667in}{1.019131in}}%
\pgfusepath{stroke}%
\end{pgfscope}%
\begin{pgfscope}%
\pgfsetbuttcap%
\pgfsetmiterjoin%
\definecolor{currentfill}{rgb}{0.000000,0.000000,0.000000}%
\pgfsetfillcolor{currentfill}%
\pgfsetlinewidth{1.003750pt}%
\definecolor{currentstroke}{rgb}{0.000000,0.000000,0.000000}%
\pgfsetstrokecolor{currentstroke}%
\pgfsetdash{}{0pt}%
\pgfsys@defobject{currentmarker}{\pgfqpoint{-0.027778in}{-0.027778in}}{\pgfqpoint{0.027778in}{0.027778in}}{%
\pgfpathmoveto{\pgfqpoint{0.027778in}{-0.000000in}}%
\pgfpathlineto{\pgfqpoint{-0.027778in}{0.027778in}}%
\pgfpathlineto{\pgfqpoint{-0.027778in}{-0.027778in}}%
\pgfpathlineto{\pgfqpoint{0.027778in}{-0.000000in}}%
\pgfpathclose%
\pgfusepath{stroke,fill}%
}%
\begin{pgfscope}%
\pgfsys@transformshift{6.171885in}{0.232212in}%
\pgfsys@useobject{currentmarker}{}%
\end{pgfscope}%
\end{pgfscope}%
\begin{pgfscope}%
\pgfsetbuttcap%
\pgfsetmiterjoin%
\definecolor{currentfill}{rgb}{0.000000,0.000000,0.000000}%
\pgfsetfillcolor{currentfill}%
\pgfsetlinewidth{1.003750pt}%
\definecolor{currentstroke}{rgb}{0.000000,0.000000,0.000000}%
\pgfsetstrokecolor{currentstroke}%
\pgfsetdash{}{0pt}%
\pgfsys@defobject{currentmarker}{\pgfqpoint{-0.027778in}{-0.027778in}}{\pgfqpoint{0.027778in}{0.027778in}}{%
\pgfpathmoveto{\pgfqpoint{0.000000in}{0.027778in}}%
\pgfpathlineto{\pgfqpoint{-0.027778in}{-0.027778in}}%
\pgfpathlineto{\pgfqpoint{0.027778in}{-0.027778in}}%
\pgfpathlineto{\pgfqpoint{0.000000in}{0.027778in}}%
\pgfpathclose%
\pgfusepath{stroke,fill}%
}%
\begin{pgfscope}%
\pgfsys@transformshift{0.396477in}{3.904501in}%
\pgfsys@useobject{currentmarker}{}%
\end{pgfscope}%
\end{pgfscope}%
\begin{pgfscope}%
\pgfsetrectcap%
\pgfsetmiterjoin%
\pgfsetlinewidth{1.505625pt}%
\definecolor{currentstroke}{rgb}{0.000000,0.000000,0.000000}%
\pgfsetstrokecolor{currentstroke}%
\pgfsetdash{}{0pt}%
\pgfpathmoveto{\pgfqpoint{0.396477in}{0.232212in}}%
\pgfpathlineto{\pgfqpoint{0.396477in}{3.904501in}}%
\pgfusepath{stroke}%
\end{pgfscope}%
\begin{pgfscope}%
\pgfsetrectcap%
\pgfsetmiterjoin%
\pgfsetlinewidth{1.505625pt}%
\definecolor{currentstroke}{rgb}{0.000000,0.000000,0.000000}%
\pgfsetstrokecolor{currentstroke}%
\pgfsetdash{}{0pt}%
\pgfpathmoveto{\pgfqpoint{0.396477in}{0.232212in}}%
\pgfpathlineto{\pgfqpoint{6.171885in}{0.232212in}}%
\pgfusepath{stroke}%
\end{pgfscope}%
\begin{pgfscope}%
\definecolor{textcolor}{rgb}{0.000000,0.000000,0.000000}%
\pgfsetstrokecolor{textcolor}%
\pgfsetfillcolor{textcolor}%
\pgftext[x=1.551558in,y=1.150284in,left,base]{\color{textcolor}{\rmfamily\fontsize{9.000000}{10.800000}\selectfont\catcode`\^=\active\def^{\ifmmode\sp\else\^{}\fi}\catcode`\%=\active\def%{\%}$v_{L_1}(x),$}}%
\end{pgfscope}%
\begin{pgfscope}%
\definecolor{textcolor}{rgb}{0.000000,0.000000,0.000000}%
\pgfsetstrokecolor{textcolor}%
\pgfsetfillcolor{textcolor}%
\pgftext[x=1.551558in,y=0.966670in,left,base]{\color{textcolor}{\rmfamily\fontsize{9.000000}{10.800000}\selectfont\catcode`\^=\active\def^{\ifmmode\sp\else\^{}\fi}\catcode`\%=\active\def%{\%}$L_1 \le x \le L_2$}}%
\end{pgfscope}%
\begin{pgfscope}%
\definecolor{textcolor}{rgb}{0.000000,0.000000,0.000000}%
\pgfsetstrokecolor{textcolor}%
\pgfsetfillcolor{textcolor}%
\pgftext[x=3.476694in,y=1.753589in,left,base]{\color{textcolor}{\rmfamily\fontsize{9.000000}{10.800000}\selectfont\catcode`\^=\active\def^{\ifmmode\sp\else\^{}\fi}\catcode`\%=\active\def%{\%}$v_{L_*}(x), x \ge 0$}}%
\end{pgfscope}%
\begin{pgfscope}%
\definecolor{textcolor}{rgb}{0.000000,0.000000,0.000000}%
\pgfsetstrokecolor{textcolor}%
\pgfsetfillcolor{textcolor}%
\pgftext[x=4.439262in,y=0.966670in,left,base]{\color{textcolor}{\rmfamily\fontsize{9.000000}{10.800000}\selectfont\catcode`\^=\active\def^{\ifmmode\sp\else\^{}\fi}\catcode`\%=\active\def%{\%}$v_{L_2}(x) = v_{L_2}(x), x \ge L_2$}}%
\end{pgfscope}%
\end{pgfpicture}%
\makeatother%
\endgroup%

	\caption{\((K-L)(\frac{x}{L})^{-\frac{2r}{\sigma^2}}\) for three values of \(L\).}\label{fig:ExerciseBoundaryValue}
\end{figure}

As \Cref{fig:ExerciseBoundaryValue} suggests, for any value of \(L\) smaller than \(L_*\), the function \(v_L(x)\) agrees with the intrinsic value for \(0 \le x \le L\), lies below the intrinsic value immediately to the right of \(L\), and lies below \(v_{L_{*}}(x)\) everywhere to the right of \(L\). For any value of \(L\) larger than \(L_*\), the function \(v_L(x)\) agrees with the intrinsic value for \(0 \le x \le L\) and lies below \(v_{L_{*}}(x)\) for all \(x \ge L_*\). Thus, among those exercise policies of the form
\begin{center}
	``Exercise the put as soon as \(S(t)\) falls to the level \(L\)''.
\end{center}
the best one is obtained by choosing \(L = L_*\). We expect therefore that \(v_{L_{*}}(x)\) is the price of the put \(v_*(x)\) of \Cref{defn:PerpetualAmericanPut}. We prove this below.

We must first determine the value of \(L_*\). We note that
\begin{equation}
	\label{eq:VLAboveL}
	v_L(x) = (K - L)L^{\frac{2r}{\sigma^2}} x^{-\frac{2r}{\sigma^2}} \quad \text{for all } x \ge L.
\end{equation}
From \Cref{fig:ExerciseBoundaryValue}, we know that \(L_*\) is the value of \(L\) that maximizes this quantity when we hold \(x\) fixed. We thus define
\begin{equation}
	\label{eq:GLDefinition}
	g(L) = (K - L)L^{\frac{2r}{\sigma^2}}
\end{equation}
and seek the value of \(L\) that maximizes this function over \(L \ge 0\). Because \(\frac{2r}{\sigma^2}\) is strictly positive, we have \(g(0) = 0\) and \(\lim_{L \to \infty} g(L) = -\infty\). Moreover,
\begin{equation}
	\label{eq:GPrimeL}
	g'(L) = -L^{\frac{2r}{\sigma^2}} + \frac{2r}{\sigma^2}(K - L)L^{\frac{2r}{\sigma^2}-1} = -\frac{2r + \sigma^2}{\sigma^2}L^{\frac{2r}{\sigma^2}} + \frac{2r}{\sigma^2} K L^{\frac{2r}{\sigma^2}-1}.
\end{equation}
Setting this equal to zero, we solve for
\begin{equation}
	\label{eq:LStar}
	L_* = \frac{2r}{2r + \sigma^2} K.
\end{equation}
This is a number between 0 and \(K\). Furthermore,
\begin{equation}
	\label{eq:GLStar}
	g(L_*) = \frac{\sigma^2}{2r + \sigma^2} \left( \frac{2r}{2r + \sigma^2} \right)^{\frac{2r}{\sigma^2}} K^{\frac{2r + \sigma^2}{\sigma^2}}
\end{equation}
is strictly positive. Therefore, the graph of \(y = g(L)\) must be as shown in \Cref{fig:GraphOfGL} and \(L_*\) given by \Cref{eq:LStar} is the point where \(g(L)\) attains its maximum.

\begin{figure}[H]
	\centering
	%% Creator: Matplotlib, PGF backend
%%
%% To include the figure in your LaTeX document, write
%%   \input{<filename>.pgf}
%%
%% Make sure the required packages are loaded in your preamble
%%   \usepackage{pgf}
%%
%% Also ensure that all the required font packages are loaded; for instance,
%% the lmodern package is sometimes necessary when using math font.
%%   \usepackage{lmodern}
%%
%% Figures using additional raster images can only be included by \input if
%% they are in the same directory as the main LaTeX file. For loading figures
%% from other directories you can use the `import` package
%%   \usepackage{import}
%%
%% and then include the figures with
%%   \import{<path to file>}{<filename>.pgf}
%%
%% Matplotlib used the following preamble
%%   \def\mathdefault#1{#1}
%%   \everymath=\expandafter{\the\everymath\displaystyle}
%%   \IfFileExists{scrextend.sty}{
%%     \usepackage[fontsize=10.000000pt]{scrextend}
%%   }{
%%     \renewcommand{\normalsize}{\fontsize{10.000000}{12.000000}\selectfont}
%%     \normalsize
%%   }
%%   \usepackage{amsmath}
%%   \usepackage{amssymb}
%%   \usepackage{amsfonts}
%%   \usepackage{libertine}
%%   \usepackage[libertine]{newtxmath}
%%   \makeatletter\@ifpackageloaded{underscore}{}{\usepackage[strings]{underscore}}\makeatother
%%
\begingroup%
\makeatletter%
\begin{pgfpicture}%
\pgfpathrectangle{\pgfpointorigin}{\pgfqpoint{5.358433in}{3.561048in}}%
\pgfusepath{use as bounding box, clip}%
\begin{pgfscope}%
\pgfsetbuttcap%
\pgfsetmiterjoin%
\definecolor{currentfill}{rgb}{1.000000,1.000000,1.000000}%
\pgfsetfillcolor{currentfill}%
\pgfsetlinewidth{0.000000pt}%
\definecolor{currentstroke}{rgb}{1.000000,1.000000,1.000000}%
\pgfsetstrokecolor{currentstroke}%
\pgfsetdash{}{0pt}%
\pgfpathmoveto{\pgfqpoint{0.000000in}{0.000000in}}%
\pgfpathlineto{\pgfqpoint{5.358433in}{0.000000in}}%
\pgfpathlineto{\pgfqpoint{5.358433in}{3.561048in}}%
\pgfpathlineto{\pgfqpoint{0.000000in}{3.561048in}}%
\pgfpathlineto{\pgfqpoint{0.000000in}{0.000000in}}%
\pgfpathclose%
\pgfusepath{fill}%
\end{pgfscope}%
\begin{pgfscope}%
\pgfsetbuttcap%
\pgfsetmiterjoin%
\definecolor{currentfill}{rgb}{1.000000,1.000000,1.000000}%
\pgfsetfillcolor{currentfill}%
\pgfsetlinewidth{0.000000pt}%
\definecolor{currentstroke}{rgb}{0.000000,0.000000,0.000000}%
\pgfsetstrokecolor{currentstroke}%
\pgfsetstrokeopacity{0.000000}%
\pgfsetdash{}{0pt}%
\pgfpathmoveto{\pgfqpoint{0.088135in}{0.050000in}}%
\pgfpathlineto{\pgfqpoint{5.125635in}{0.050000in}}%
\pgfpathlineto{\pgfqpoint{5.125635in}{3.303250in}}%
\pgfpathlineto{\pgfqpoint{0.088135in}{3.303250in}}%
\pgfpathlineto{\pgfqpoint{0.088135in}{0.050000in}}%
\pgfpathclose%
\pgfusepath{fill}%
\end{pgfscope}%
\begin{pgfscope}%
\pgfsetbuttcap%
\pgfsetroundjoin%
\definecolor{currentfill}{rgb}{0.000000,0.000000,0.000000}%
\pgfsetfillcolor{currentfill}%
\pgfsetlinewidth{1.204500pt}%
\definecolor{currentstroke}{rgb}{0.000000,0.000000,0.000000}%
\pgfsetstrokecolor{currentstroke}%
\pgfsetdash{}{0pt}%
\pgfsys@defobject{currentmarker}{\pgfqpoint{0.000000in}{-0.027778in}}{\pgfqpoint{0.000000in}{0.027778in}}{%
\pgfpathmoveto{\pgfqpoint{0.000000in}{-0.027778in}}%
\pgfpathlineto{\pgfqpoint{0.000000in}{0.027778in}}%
\pgfusepath{stroke,fill}%
}%
\begin{pgfscope}%
\pgfsys@transformshift{2.530201in}{2.988005in}%
\pgfsys@useobject{currentmarker}{}%
\end{pgfscope}%
\end{pgfscope}%
\begin{pgfscope}%
\pgfsetbuttcap%
\pgfsetroundjoin%
\definecolor{currentfill}{rgb}{0.000000,0.000000,0.000000}%
\pgfsetfillcolor{currentfill}%
\pgfsetlinewidth{1.204500pt}%
\definecolor{currentstroke}{rgb}{0.000000,0.000000,0.000000}%
\pgfsetstrokecolor{currentstroke}%
\pgfsetdash{}{0pt}%
\pgfsys@defobject{currentmarker}{\pgfqpoint{0.000000in}{-0.027778in}}{\pgfqpoint{0.000000in}{0.027778in}}{%
\pgfpathmoveto{\pgfqpoint{0.000000in}{-0.027778in}}%
\pgfpathlineto{\pgfqpoint{0.000000in}{0.027778in}}%
\pgfusepath{stroke,fill}%
}%
\begin{pgfscope}%
\pgfsys@transformshift{2.530201in}{3.303250in}%
\pgfsys@useobject{currentmarker}{}%
\end{pgfscope}%
\end{pgfscope}%
\begin{pgfscope}%
\definecolor{textcolor}{rgb}{0.000000,0.000000,0.000000}%
\pgfsetstrokecolor{textcolor}%
\pgfsetfillcolor{textcolor}%
\pgftext[x=2.530201in,y=2.890783in,,top]{\color{textcolor}{\rmfamily\fontsize{10.000000}{12.000000}\selectfont\catcode`\^=\active\def^{\ifmmode\sp\else\^{}\fi}\catcode`\%=\active\def%{\%}$L_*$}}%
\end{pgfscope}%
\begin{pgfscope}%
\pgfpathrectangle{\pgfqpoint{0.088135in}{0.050000in}}{\pgfqpoint{5.037500in}{3.253250in}}%
\pgfusepath{clip}%
\pgfsetrectcap%
\pgfsetroundjoin%
\pgfsetlinewidth{1.204500pt}%
\definecolor{currentstroke}{rgb}{0.000000,0.000000,0.000000}%
\pgfsetstrokecolor{currentstroke}%
\pgfsetdash{}{0pt}%
\pgfpathmoveto{\pgfqpoint{0.170717in}{2.988005in}}%
\pgfpathlineto{\pgfqpoint{0.334393in}{2.989113in}}%
\pgfpathlineto{\pgfqpoint{0.453430in}{2.992184in}}%
\pgfpathlineto{\pgfqpoint{0.567507in}{2.997388in}}%
\pgfpathlineto{\pgfqpoint{0.676624in}{3.004581in}}%
\pgfpathlineto{\pgfqpoint{0.785742in}{3.013962in}}%
\pgfpathlineto{\pgfqpoint{0.899819in}{3.026038in}}%
\pgfpathlineto{\pgfqpoint{1.018856in}{3.040946in}}%
\pgfpathlineto{\pgfqpoint{1.142853in}{3.058708in}}%
\pgfpathlineto{\pgfqpoint{1.281729in}{3.080846in}}%
\pgfpathlineto{\pgfqpoint{1.455325in}{3.110902in}}%
\pgfpathlineto{\pgfqpoint{1.921554in}{3.192969in}}%
\pgfpathlineto{\pgfqpoint{2.045551in}{3.211783in}}%
\pgfpathlineto{\pgfqpoint{2.154668in}{3.226087in}}%
\pgfpathlineto{\pgfqpoint{2.248906in}{3.236277in}}%
\pgfpathlineto{\pgfqpoint{2.338183in}{3.243707in}}%
\pgfpathlineto{\pgfqpoint{2.422501in}{3.248423in}}%
\pgfpathlineto{\pgfqpoint{2.501859in}{3.250545in}}%
\pgfpathlineto{\pgfqpoint{2.576257in}{3.250263in}}%
\pgfpathlineto{\pgfqpoint{2.645696in}{3.247819in}}%
\pgfpathlineto{\pgfqpoint{2.715134in}{3.243087in}}%
\pgfpathlineto{\pgfqpoint{2.779612in}{3.236483in}}%
\pgfpathlineto{\pgfqpoint{2.844091in}{3.227596in}}%
\pgfpathlineto{\pgfqpoint{2.903609in}{3.217235in}}%
\pgfpathlineto{\pgfqpoint{2.963128in}{3.204678in}}%
\pgfpathlineto{\pgfqpoint{3.022646in}{3.189798in}}%
\pgfpathlineto{\pgfqpoint{3.082165in}{3.172472in}}%
\pgfpathlineto{\pgfqpoint{3.136724in}{3.154330in}}%
\pgfpathlineto{\pgfqpoint{3.191282in}{3.133925in}}%
\pgfpathlineto{\pgfqpoint{3.245841in}{3.111155in}}%
\pgfpathlineto{\pgfqpoint{3.300400in}{3.085918in}}%
\pgfpathlineto{\pgfqpoint{3.354958in}{3.058110in}}%
\pgfpathlineto{\pgfqpoint{3.409517in}{3.027626in}}%
\pgfpathlineto{\pgfqpoint{3.464076in}{2.994361in}}%
\pgfpathlineto{\pgfqpoint{3.513674in}{2.961616in}}%
\pgfpathlineto{\pgfqpoint{3.563273in}{2.926402in}}%
\pgfpathlineto{\pgfqpoint{3.612872in}{2.888636in}}%
\pgfpathlineto{\pgfqpoint{3.662471in}{2.848237in}}%
\pgfpathlineto{\pgfqpoint{3.712070in}{2.805119in}}%
\pgfpathlineto{\pgfqpoint{3.761668in}{2.759199in}}%
\pgfpathlineto{\pgfqpoint{3.811267in}{2.710390in}}%
\pgfpathlineto{\pgfqpoint{3.860866in}{2.658607in}}%
\pgfpathlineto{\pgfqpoint{3.910465in}{2.603763in}}%
\pgfpathlineto{\pgfqpoint{3.960063in}{2.545769in}}%
\pgfpathlineto{\pgfqpoint{4.009662in}{2.484539in}}%
\pgfpathlineto{\pgfqpoint{4.059261in}{2.419981in}}%
\pgfpathlineto{\pgfqpoint{4.108860in}{2.352007in}}%
\pgfpathlineto{\pgfqpoint{4.158459in}{2.280526in}}%
\pgfpathlineto{\pgfqpoint{4.208057in}{2.205446in}}%
\pgfpathlineto{\pgfqpoint{4.262616in}{2.118591in}}%
\pgfpathlineto{\pgfqpoint{4.317175in}{2.027147in}}%
\pgfpathlineto{\pgfqpoint{4.371733in}{1.930989in}}%
\pgfpathlineto{\pgfqpoint{4.426292in}{1.829991in}}%
\pgfpathlineto{\pgfqpoint{4.480851in}{1.724026in}}%
\pgfpathlineto{\pgfqpoint{4.535409in}{1.612966in}}%
\pgfpathlineto{\pgfqpoint{4.589968in}{1.496683in}}%
\pgfpathlineto{\pgfqpoint{4.644527in}{1.375048in}}%
\pgfpathlineto{\pgfqpoint{4.699085in}{1.247928in}}%
\pgfpathlineto{\pgfqpoint{4.753644in}{1.115193in}}%
\pgfpathlineto{\pgfqpoint{4.808203in}{0.976710in}}%
\pgfpathlineto{\pgfqpoint{4.862761in}{0.832344in}}%
\pgfpathlineto{\pgfqpoint{4.917320in}{0.681962in}}%
\pgfpathlineto{\pgfqpoint{4.971879in}{0.525428in}}%
\pgfpathlineto{\pgfqpoint{5.026437in}{0.362604in}}%
\pgfpathlineto{\pgfqpoint{5.085956in}{0.177644in}}%
\pgfpathlineto{\pgfqpoint{5.125635in}{0.050000in}}%
\pgfpathlineto{\pgfqpoint{5.125635in}{0.050000in}}%
\pgfusepath{stroke}%
\end{pgfscope}%
\begin{pgfscope}%
\pgfpathrectangle{\pgfqpoint{0.088135in}{0.050000in}}{\pgfqpoint{5.037500in}{3.253250in}}%
\pgfusepath{clip}%
\pgfsetbuttcap%
\pgfsetroundjoin%
\pgfsetlinewidth{1.003750pt}%
\definecolor{currentstroke}{rgb}{0.000000,0.000000,0.000000}%
\pgfsetstrokecolor{currentstroke}%
\pgfsetdash{{1.000000pt}{1.650000pt}}{0.000000pt}%
\pgfpathmoveto{\pgfqpoint{2.530201in}{2.988005in}}%
\pgfpathlineto{\pgfqpoint{2.530201in}{3.250709in}}%
\pgfusepath{stroke}%
\end{pgfscope}%
\begin{pgfscope}%
\pgfsetbuttcap%
\pgfsetmiterjoin%
\definecolor{currentfill}{rgb}{0.000000,0.000000,0.000000}%
\pgfsetfillcolor{currentfill}%
\pgfsetlinewidth{1.003750pt}%
\definecolor{currentstroke}{rgb}{0.000000,0.000000,0.000000}%
\pgfsetstrokecolor{currentstroke}%
\pgfsetdash{}{0pt}%
\pgfsys@defobject{currentmarker}{\pgfqpoint{-0.027778in}{-0.027778in}}{\pgfqpoint{0.027778in}{0.027778in}}{%
\pgfpathmoveto{\pgfqpoint{0.027778in}{-0.000000in}}%
\pgfpathlineto{\pgfqpoint{-0.027778in}{0.027778in}}%
\pgfpathlineto{\pgfqpoint{-0.027778in}{-0.027778in}}%
\pgfpathlineto{\pgfqpoint{0.027778in}{-0.000000in}}%
\pgfpathclose%
\pgfusepath{stroke,fill}%
}%
\begin{pgfscope}%
\pgfsys@transformshift{5.125635in}{2.988005in}%
\pgfsys@useobject{currentmarker}{}%
\end{pgfscope}%
\end{pgfscope}%
\begin{pgfscope}%
\pgfsetbuttcap%
\pgfsetmiterjoin%
\definecolor{currentfill}{rgb}{0.000000,0.000000,0.000000}%
\pgfsetfillcolor{currentfill}%
\pgfsetlinewidth{1.003750pt}%
\definecolor{currentstroke}{rgb}{0.000000,0.000000,0.000000}%
\pgfsetstrokecolor{currentstroke}%
\pgfsetdash{}{0pt}%
\pgfsys@defobject{currentmarker}{\pgfqpoint{-0.027778in}{-0.027778in}}{\pgfqpoint{0.027778in}{0.027778in}}{%
\pgfpathmoveto{\pgfqpoint{0.000000in}{0.027778in}}%
\pgfpathlineto{\pgfqpoint{-0.027778in}{-0.027778in}}%
\pgfpathlineto{\pgfqpoint{0.027778in}{-0.027778in}}%
\pgfpathlineto{\pgfqpoint{0.000000in}{0.027778in}}%
\pgfpathclose%
\pgfusepath{stroke,fill}%
}%
\begin{pgfscope}%
\pgfsys@transformshift{0.170717in}{3.303250in}%
\pgfsys@useobject{currentmarker}{}%
\end{pgfscope}%
\end{pgfscope}%
\begin{pgfscope}%
\pgfsetrectcap%
\pgfsetmiterjoin%
\pgfsetlinewidth{1.204500pt}%
\definecolor{currentstroke}{rgb}{0.000000,0.000000,0.000000}%
\pgfsetstrokecolor{currentstroke}%
\pgfsetdash{}{0pt}%
\pgfpathmoveto{\pgfqpoint{0.170717in}{0.050000in}}%
\pgfpathlineto{\pgfqpoint{0.170717in}{3.303250in}}%
\pgfusepath{stroke}%
\end{pgfscope}%
\begin{pgfscope}%
\pgfsetrectcap%
\pgfsetmiterjoin%
\pgfsetlinewidth{0.000000pt}%
\definecolor{currentstroke}{rgb}{0.000000,0.000000,0.000000}%
\pgfsetstrokecolor{currentstroke}%
\pgfsetstrokeopacity{0.000000}%
\pgfsetdash{}{0pt}%
\pgfpathmoveto{\pgfqpoint{5.125635in}{0.050000in}}%
\pgfpathlineto{\pgfqpoint{5.125635in}{3.303250in}}%
\pgfusepath{}%
\end{pgfscope}%
\begin{pgfscope}%
\pgfsetrectcap%
\pgfsetmiterjoin%
\pgfsetlinewidth{1.204500pt}%
\definecolor{currentstroke}{rgb}{0.000000,0.000000,0.000000}%
\pgfsetstrokecolor{currentstroke}%
\pgfsetdash{}{0pt}%
\pgfpathmoveto{\pgfqpoint{0.088135in}{2.988005in}}%
\pgfpathlineto{\pgfqpoint{5.125635in}{2.988005in}}%
\pgfusepath{stroke}%
\end{pgfscope}%
\begin{pgfscope}%
\pgfsetrectcap%
\pgfsetmiterjoin%
\pgfsetlinewidth{0.000000pt}%
\definecolor{currentstroke}{rgb}{0.000000,0.000000,0.000000}%
\pgfsetstrokecolor{currentstroke}%
\pgfsetstrokeopacity{0.000000}%
\pgfsetdash{}{0pt}%
\pgfpathmoveto{\pgfqpoint{0.088135in}{3.303250in}}%
\pgfpathlineto{\pgfqpoint{5.125635in}{3.303250in}}%
\pgfusepath{}%
\end{pgfscope}%
\begin{pgfscope}%
\definecolor{textcolor}{rgb}{0.000000,0.000000,0.000000}%
\pgfsetstrokecolor{textcolor}%
\pgfsetfillcolor{textcolor}%
\pgftext[x=0.088135in,y=3.368315in,,bottom]{\color{textcolor}{\rmfamily\fontsize{11.000000}{13.200000}\selectfont\catcode`\^=\active\def^{\ifmmode\sp\else\^{}\fi}\catcode`\%=\active\def%{\%}$y$}}%
\end{pgfscope}%
\begin{pgfscope}%
\definecolor{textcolor}{rgb}{0.000000,0.000000,0.000000}%
\pgfsetstrokecolor{textcolor}%
\pgfsetfillcolor{textcolor}%
\pgftext[x=5.226385in,y=1.091040in,left,]{\color{textcolor}{\rmfamily\fontsize{11.000000}{13.200000}\selectfont\catcode`\^=\active\def^{\ifmmode\sp\else\^{}\fi}\catcode`\%=\active\def%{\%}$L$}}%
\end{pgfscope}%
\begin{pgfscope}%
\definecolor{textcolor}{rgb}{0.000000,0.000000,0.000000}%
\pgfsetstrokecolor{textcolor}%
\pgfsetfillcolor{textcolor}%
\pgftext[x=3.709944in,y=3.119357in,left,base]{\color{textcolor}{\rmfamily\fontsize{11.000000}{13.200000}\selectfont\catcode`\^=\active\def^{\ifmmode\sp\else\^{}\fi}\catcode`\%=\active\def%{\%}$y = g(L)$}}%
\end{pgfscope}%
\end{pgfpicture}%
\makeatother%
\endgroup%

	\caption{Graph of \(g(L)\).}\label{fig:GraphOfGL}
\end{figure}

\subsubsection{Analytical Characterisation of the Put Price}

We have
\begin{align} \label{eq:PiecewiseValue}
	v_{L_*}(x) = \begin{cases} K - x, & 0 \le x \le L_*, \\ (K - L_*) \left( \frac{x}{L_*} \right)^{-\frac{2r}{\sigma^2}}, & x \ge L_*, \end{cases}
\end{align}
so that
\begin{align} \label{eq:PiecewiseDerivative}
	v'_{L_*}(x) = \begin{cases} -1, & 0 \le x \le L_*, \\ -(K - L_*) \frac{2r}{\sigma^2 x} \left( \frac{x}{L_*} \right)^{-\frac{2r}{\sigma^2}}, & x \ge L_*. \end{cases}
\end{align}
If we evaluate the second line in \Cref{eq:PiecewiseDerivative} at \(x = L_*\), we get the right-hand derivative
\[
	v'_{L_*}(L_*+) = - \frac{2r}{\sigma^2 L_*} (K - L_*) = - \frac{2rK}{\sigma^2 L_*} + \frac{2r}{\sigma^2} = - \frac{2r}{\sigma^2} \cdot \frac{2r + \sigma^2}{2r} + \frac{2r}{\sigma^2} = -1,
\]
which agrees with the left-hand derivative \(v'_{L_*}(L_*-) = -1\) provided by the first line in \Cref{eq:PiecewiseDerivative}. The derivative of \(v_{L_*}(x)\) is continuous at \(x = L_*\). This is known as \textit{smooth pasting}. The two parts of the definition of \(v_{L_*}(x)\) fit together at \(x = L_*\) so that both \(v_{L_*}(x)\) and \(v'_{L_*}(x)\) are continuous. This is because the graph of the function \(y = (K - L_*)(\frac{x}{L_*})^{-\frac{2r}{\sigma^2}}\) is tangent to the line \(y = K - x\) at \(x = L_*\), as one can see from \Cref{fig:ExerciseBoundaryValue}. In fact, we could have used the smooth pasting condition to solve for \(L_*\).

The second derivative of \(v(x)\) has a jump at \(x = L_*\), and hence is undefined at this point. Indeed,
\begin{align} \label{eq:PiecewiseSecondDeriv}
	v''_{L_*}(x) = \begin{cases} 0, & 0 \le x < L_*, \\ (K - L_*) \frac{2r(2r + \sigma^2)}{\sigma^4 x^2} \left( \frac{x}{L_*} \right)^{-\frac{2r}{\sigma^2}}, & x > L_*. \end{cases}
\end{align}
The left-hand and right-hand second derivatives at \(x = L_*\) are \(v''(L_*-) = 0\) and \(v''(L_*+) = (K - L_*) \frac{2r(2r+\sigma^2)}{\sigma^4 L_*^2} > 0\).

For \(x > L_*\), we can verify by direct computation that
\begin{align} \label{eq:OperatorCheckActionRegion}
	rv_{L_*}(x) - rxv'_{L_*}(x) - \frac{1}{2}\sigma^2 x^2 v''_{L_*}(x) \notag \\
	= (K - L_*) \left( r + \frac{2r^2}{\sigma^2} - \frac{r(2r + \sigma^2)}{\sigma^2} \right) \left( \frac{x}{L_*} \right)^{-\frac{2r}{\sigma^2}} = 0.
\end{align}

On the other hand, for \(0 \le x < L_*\),
\begin{align} \label{eq:OperatorCheckWaitRegion}
	rv_{L_*}(x) - rxv'_{L_*}(x) - \frac{1}{2}\sigma^2 x^2 v''_{L_*}(x) = r(K - x) + rx = rK.
\end{align}

In particular, we see that \(v_{L_*}(x)\) satisfies the so-called \textit{linear complementarity conditions}
\begin{align}
	v(x) \ge (K - x)^+ \text{ for all } x \ge 0, \label{eq:LCCPositivity}                                                 \\
	rv(x) - rxv'(x) - \frac{1}{2}\sigma^2 x^2 v''(x) \ge 0 \text{ for all } x \ge 0, \text{ and} \label{eq:LCCInequality} \\
	\text{for each } x \ge 0, \text{ equality holds in either \Cref{eq:LCCPositivity} or \Cref{eq:LCCInequality}.} \label{eq:LCCSlackness}
\end{align}
The point \(L_*\) is slightly problematical in \Cref{eq:LCCInequality} since \(v''_{L_*}(L_*)\) is undefined. However, if we replace \(v''_{L_*}(L_*)\) in \Cref{eq:LCCInequality} by either \(v''_{L_*}(L_*-)\) or \(v''_{L_*}(L_*+)\), the inequality holds.

The linear complementarity conditions \Crefrange{eq:LCCPositivity}{eq:LCCSlackness} determine the function \(v_{L_*}(x)\). More precisely, the function \(v_{L_*}(x)\) given by \Cref{eq:PiecewiseValue} is the only bounded continuous function having a continuous derivative that satisfies these conditions.

\subsubsection{Probabilistic Characterisation of the Put Price}

\begin{thm}{Stock Price Supermartingale }{StockPriceSupermartingale}
	Let \(S(t)\) be the stock price given by \Cref{eq:GBMRiskNeutral} and let \(\tau_{L_*}\) be given by \Cref{eq:StoppingTime} with \(L = L_*\). Then \(e^{-rt}v_{L_*}(S(t))\) is a supermartingale under \(\widetilde{\Prob}\), and the stopped process \(e^{-r(t \land \tau_{L_*})}v_{L_*}(S(t \land \tau_{L_*}))\) is a martingale.
\end{thm}

\begin{proof}
	Fortunately, the \Ito{} formula applies to functions whose second derivatives have jumps, provided the first derivative is continuous. We may thus compute
	\begin{align*}
		d\left[e^{-rt}v_{L_*}(S(t))\right]
		 & = e^{-rt}\left[-rv_{L_*}(S(t)) \, dt + v'_{L_*}(S(t)) \, dS(t) + \frac{1}{2}v''_{L_*}(S(t)) \, dS(t) \, dS(t)\right] \\
		 & = e^{-rt}\left[-rv_{L_*}(S(t)) + rS(t)v'_{L_*}(S(t)) + \frac{1}{2}\sigma^2 S^2(t)v''_{L_*}(S(t))\right] \, dt        \\
		 & \quad + e^{-rt}\sigma S(t)v'_{L_*}(S(t)) \, d\widetilde{W}(t).
	\end{align*}
	Because of \Crefrange{eq:OperatorCheckActionRegion}{eq:OperatorCheckWaitRegion}, the \(dt\) term in this expression is either \(0\) or \(-rK\), depending on whether \(S(t) > L_*\) or \(S(t) < L_*\). If \(S(t) = L_*\), \(v''_{L_*}(S(t))\) is undefined, but the probability \(S(t) = L_*\) is zero so this does not matter. We thus have
	\begin{equation} \label{eq:SupermartingaleDifferential}
		d\left[e^{-rt}v_{L_*}(S(t))\right] = -e^{-rt}rK\mathbf{I}_{\{S(t)<L_*\}} \, dt + e^{-rt}\sigma S(t)v'_{L_*}(S(t)) \, d\widetilde{W}(t).
	\end{equation}
	Because the \(dt\) term in \Cref{eq:SupermartingaleDifferential} is less than or equal to zero, \(e^{-rt}v_{L_*}(S(t))\) is a supermartingale; when \(S(t) < L_*\) it has a downward tendency. If the initial stock price is above \(L_*\), then prior to the time \(\tau_{L_*}\) when the stock price first reaches \(L_*\), the \(dt\) term in \Cref{eq:SupermartingaleDifferential} is zero and hence \(e^{-r(t \land \tau_{L_*})}v(S(t \land \tau_{L_*}))\) is a martingale. Indeed, integration of \Cref{eq:SupermartingaleDifferential} yields
	\[
		e^{-r(t\land\tau_{L_*})}v_{L_*}(S(t \land \tau_{L_*})) = v_{L_*}(S(0)) + \int_0^{t\land\tau_{L_*}} e^{-ru}\sigma S(u)v'_{L_*}(S(u)) \, d\widetilde{W}(u).
	\]
	\Ito\ integrals are martingales, and hence the \Ito\ integral above stopped at the stopping time \(\tau_{L_*}\) is a martingale.
\end{proof}

\begin{cor}{Perpetual American Put Price}{PerpetualAmericanPutPrice}
	Recall that \(\mathcal{T}\) is the set of all stopping times, not just those of the form (8.3.9). We have
	\[
		v_{L_*}(x) = \max_{\tau \in \mathcal{T}} \widetilde{\E} \left[e^{-r\tau} (K - S(\tau))\right],
	\]
	where \(x = S(0)\) is the initial stock price. In other words, \(v_{L_*}(x)\) is the perpetual American put price of \Cref{defn:PerpetualAmericanPut}.
\end{cor}

\begin{proof}
	Because \(e^{-rt}v_{L_*}(S(t))\) is a supermartingale under \(\widetilde{\Prob}\), we have from \Cref{thm:DoobsStoppingTheoremForContinuousMartingales} that, for every stopping time \(\tau \in \mathcal{T}\),
	\begin{equation} \label{SupermartingaleOptionalSampling}
		v_{L_*}(x) = v_{L_*}(S(0)) \ge \widetilde{\E} \left[e^{-r(t\land\tau)}v_{L_*}(S(t \land \tau))\right].
	\end{equation}
	Because \(v_{L_*}(S(t \land \tau))\) is bounded, we may let \(t \to \infty\) in \Cref{SupermartingaleOptionalSampling}, using the Dominated Convergence Theorem~\Cref{thm:DominatedAndBoundedConvergence}, to conclude that
	\[
		v_{L_*}(x) \ge \widetilde{\E} \left[e^{-r\tau}v_{L_*}(S(\tau))\right] \ge \widetilde{\E} \left[e^{-r\tau}(K - S(\tau))\right],
	\]
	where we have gotten the last inequality from \Cref{eq:LCCPositivity}. Because this inequality holds for every \(\tau \in \mathcal{T}\), we have
	\[
		v_{L_*}(x) \ge \max_{\tau \in \mathcal{T}} \widetilde{\E} \left[e^{-r\tau}(K - S(\tau))\right].
	\]

	On the other hand, if we replace \(\tau\) by \(\tau_{L_*}\), we obtain equality in \Cref{SupermartingaleOptionalSampling} because \(e^{-r(t\land\tau_{L_*})}v(S(t \land \tau_{L_*}))\) is a martingale under \(\widetilde{\Prob}\). Letting \(t \to \infty\) and using the Dominated Convergence Theorem~\Cref{thm:DominatedAndBoundedConvergence}, we obtain
	\[
		v_{L_*}(x) = \widetilde{\E} \left[e^{-r\tau_{L_*}}v_{L_*}(S(\tau_{L_*}))\right].
	\]
	Since
	\[
		e^{-r\tau_{L_*}}v_{L_*}(S(\tau_{L_*})) = e^{-r\tau_{L_*}}v_{L_*}(L_*) = e^{-r\tau_{L_*}}(K - L_*) = e^{-r\tau_{L_*}}(K - S(\tau_{L_*}))
	\]
	if \(\tau_{L_*} < \infty\) (and is interpreted to be zero if \(\tau_{L_*} = \infty\)), we see that
	\begin{equation} \label{eq:ValueAtOptimalStopping}
		v_{L_*}(x) = \widetilde{\E} \left[e^{-r\tau_{L_*}}(K - S(\tau_{L_*}))\right].
	\end{equation}
	It follows that
	\[
		v_{L_*}(x) \le \max_{\tau \in \mathcal{T}} \widetilde{\E} \left[e^{-r\tau}(K - S(\tau))\right].
	\]
\end{proof}

Discounted European option prices are martingales under the risk-neutral probability measure. Discounted American option prices are martingales up to the time they should be exercised. If they are not exercised when they should be, they tend downward. Since a martingale is a special case of a supermartingale, and processes that tend downward are supermartingales, discounted American option prices are supermartingales. An agent who is short an American option can hedge that short position in the usual way during the time the discounted option price is a martingale. If the option is not exercised when it should be, then the agent can continue the hedge and take money off the table. The following corollary illustrates this for the perpetual American put of this section.

\begin{cor}{Perpetual American Put Hedging}{PerpetualAmericanPutHedging}
	Consider an agent with initial capital \(X(0) = v_{L_*}(S(0))\), the initial perpetual American put price. Suppose this agent uses the portfolio process \(\Delta(t) = v'_{L_*}(S(t))\) and consumes cash at rate \(C(t) = rK\mathbb{I}_{\{S(t)<L_*\}}\) (i.e., consumes cash at rate \(rK\) whenever \(S(t) < L_*\)). Then the value \(X(t)\) of the agent's portfolio agrees with the option price \(v_{L_*}(S(t))\) for all times \(t\) until the option is exercised. In particular, \(X(t) \ge (K - S(t))^+\) for all \(t\) until the option is exercised, so the agent can pay off a short option position regardless of when the option is exercised.
\end{cor}

\begin{proof}
	The differential of the agent's portfolio value process is
	\[
		dX(t) = \Delta(t) \, dS(t) + r(X(t) - \Delta(t)S(t)) \, dt - C(t) \, dt,
	\]
	so the differential of the discounted portfolio value process is
	\begin{align}
		d(e^{-rt}X(t)) & = e^{-rt}(-rX(t) \, dt + dX(t)) \nonumber                                                                  \\
		               & = e^{-rt}(\Delta(t) \, dS(t) - r\Delta(t)S(t) \, dt - C(t) \, dt) \nonumber                                \\
		               & = e^{-rt}(\Delta(t)\sigma S(t) \, d\widetilde{W}(t) - C(t) \, dt). \label{DiscountedPortfolioDifferential}
	\end{align}
	Substituting \(\Delta(t) = v'_{L_*}(S(t))\) and \(C(t) = rK\mathbb{I}_{\{S(t)<L_*\}}\) into \Cref{DiscountedPortfolioDifferential} and comparing it to \Cref{eq:SupermartingaleDifferential}, we see that \(d(e^{-rt}X(t)) = d[e^{-rt}v_{L_*}(S(t))]\). Integrating both sides of this equation and using the initial equality \(X(0) = v_{L_*}(S(0))\), we obtain \(X(t) = v_{L_*}(S(t))\) for all \(t\) prior to exercise.
\end{proof}

\begin{xtip}{Hedging Mechanics}{HedgingMechanics}
	During any period in which \(S(t) < L_*\), the agent in \Cref{cor:PerpetualAmericanPutHedging} has stock position \(\Delta(t) = v'_{L_*}(S(t)) = -1\) (i.e., is short one share of stock) and has a total portfolio value \(X(t) = v_{L_*}(S(t)) = K - S(t)\). Therefore, the agent has \(K\) invested in the money market. If the owner of the put exercises, the agent in \Cref{cor:PerpetualAmericanPutHedging} receives a share of stock, which covers his short position, and pays out \(K\) from his money market account. If the owner of the put does not exercise, the agent holds his position and consumes the interest from the money market investment (i.e., consumes cash at rate \(rK\) per unit time).
\end{xtip}

The argument in \Cref{cor:PerpetualAmericanPutHedging} applies generally. In a complete market, whenever some discounted price process is a supermartingale, it is possible to construct a hedging portfolio whose value tracks the price process. This portfolio may sometimes consume. In the case of the perpetual American put, the supermartingale property for the discounted put price follows from \Cref{eq:LCCInequality}. If, in addition, the price process dominates some intrinsic value (see \Cref{eq:LCCPositivity} for the perpetual American put), then a short position in the American option with that intrinsic value can be hedged. There are always two conditions on the price of any American option, corresponding to \Cref{eq:LCCPositivity} and \Cref{eq:LCCInequality}. These conditions guarantee that the price is sufficient to satisfy the seller of the put.

However, \Crefrange{eq:LCCPositivity}{eq:LCCInequality} alone are not enough to determine the price of the perpetual American put. There can be functions that satisfy these conditions but are strictly greater than the price \(v_{L_*}(x)\) we constructed in \Cref{eq:PiecewiseValue}. There must be some additional condition that guarantees that the price is satisfactory for the purchaser of the put. One version of this condition for the perpetual American put is \Cref{eq:LCCSlackness}, guaranteeing that there exists an exercise strategy that permits the owner of the put to capture the full value of the put. It says that if we divide the half-line \([0,\infty)\) into two sets, the \textit{stopping set}

\begin{equation} \label{StoppingSetDefinition}
	S = \set[\big]{x \ge 0; v_{L_*}(x) = (K-x)^+}
\end{equation}
and the \textit{continuation set}
\begin{equation} \label{ContinuationSetDefinition}
	C = \set[\big]{x \ge 0; v_{L_*}(x) > (K-x)^+},
\end{equation}
then equality holds in \Cref{eq:LCCInequality} for \(x \in C\). If the initial stock price is in \(S\), then the owner of the put can get full value by exercising it immediately. On the other hand, if the initial stock price is in \(C\), then the put is more valuable than its intrinsic value, and the owner of the put can capture this extra value by waiting until the stock price enters \(S\) to exercise, if it ever does enter \(S\). The time of entry into the set \(S\) is in fact \(\tau_{L_*}\) in \Cref{thm:StockPriceSupermartingale}. We saw in \Cref{eq:ValueAtOptimalStopping} that
\begin{equation*}
	v(S(0)) = \tilde{\E} \left[ e^{-r\tau^*} v(S(\tau^*)) \right] = \tilde{\E} \left[ e^{-r\tau^*} (K - S(\tau^*)) \right].
\end{equation*}
In conclusion, the three linear complementarity conditions have counterparts that can be stated probabilistically rather than analytically (i.e., without writing conditions on the derivatives of \(v(x)\)). Let \(V(t) = e^{-rt}v(S(t))\) be the value of the perpetual American put. The stochastic process \(V(t)\) satisfies the following three conditions:
\begin{enumerate}[label=(\roman*)]
	\item \(V(t) \ge (K - S(t))^+\) for all \(t \ge 0\),
	\item \(e^{-rt}V(t)\) is a supermartingale under \(\tilde{\Prob}\), and
	\item there exists a stopping time \(\tau_*\) such that
\end{enumerate}
\begin{equation*}
	V(0) = \tilde{\E} \left[ e^{-r\tau_*} (K - S(\tau_*))^+ \right].
\end{equation*}
These three conditions determine the value of \(V(0)\).

\subsection{Finite-Expiration American Put}

\subsubsection{Analytical Characterisation of the Put Price}

\subsubsection{Probabilistic Characterisation of the Put Price}

\subsection{American Call}

We treat the American call, first on the usual geometric Brownian motion asset, \Cref{eq:GBMRiskNeutral} and then on a variation of this asset that pays dividends at discrete dates. In the first case, we see that the American call price is the same as the European call price. In the second case, we provide a recursion formula for computing the American call price.

\subsubsection{Underlying Asset Pays No Dividends}

We begin with a case slightly more general than a call option. Consider a stock whose price process \(S(t)\) is given by
\begin{equation}
	\label{eq:GeneralGBMRiskNeutral}
	dS(t) = rS(t) \, \di t + \sigma S(t) \, \di \widetilde{W}(t),
\end{equation}
where the interest rate \(r\) and the volatility \(\sigma\) are strictly positive and \(\widetilde{W}(t)\) is a Brownian motion under the risk-neutral probability measure \(\widetilde{\Prob}\).

\begin{lem}{Discounted Intrinsic Value Submartingale}{DiscountedIntrinsicValueSubmartingale}
	Let \(h(x)\) be a nonnegative, convex function of \(x \ge 0\) satisfying \(h(0) = 0\). Then the discounted intrinsic value \(e^{-rt}h(S(t))\) of the American derivative security that pays \(h(S(t))\) upon exercise is a submartingale.
\end{lem}

\begin{proof}
	Because \(h(x)\) is convex, for \(0 \le \lambda \le 1\) and \(0 \le x_1 \le x_2\), we have
	\begin{equation}
		h((1-\lambda)x_1 + \lambda x_2) \le (1-\lambda)h(x_1) + \lambda h(x_2). \label{eq:ConvexityInequalityForH}
	\end{equation}
	See Figure 8.5.1 for the case of a call payoff, \(h(x) = (x-K)^+\).
	Taking \(x_1 = 0\), \(x_2 = x\), and using the fact that \(h(0) = 0\), we obtain from \Cref{eq:ConvexityInequalityForH} that
	\begin{equation}
		h(\lambda x) \le \lambda h(x) \quad \text{for all } x \ge 0,\ 0 \le \lambda \le 1. \label{eq:ScalingInequalityForConvexH}
	\end{equation}
	For \(0 \le u \le t \le T\), we have \(0 \le e^{-r(t-u)} \le 1\), and \Cref{eq:ScalingInequalityForConvexH} implies
	\begin{equation}
		\widetilde{\E}\left[e^{-r(t-u)}h(S(t)) \middle| \cF(u)\right] \ge \widetilde{\E}\left[h\left(e^{-r(t-u)}S(t)\right) \middle| \cF(u)\right]. \label{eq:ConditionalExpectationMonotonicity}
	\end{equation}
	The conditional Jensen's inequality~\Cref{lem:JensensInequality} implies
	\begin{equation}
		\begin{aligned}
			\widetilde{\E}\left[h\left(e^{-r(t-u)}S(t)\right) \middle| \cF(u)\right] & \ge h\left(\widetilde{\E}\left[e^{-r(t-u)}S(t) \middle| \cF(u)\right]\right) \\ = h\left(e^{ru}\widetilde{\E}\left[e^{-rt}S(t) \middle| \cF(u)\right]\right).
		\end{aligned}
		\label{eq:ConditionalJensenInequality}
	\end{equation}
	Because \(e^{-rt}S(t)\) is a martingale under \(\widetilde{\Prob}\), we have
	\begin{equation}
		h\left(e^{ru}\widetilde{\E}\left[e^{-rt}S(t) \middle| \cF(u)\right]\right) = h\left(e^{ru}e^{-ru}S(u)\right) = h(S(u)). \label{eq:MartingaleSubstitutionStep}
	\end{equation}
	Putting \Crefrange{eq:ConditionalExpectationMonotonicity}{eq:MartingaleSubstitutionStep} together, we conclude that
	\begin{equation}
		\widetilde{\E}\left[e^{-r(t-u)}h(S(t)) \middle| \cF(u)\right] \ge h(S(u)) \label{eq:PreDiscountedSubmartingaleInequality}
	\end{equation}
	or, equivalently,
	\begin{equation}
		\widetilde{\E}\left[e^{-rt}h(S(t)) \middle| \cF(u)\right] \ge e^{-ru}h(S(u)). \label{eq:DiscountedIntrinsicSubmartingaleProperty}
	\end{equation}
	This is the submartingale property for \(e^{-rt}h(S(t))\).
\end{proof}

\begin{figure}[H]
	\centering
	%% Creator: Matplotlib, PGF backend
%%
%% To include the figure in your LaTeX document, write
%%   \input{<filename>.pgf}
%%
%% Make sure the required packages are loaded in your preamble
%%   \usepackage{pgf}
%%
%% Also ensure that all the required font packages are loaded; for instance,
%% the lmodern package is sometimes necessary when using math font.
%%   \usepackage{lmodern}
%%
%% Figures using additional raster images can only be included by \input if
%% they are in the same directory as the main LaTeX file. For loading figures
%% from other directories you can use the `import` package
%%   \usepackage{import}
%%
%% and then include the figures with
%%   \import{<path to file>}{<filename>.pgf}
%%
%% Matplotlib used the following preamble
%%   \def\mathdefault#1{#1}
%%   \everymath=\expandafter{\the\everymath\displaystyle}
%%   \IfFileExists{scrextend.sty}{
%%     \usepackage[fontsize=10.000000pt]{scrextend}
%%   }{
%%     \renewcommand{\normalsize}{\fontsize{10.000000}{12.000000}\selectfont}
%%     \normalsize
%%   }
%%   \usepackage{amsmath}
%%   \usepackage{amssymb}
%%   \usepackage{amsfonts}
%%   \usepackage{libertine}
%%   \usepackage[libertine]{newtxmath}
%%   \makeatletter\@ifpackageloaded{underscore}{}{\usepackage[strings]{underscore}}\makeatother
%%
\begingroup%
\makeatletter%
\begin{pgfpicture}%
\pgfpathrectangle{\pgfpointorigin}{\pgfqpoint{6.195516in}{3.990743in}}%
\pgfusepath{use as bounding box, clip}%
\begin{pgfscope}%
\pgfsetbuttcap%
\pgfsetmiterjoin%
\definecolor{currentfill}{rgb}{1.000000,1.000000,1.000000}%
\pgfsetfillcolor{currentfill}%
\pgfsetlinewidth{0.000000pt}%
\definecolor{currentstroke}{rgb}{1.000000,1.000000,1.000000}%
\pgfsetstrokecolor{currentstroke}%
\pgfsetdash{}{0pt}%
\pgfpathmoveto{\pgfqpoint{0.000000in}{0.000000in}}%
\pgfpathlineto{\pgfqpoint{6.195516in}{0.000000in}}%
\pgfpathlineto{\pgfqpoint{6.195516in}{3.990743in}}%
\pgfpathlineto{\pgfqpoint{0.000000in}{3.990743in}}%
\pgfpathlineto{\pgfqpoint{0.000000in}{0.000000in}}%
\pgfpathclose%
\pgfusepath{fill}%
\end{pgfscope}%
\begin{pgfscope}%
\pgfsetbuttcap%
\pgfsetmiterjoin%
\definecolor{currentfill}{rgb}{1.000000,1.000000,1.000000}%
\pgfsetfillcolor{currentfill}%
\pgfsetlinewidth{0.000000pt}%
\definecolor{currentstroke}{rgb}{0.000000,0.000000,0.000000}%
\pgfsetstrokecolor{currentstroke}%
\pgfsetstrokeopacity{0.000000}%
\pgfsetdash{}{0pt}%
\pgfpathmoveto{\pgfqpoint{1.355382in}{0.155832in}}%
\pgfpathlineto{\pgfqpoint{6.103850in}{0.155832in}}%
\pgfpathlineto{\pgfqpoint{6.103850in}{3.899076in}}%
\pgfpathlineto{\pgfqpoint{1.355382in}{3.899076in}}%
\pgfpathlineto{\pgfqpoint{1.355382in}{0.155832in}}%
\pgfpathclose%
\pgfusepath{fill}%
\end{pgfscope}%
\begin{pgfscope}%
\pgfsetbuttcap%
\pgfsetroundjoin%
\definecolor{currentfill}{rgb}{0.000000,0.000000,0.000000}%
\pgfsetfillcolor{currentfill}%
\pgfsetlinewidth{1.003750pt}%
\definecolor{currentstroke}{rgb}{0.000000,0.000000,0.000000}%
\pgfsetstrokecolor{currentstroke}%
\pgfsetdash{}{0pt}%
\pgfsys@defobject{currentmarker}{\pgfqpoint{0.000000in}{-0.027778in}}{\pgfqpoint{0.000000in}{0.027778in}}{%
\pgfpathmoveto{\pgfqpoint{0.000000in}{-0.027778in}}%
\pgfpathlineto{\pgfqpoint{0.000000in}{0.027778in}}%
\pgfusepath{stroke,fill}%
}%
\begin{pgfscope}%
\pgfsys@transformshift{2.387658in}{0.405382in}%
\pgfsys@useobject{currentmarker}{}%
\end{pgfscope}%
\end{pgfscope}%
\begin{pgfscope}%
\pgfsetbuttcap%
\pgfsetroundjoin%
\definecolor{currentfill}{rgb}{0.000000,0.000000,0.000000}%
\pgfsetfillcolor{currentfill}%
\pgfsetlinewidth{1.003750pt}%
\definecolor{currentstroke}{rgb}{0.000000,0.000000,0.000000}%
\pgfsetstrokecolor{currentstroke}%
\pgfsetdash{}{0pt}%
\pgfsys@defobject{currentmarker}{\pgfqpoint{0.000000in}{-0.027778in}}{\pgfqpoint{0.000000in}{0.027778in}}{%
\pgfpathmoveto{\pgfqpoint{0.000000in}{-0.027778in}}%
\pgfpathlineto{\pgfqpoint{0.000000in}{0.027778in}}%
\pgfusepath{stroke,fill}%
}%
\begin{pgfscope}%
\pgfsys@transformshift{2.387658in}{3.899076in}%
\pgfsys@useobject{currentmarker}{}%
\end{pgfscope}%
\end{pgfscope}%
\begin{pgfscope}%
\definecolor{textcolor}{rgb}{0.000000,0.000000,0.000000}%
\pgfsetstrokecolor{textcolor}%
\pgfsetfillcolor{textcolor}%
\pgftext[x=2.387658in,y=0.328993in,,top]{\color{textcolor}{\rmfamily\fontsize{10.000000}{12.000000}\selectfont\catcode`\^=\active\def^{\ifmmode\sp\else\^{}\fi}\catcode`\%=\active\def%{\%}$x_1$}}%
\end{pgfscope}%
\begin{pgfscope}%
\pgfsetbuttcap%
\pgfsetroundjoin%
\definecolor{currentfill}{rgb}{0.000000,0.000000,0.000000}%
\pgfsetfillcolor{currentfill}%
\pgfsetlinewidth{1.003750pt}%
\definecolor{currentstroke}{rgb}{0.000000,0.000000,0.000000}%
\pgfsetstrokecolor{currentstroke}%
\pgfsetdash{}{0pt}%
\pgfsys@defobject{currentmarker}{\pgfqpoint{0.000000in}{-0.027778in}}{\pgfqpoint{0.000000in}{0.027778in}}{%
\pgfpathmoveto{\pgfqpoint{0.000000in}{-0.027778in}}%
\pgfpathlineto{\pgfqpoint{0.000000in}{0.027778in}}%
\pgfusepath{stroke,fill}%
}%
\begin{pgfscope}%
\pgfsys@transformshift{3.419933in}{0.405382in}%
\pgfsys@useobject{currentmarker}{}%
\end{pgfscope}%
\end{pgfscope}%
\begin{pgfscope}%
\pgfsetbuttcap%
\pgfsetroundjoin%
\definecolor{currentfill}{rgb}{0.000000,0.000000,0.000000}%
\pgfsetfillcolor{currentfill}%
\pgfsetlinewidth{1.003750pt}%
\definecolor{currentstroke}{rgb}{0.000000,0.000000,0.000000}%
\pgfsetstrokecolor{currentstroke}%
\pgfsetdash{}{0pt}%
\pgfsys@defobject{currentmarker}{\pgfqpoint{0.000000in}{-0.027778in}}{\pgfqpoint{0.000000in}{0.027778in}}{%
\pgfpathmoveto{\pgfqpoint{0.000000in}{-0.027778in}}%
\pgfpathlineto{\pgfqpoint{0.000000in}{0.027778in}}%
\pgfusepath{stroke,fill}%
}%
\begin{pgfscope}%
\pgfsys@transformshift{3.419933in}{3.899076in}%
\pgfsys@useobject{currentmarker}{}%
\end{pgfscope}%
\end{pgfscope}%
\begin{pgfscope}%
\definecolor{textcolor}{rgb}{0.000000,0.000000,0.000000}%
\pgfsetstrokecolor{textcolor}%
\pgfsetfillcolor{textcolor}%
\pgftext[x=3.419933in,y=0.328993in,,top]{\color{textcolor}{\rmfamily\fontsize{10.000000}{12.000000}\selectfont\catcode`\^=\active\def^{\ifmmode\sp\else\^{}\fi}\catcode`\%=\active\def%{\%}$K$}}%
\end{pgfscope}%
\begin{pgfscope}%
\pgfsetbuttcap%
\pgfsetroundjoin%
\definecolor{currentfill}{rgb}{0.000000,0.000000,0.000000}%
\pgfsetfillcolor{currentfill}%
\pgfsetlinewidth{1.003750pt}%
\definecolor{currentstroke}{rgb}{0.000000,0.000000,0.000000}%
\pgfsetstrokecolor{currentstroke}%
\pgfsetdash{}{0pt}%
\pgfsys@defobject{currentmarker}{\pgfqpoint{0.000000in}{-0.027778in}}{\pgfqpoint{0.000000in}{0.027778in}}{%
\pgfpathmoveto{\pgfqpoint{0.000000in}{-0.027778in}}%
\pgfpathlineto{\pgfqpoint{0.000000in}{0.027778in}}%
\pgfusepath{stroke,fill}%
}%
\begin{pgfscope}%
\pgfsys@transformshift{4.245754in}{0.405382in}%
\pgfsys@useobject{currentmarker}{}%
\end{pgfscope}%
\end{pgfscope}%
\begin{pgfscope}%
\pgfsetbuttcap%
\pgfsetroundjoin%
\definecolor{currentfill}{rgb}{0.000000,0.000000,0.000000}%
\pgfsetfillcolor{currentfill}%
\pgfsetlinewidth{1.003750pt}%
\definecolor{currentstroke}{rgb}{0.000000,0.000000,0.000000}%
\pgfsetstrokecolor{currentstroke}%
\pgfsetdash{}{0pt}%
\pgfsys@defobject{currentmarker}{\pgfqpoint{0.000000in}{-0.027778in}}{\pgfqpoint{0.000000in}{0.027778in}}{%
\pgfpathmoveto{\pgfqpoint{0.000000in}{-0.027778in}}%
\pgfpathlineto{\pgfqpoint{0.000000in}{0.027778in}}%
\pgfusepath{stroke,fill}%
}%
\begin{pgfscope}%
\pgfsys@transformshift{4.245754in}{3.899076in}%
\pgfsys@useobject{currentmarker}{}%
\end{pgfscope}%
\end{pgfscope}%
\begin{pgfscope}%
\definecolor{textcolor}{rgb}{0.000000,0.000000,0.000000}%
\pgfsetstrokecolor{textcolor}%
\pgfsetfillcolor{textcolor}%
\pgftext[x=3.985693in, y=0.231597in, left, base]{\color{textcolor}{\rmfamily\fontsize{10.000000}{12.000000}\selectfont\catcode`\^=\active\def^{\ifmmode\sp\else\^{}\fi}\catcode`\%=\active\def%{\%}$(1-\lambda)x_1$}}%
\end{pgfscope}%
\begin{pgfscope}%
\definecolor{textcolor}{rgb}{0.000000,0.000000,0.000000}%
\pgfsetstrokecolor{textcolor}%
\pgfsetfillcolor{textcolor}%
\pgftext[x=4.104820in, y=0.082361in, left, base]{\color{textcolor}{\rmfamily\fontsize{10.000000}{12.000000}\selectfont\catcode`\^=\active\def^{\ifmmode\sp\else\^{}\fi}\catcode`\%=\active\def%{\%}$+\lambda x_2$}}%
\end{pgfscope}%
\begin{pgfscope}%
\pgfsetbuttcap%
\pgfsetroundjoin%
\definecolor{currentfill}{rgb}{0.000000,0.000000,0.000000}%
\pgfsetfillcolor{currentfill}%
\pgfsetlinewidth{1.003750pt}%
\definecolor{currentstroke}{rgb}{0.000000,0.000000,0.000000}%
\pgfsetstrokecolor{currentstroke}%
\pgfsetdash{}{0pt}%
\pgfsys@defobject{currentmarker}{\pgfqpoint{0.000000in}{-0.027778in}}{\pgfqpoint{0.000000in}{0.027778in}}{%
\pgfpathmoveto{\pgfqpoint{0.000000in}{-0.027778in}}%
\pgfpathlineto{\pgfqpoint{0.000000in}{0.027778in}}%
\pgfusepath{stroke,fill}%
}%
\begin{pgfscope}%
\pgfsys@transformshift{5.484484in}{0.405382in}%
\pgfsys@useobject{currentmarker}{}%
\end{pgfscope}%
\end{pgfscope}%
\begin{pgfscope}%
\pgfsetbuttcap%
\pgfsetroundjoin%
\definecolor{currentfill}{rgb}{0.000000,0.000000,0.000000}%
\pgfsetfillcolor{currentfill}%
\pgfsetlinewidth{1.003750pt}%
\definecolor{currentstroke}{rgb}{0.000000,0.000000,0.000000}%
\pgfsetstrokecolor{currentstroke}%
\pgfsetdash{}{0pt}%
\pgfsys@defobject{currentmarker}{\pgfqpoint{0.000000in}{-0.027778in}}{\pgfqpoint{0.000000in}{0.027778in}}{%
\pgfpathmoveto{\pgfqpoint{0.000000in}{-0.027778in}}%
\pgfpathlineto{\pgfqpoint{0.000000in}{0.027778in}}%
\pgfusepath{stroke,fill}%
}%
\begin{pgfscope}%
\pgfsys@transformshift{5.484484in}{3.899076in}%
\pgfsys@useobject{currentmarker}{}%
\end{pgfscope}%
\end{pgfscope}%
\begin{pgfscope}%
\definecolor{textcolor}{rgb}{0.000000,0.000000,0.000000}%
\pgfsetstrokecolor{textcolor}%
\pgfsetfillcolor{textcolor}%
\pgftext[x=5.484484in,y=0.328993in,,top]{\color{textcolor}{\rmfamily\fontsize{10.000000}{12.000000}\selectfont\catcode`\^=\active\def^{\ifmmode\sp\else\^{}\fi}\catcode`\%=\active\def%{\%}$x_2$}}%
\end{pgfscope}%
\begin{pgfscope}%
\pgfsetbuttcap%
\pgfsetroundjoin%
\definecolor{currentfill}{rgb}{0.000000,0.000000,0.000000}%
\pgfsetfillcolor{currentfill}%
\pgfsetlinewidth{0.501875pt}%
\definecolor{currentstroke}{rgb}{0.000000,0.000000,0.000000}%
\pgfsetstrokecolor{currentstroke}%
\pgfsetdash{}{0pt}%
\pgfsys@defobject{currentmarker}{\pgfqpoint{0.000000in}{0.000000in}}{\pgfqpoint{0.000000in}{0.020833in}}{%
\pgfpathmoveto{\pgfqpoint{0.000000in}{0.000000in}}%
\pgfpathlineto{\pgfqpoint{0.000000in}{0.020833in}}%
\pgfusepath{stroke,fill}%
}%
\begin{pgfscope}%
\pgfsys@transformshift{1.355382in}{0.405382in}%
\pgfsys@useobject{currentmarker}{}%
\end{pgfscope}%
\end{pgfscope}%
\begin{pgfscope}%
\pgfsetbuttcap%
\pgfsetroundjoin%
\definecolor{currentfill}{rgb}{0.000000,0.000000,0.000000}%
\pgfsetfillcolor{currentfill}%
\pgfsetlinewidth{0.501875pt}%
\definecolor{currentstroke}{rgb}{0.000000,0.000000,0.000000}%
\pgfsetstrokecolor{currentstroke}%
\pgfsetdash{}{0pt}%
\pgfsys@defobject{currentmarker}{\pgfqpoint{0.000000in}{-0.020833in}}{\pgfqpoint{0.000000in}{0.000000in}}{%
\pgfpathmoveto{\pgfqpoint{0.000000in}{0.000000in}}%
\pgfpathlineto{\pgfqpoint{0.000000in}{-0.020833in}}%
\pgfusepath{stroke,fill}%
}%
\begin{pgfscope}%
\pgfsys@transformshift{1.355382in}{3.899076in}%
\pgfsys@useobject{currentmarker}{}%
\end{pgfscope}%
\end{pgfscope}%
\begin{pgfscope}%
\pgfsetbuttcap%
\pgfsetroundjoin%
\definecolor{currentfill}{rgb}{0.000000,0.000000,0.000000}%
\pgfsetfillcolor{currentfill}%
\pgfsetlinewidth{0.501875pt}%
\definecolor{currentstroke}{rgb}{0.000000,0.000000,0.000000}%
\pgfsetstrokecolor{currentstroke}%
\pgfsetdash{}{0pt}%
\pgfsys@defobject{currentmarker}{\pgfqpoint{0.000000in}{0.000000in}}{\pgfqpoint{0.000000in}{0.020833in}}{%
\pgfpathmoveto{\pgfqpoint{0.000000in}{0.000000in}}%
\pgfpathlineto{\pgfqpoint{0.000000in}{0.020833in}}%
\pgfusepath{stroke,fill}%
}%
\begin{pgfscope}%
\pgfsys@transformshift{1.561837in}{0.405382in}%
\pgfsys@useobject{currentmarker}{}%
\end{pgfscope}%
\end{pgfscope}%
\begin{pgfscope}%
\pgfsetbuttcap%
\pgfsetroundjoin%
\definecolor{currentfill}{rgb}{0.000000,0.000000,0.000000}%
\pgfsetfillcolor{currentfill}%
\pgfsetlinewidth{0.501875pt}%
\definecolor{currentstroke}{rgb}{0.000000,0.000000,0.000000}%
\pgfsetstrokecolor{currentstroke}%
\pgfsetdash{}{0pt}%
\pgfsys@defobject{currentmarker}{\pgfqpoint{0.000000in}{-0.020833in}}{\pgfqpoint{0.000000in}{0.000000in}}{%
\pgfpathmoveto{\pgfqpoint{0.000000in}{0.000000in}}%
\pgfpathlineto{\pgfqpoint{0.000000in}{-0.020833in}}%
\pgfusepath{stroke,fill}%
}%
\begin{pgfscope}%
\pgfsys@transformshift{1.561837in}{3.899076in}%
\pgfsys@useobject{currentmarker}{}%
\end{pgfscope}%
\end{pgfscope}%
\begin{pgfscope}%
\pgfsetbuttcap%
\pgfsetroundjoin%
\definecolor{currentfill}{rgb}{0.000000,0.000000,0.000000}%
\pgfsetfillcolor{currentfill}%
\pgfsetlinewidth{0.501875pt}%
\definecolor{currentstroke}{rgb}{0.000000,0.000000,0.000000}%
\pgfsetstrokecolor{currentstroke}%
\pgfsetdash{}{0pt}%
\pgfsys@defobject{currentmarker}{\pgfqpoint{0.000000in}{0.000000in}}{\pgfqpoint{0.000000in}{0.020833in}}{%
\pgfpathmoveto{\pgfqpoint{0.000000in}{0.000000in}}%
\pgfpathlineto{\pgfqpoint{0.000000in}{0.020833in}}%
\pgfusepath{stroke,fill}%
}%
\begin{pgfscope}%
\pgfsys@transformshift{1.768292in}{0.405382in}%
\pgfsys@useobject{currentmarker}{}%
\end{pgfscope}%
\end{pgfscope}%
\begin{pgfscope}%
\pgfsetbuttcap%
\pgfsetroundjoin%
\definecolor{currentfill}{rgb}{0.000000,0.000000,0.000000}%
\pgfsetfillcolor{currentfill}%
\pgfsetlinewidth{0.501875pt}%
\definecolor{currentstroke}{rgb}{0.000000,0.000000,0.000000}%
\pgfsetstrokecolor{currentstroke}%
\pgfsetdash{}{0pt}%
\pgfsys@defobject{currentmarker}{\pgfqpoint{0.000000in}{-0.020833in}}{\pgfqpoint{0.000000in}{0.000000in}}{%
\pgfpathmoveto{\pgfqpoint{0.000000in}{0.000000in}}%
\pgfpathlineto{\pgfqpoint{0.000000in}{-0.020833in}}%
\pgfusepath{stroke,fill}%
}%
\begin{pgfscope}%
\pgfsys@transformshift{1.768292in}{3.899076in}%
\pgfsys@useobject{currentmarker}{}%
\end{pgfscope}%
\end{pgfscope}%
\begin{pgfscope}%
\pgfsetbuttcap%
\pgfsetroundjoin%
\definecolor{currentfill}{rgb}{0.000000,0.000000,0.000000}%
\pgfsetfillcolor{currentfill}%
\pgfsetlinewidth{0.501875pt}%
\definecolor{currentstroke}{rgb}{0.000000,0.000000,0.000000}%
\pgfsetstrokecolor{currentstroke}%
\pgfsetdash{}{0pt}%
\pgfsys@defobject{currentmarker}{\pgfqpoint{0.000000in}{0.000000in}}{\pgfqpoint{0.000000in}{0.020833in}}{%
\pgfpathmoveto{\pgfqpoint{0.000000in}{0.000000in}}%
\pgfpathlineto{\pgfqpoint{0.000000in}{0.020833in}}%
\pgfusepath{stroke,fill}%
}%
\begin{pgfscope}%
\pgfsys@transformshift{1.974748in}{0.405382in}%
\pgfsys@useobject{currentmarker}{}%
\end{pgfscope}%
\end{pgfscope}%
\begin{pgfscope}%
\pgfsetbuttcap%
\pgfsetroundjoin%
\definecolor{currentfill}{rgb}{0.000000,0.000000,0.000000}%
\pgfsetfillcolor{currentfill}%
\pgfsetlinewidth{0.501875pt}%
\definecolor{currentstroke}{rgb}{0.000000,0.000000,0.000000}%
\pgfsetstrokecolor{currentstroke}%
\pgfsetdash{}{0pt}%
\pgfsys@defobject{currentmarker}{\pgfqpoint{0.000000in}{-0.020833in}}{\pgfqpoint{0.000000in}{0.000000in}}{%
\pgfpathmoveto{\pgfqpoint{0.000000in}{0.000000in}}%
\pgfpathlineto{\pgfqpoint{0.000000in}{-0.020833in}}%
\pgfusepath{stroke,fill}%
}%
\begin{pgfscope}%
\pgfsys@transformshift{1.974748in}{3.899076in}%
\pgfsys@useobject{currentmarker}{}%
\end{pgfscope}%
\end{pgfscope}%
\begin{pgfscope}%
\pgfsetbuttcap%
\pgfsetroundjoin%
\definecolor{currentfill}{rgb}{0.000000,0.000000,0.000000}%
\pgfsetfillcolor{currentfill}%
\pgfsetlinewidth{0.501875pt}%
\definecolor{currentstroke}{rgb}{0.000000,0.000000,0.000000}%
\pgfsetstrokecolor{currentstroke}%
\pgfsetdash{}{0pt}%
\pgfsys@defobject{currentmarker}{\pgfqpoint{0.000000in}{0.000000in}}{\pgfqpoint{0.000000in}{0.020833in}}{%
\pgfpathmoveto{\pgfqpoint{0.000000in}{0.000000in}}%
\pgfpathlineto{\pgfqpoint{0.000000in}{0.020833in}}%
\pgfusepath{stroke,fill}%
}%
\begin{pgfscope}%
\pgfsys@transformshift{2.181203in}{0.405382in}%
\pgfsys@useobject{currentmarker}{}%
\end{pgfscope}%
\end{pgfscope}%
\begin{pgfscope}%
\pgfsetbuttcap%
\pgfsetroundjoin%
\definecolor{currentfill}{rgb}{0.000000,0.000000,0.000000}%
\pgfsetfillcolor{currentfill}%
\pgfsetlinewidth{0.501875pt}%
\definecolor{currentstroke}{rgb}{0.000000,0.000000,0.000000}%
\pgfsetstrokecolor{currentstroke}%
\pgfsetdash{}{0pt}%
\pgfsys@defobject{currentmarker}{\pgfqpoint{0.000000in}{-0.020833in}}{\pgfqpoint{0.000000in}{0.000000in}}{%
\pgfpathmoveto{\pgfqpoint{0.000000in}{0.000000in}}%
\pgfpathlineto{\pgfqpoint{0.000000in}{-0.020833in}}%
\pgfusepath{stroke,fill}%
}%
\begin{pgfscope}%
\pgfsys@transformshift{2.181203in}{3.899076in}%
\pgfsys@useobject{currentmarker}{}%
\end{pgfscope}%
\end{pgfscope}%
\begin{pgfscope}%
\pgfsetbuttcap%
\pgfsetroundjoin%
\definecolor{currentfill}{rgb}{0.000000,0.000000,0.000000}%
\pgfsetfillcolor{currentfill}%
\pgfsetlinewidth{0.501875pt}%
\definecolor{currentstroke}{rgb}{0.000000,0.000000,0.000000}%
\pgfsetstrokecolor{currentstroke}%
\pgfsetdash{}{0pt}%
\pgfsys@defobject{currentmarker}{\pgfqpoint{0.000000in}{0.000000in}}{\pgfqpoint{0.000000in}{0.020833in}}{%
\pgfpathmoveto{\pgfqpoint{0.000000in}{0.000000in}}%
\pgfpathlineto{\pgfqpoint{0.000000in}{0.020833in}}%
\pgfusepath{stroke,fill}%
}%
\begin{pgfscope}%
\pgfsys@transformshift{2.594113in}{0.405382in}%
\pgfsys@useobject{currentmarker}{}%
\end{pgfscope}%
\end{pgfscope}%
\begin{pgfscope}%
\pgfsetbuttcap%
\pgfsetroundjoin%
\definecolor{currentfill}{rgb}{0.000000,0.000000,0.000000}%
\pgfsetfillcolor{currentfill}%
\pgfsetlinewidth{0.501875pt}%
\definecolor{currentstroke}{rgb}{0.000000,0.000000,0.000000}%
\pgfsetstrokecolor{currentstroke}%
\pgfsetdash{}{0pt}%
\pgfsys@defobject{currentmarker}{\pgfqpoint{0.000000in}{-0.020833in}}{\pgfqpoint{0.000000in}{0.000000in}}{%
\pgfpathmoveto{\pgfqpoint{0.000000in}{0.000000in}}%
\pgfpathlineto{\pgfqpoint{0.000000in}{-0.020833in}}%
\pgfusepath{stroke,fill}%
}%
\begin{pgfscope}%
\pgfsys@transformshift{2.594113in}{3.899076in}%
\pgfsys@useobject{currentmarker}{}%
\end{pgfscope}%
\end{pgfscope}%
\begin{pgfscope}%
\pgfsetbuttcap%
\pgfsetroundjoin%
\definecolor{currentfill}{rgb}{0.000000,0.000000,0.000000}%
\pgfsetfillcolor{currentfill}%
\pgfsetlinewidth{0.501875pt}%
\definecolor{currentstroke}{rgb}{0.000000,0.000000,0.000000}%
\pgfsetstrokecolor{currentstroke}%
\pgfsetdash{}{0pt}%
\pgfsys@defobject{currentmarker}{\pgfqpoint{0.000000in}{0.000000in}}{\pgfqpoint{0.000000in}{0.020833in}}{%
\pgfpathmoveto{\pgfqpoint{0.000000in}{0.000000in}}%
\pgfpathlineto{\pgfqpoint{0.000000in}{0.020833in}}%
\pgfusepath{stroke,fill}%
}%
\begin{pgfscope}%
\pgfsys@transformshift{2.800568in}{0.405382in}%
\pgfsys@useobject{currentmarker}{}%
\end{pgfscope}%
\end{pgfscope}%
\begin{pgfscope}%
\pgfsetbuttcap%
\pgfsetroundjoin%
\definecolor{currentfill}{rgb}{0.000000,0.000000,0.000000}%
\pgfsetfillcolor{currentfill}%
\pgfsetlinewidth{0.501875pt}%
\definecolor{currentstroke}{rgb}{0.000000,0.000000,0.000000}%
\pgfsetstrokecolor{currentstroke}%
\pgfsetdash{}{0pt}%
\pgfsys@defobject{currentmarker}{\pgfqpoint{0.000000in}{-0.020833in}}{\pgfqpoint{0.000000in}{0.000000in}}{%
\pgfpathmoveto{\pgfqpoint{0.000000in}{0.000000in}}%
\pgfpathlineto{\pgfqpoint{0.000000in}{-0.020833in}}%
\pgfusepath{stroke,fill}%
}%
\begin{pgfscope}%
\pgfsys@transformshift{2.800568in}{3.899076in}%
\pgfsys@useobject{currentmarker}{}%
\end{pgfscope}%
\end{pgfscope}%
\begin{pgfscope}%
\pgfsetbuttcap%
\pgfsetroundjoin%
\definecolor{currentfill}{rgb}{0.000000,0.000000,0.000000}%
\pgfsetfillcolor{currentfill}%
\pgfsetlinewidth{0.501875pt}%
\definecolor{currentstroke}{rgb}{0.000000,0.000000,0.000000}%
\pgfsetstrokecolor{currentstroke}%
\pgfsetdash{}{0pt}%
\pgfsys@defobject{currentmarker}{\pgfqpoint{0.000000in}{0.000000in}}{\pgfqpoint{0.000000in}{0.020833in}}{%
\pgfpathmoveto{\pgfqpoint{0.000000in}{0.000000in}}%
\pgfpathlineto{\pgfqpoint{0.000000in}{0.020833in}}%
\pgfusepath{stroke,fill}%
}%
\begin{pgfscope}%
\pgfsys@transformshift{3.007023in}{0.405382in}%
\pgfsys@useobject{currentmarker}{}%
\end{pgfscope}%
\end{pgfscope}%
\begin{pgfscope}%
\pgfsetbuttcap%
\pgfsetroundjoin%
\definecolor{currentfill}{rgb}{0.000000,0.000000,0.000000}%
\pgfsetfillcolor{currentfill}%
\pgfsetlinewidth{0.501875pt}%
\definecolor{currentstroke}{rgb}{0.000000,0.000000,0.000000}%
\pgfsetstrokecolor{currentstroke}%
\pgfsetdash{}{0pt}%
\pgfsys@defobject{currentmarker}{\pgfqpoint{0.000000in}{-0.020833in}}{\pgfqpoint{0.000000in}{0.000000in}}{%
\pgfpathmoveto{\pgfqpoint{0.000000in}{0.000000in}}%
\pgfpathlineto{\pgfqpoint{0.000000in}{-0.020833in}}%
\pgfusepath{stroke,fill}%
}%
\begin{pgfscope}%
\pgfsys@transformshift{3.007023in}{3.899076in}%
\pgfsys@useobject{currentmarker}{}%
\end{pgfscope}%
\end{pgfscope}%
\begin{pgfscope}%
\pgfsetbuttcap%
\pgfsetroundjoin%
\definecolor{currentfill}{rgb}{0.000000,0.000000,0.000000}%
\pgfsetfillcolor{currentfill}%
\pgfsetlinewidth{0.501875pt}%
\definecolor{currentstroke}{rgb}{0.000000,0.000000,0.000000}%
\pgfsetstrokecolor{currentstroke}%
\pgfsetdash{}{0pt}%
\pgfsys@defobject{currentmarker}{\pgfqpoint{0.000000in}{0.000000in}}{\pgfqpoint{0.000000in}{0.020833in}}{%
\pgfpathmoveto{\pgfqpoint{0.000000in}{0.000000in}}%
\pgfpathlineto{\pgfqpoint{0.000000in}{0.020833in}}%
\pgfusepath{stroke,fill}%
}%
\begin{pgfscope}%
\pgfsys@transformshift{3.213478in}{0.405382in}%
\pgfsys@useobject{currentmarker}{}%
\end{pgfscope}%
\end{pgfscope}%
\begin{pgfscope}%
\pgfsetbuttcap%
\pgfsetroundjoin%
\definecolor{currentfill}{rgb}{0.000000,0.000000,0.000000}%
\pgfsetfillcolor{currentfill}%
\pgfsetlinewidth{0.501875pt}%
\definecolor{currentstroke}{rgb}{0.000000,0.000000,0.000000}%
\pgfsetstrokecolor{currentstroke}%
\pgfsetdash{}{0pt}%
\pgfsys@defobject{currentmarker}{\pgfqpoint{0.000000in}{-0.020833in}}{\pgfqpoint{0.000000in}{0.000000in}}{%
\pgfpathmoveto{\pgfqpoint{0.000000in}{0.000000in}}%
\pgfpathlineto{\pgfqpoint{0.000000in}{-0.020833in}}%
\pgfusepath{stroke,fill}%
}%
\begin{pgfscope}%
\pgfsys@transformshift{3.213478in}{3.899076in}%
\pgfsys@useobject{currentmarker}{}%
\end{pgfscope}%
\end{pgfscope}%
\begin{pgfscope}%
\pgfsetbuttcap%
\pgfsetroundjoin%
\definecolor{currentfill}{rgb}{0.000000,0.000000,0.000000}%
\pgfsetfillcolor{currentfill}%
\pgfsetlinewidth{0.501875pt}%
\definecolor{currentstroke}{rgb}{0.000000,0.000000,0.000000}%
\pgfsetstrokecolor{currentstroke}%
\pgfsetdash{}{0pt}%
\pgfsys@defobject{currentmarker}{\pgfqpoint{0.000000in}{0.000000in}}{\pgfqpoint{0.000000in}{0.020833in}}{%
\pgfpathmoveto{\pgfqpoint{0.000000in}{0.000000in}}%
\pgfpathlineto{\pgfqpoint{0.000000in}{0.020833in}}%
\pgfusepath{stroke,fill}%
}%
\begin{pgfscope}%
\pgfsys@transformshift{3.626388in}{0.405382in}%
\pgfsys@useobject{currentmarker}{}%
\end{pgfscope}%
\end{pgfscope}%
\begin{pgfscope}%
\pgfsetbuttcap%
\pgfsetroundjoin%
\definecolor{currentfill}{rgb}{0.000000,0.000000,0.000000}%
\pgfsetfillcolor{currentfill}%
\pgfsetlinewidth{0.501875pt}%
\definecolor{currentstroke}{rgb}{0.000000,0.000000,0.000000}%
\pgfsetstrokecolor{currentstroke}%
\pgfsetdash{}{0pt}%
\pgfsys@defobject{currentmarker}{\pgfqpoint{0.000000in}{-0.020833in}}{\pgfqpoint{0.000000in}{0.000000in}}{%
\pgfpathmoveto{\pgfqpoint{0.000000in}{0.000000in}}%
\pgfpathlineto{\pgfqpoint{0.000000in}{-0.020833in}}%
\pgfusepath{stroke,fill}%
}%
\begin{pgfscope}%
\pgfsys@transformshift{3.626388in}{3.899076in}%
\pgfsys@useobject{currentmarker}{}%
\end{pgfscope}%
\end{pgfscope}%
\begin{pgfscope}%
\pgfsetbuttcap%
\pgfsetroundjoin%
\definecolor{currentfill}{rgb}{0.000000,0.000000,0.000000}%
\pgfsetfillcolor{currentfill}%
\pgfsetlinewidth{0.501875pt}%
\definecolor{currentstroke}{rgb}{0.000000,0.000000,0.000000}%
\pgfsetstrokecolor{currentstroke}%
\pgfsetdash{}{0pt}%
\pgfsys@defobject{currentmarker}{\pgfqpoint{0.000000in}{0.000000in}}{\pgfqpoint{0.000000in}{0.020833in}}{%
\pgfpathmoveto{\pgfqpoint{0.000000in}{0.000000in}}%
\pgfpathlineto{\pgfqpoint{0.000000in}{0.020833in}}%
\pgfusepath{stroke,fill}%
}%
\begin{pgfscope}%
\pgfsys@transformshift{3.832844in}{0.405382in}%
\pgfsys@useobject{currentmarker}{}%
\end{pgfscope}%
\end{pgfscope}%
\begin{pgfscope}%
\pgfsetbuttcap%
\pgfsetroundjoin%
\definecolor{currentfill}{rgb}{0.000000,0.000000,0.000000}%
\pgfsetfillcolor{currentfill}%
\pgfsetlinewidth{0.501875pt}%
\definecolor{currentstroke}{rgb}{0.000000,0.000000,0.000000}%
\pgfsetstrokecolor{currentstroke}%
\pgfsetdash{}{0pt}%
\pgfsys@defobject{currentmarker}{\pgfqpoint{0.000000in}{-0.020833in}}{\pgfqpoint{0.000000in}{0.000000in}}{%
\pgfpathmoveto{\pgfqpoint{0.000000in}{0.000000in}}%
\pgfpathlineto{\pgfqpoint{0.000000in}{-0.020833in}}%
\pgfusepath{stroke,fill}%
}%
\begin{pgfscope}%
\pgfsys@transformshift{3.832844in}{3.899076in}%
\pgfsys@useobject{currentmarker}{}%
\end{pgfscope}%
\end{pgfscope}%
\begin{pgfscope}%
\pgfsetbuttcap%
\pgfsetroundjoin%
\definecolor{currentfill}{rgb}{0.000000,0.000000,0.000000}%
\pgfsetfillcolor{currentfill}%
\pgfsetlinewidth{0.501875pt}%
\definecolor{currentstroke}{rgb}{0.000000,0.000000,0.000000}%
\pgfsetstrokecolor{currentstroke}%
\pgfsetdash{}{0pt}%
\pgfsys@defobject{currentmarker}{\pgfqpoint{0.000000in}{0.000000in}}{\pgfqpoint{0.000000in}{0.020833in}}{%
\pgfpathmoveto{\pgfqpoint{0.000000in}{0.000000in}}%
\pgfpathlineto{\pgfqpoint{0.000000in}{0.020833in}}%
\pgfusepath{stroke,fill}%
}%
\begin{pgfscope}%
\pgfsys@transformshift{4.039299in}{0.405382in}%
\pgfsys@useobject{currentmarker}{}%
\end{pgfscope}%
\end{pgfscope}%
\begin{pgfscope}%
\pgfsetbuttcap%
\pgfsetroundjoin%
\definecolor{currentfill}{rgb}{0.000000,0.000000,0.000000}%
\pgfsetfillcolor{currentfill}%
\pgfsetlinewidth{0.501875pt}%
\definecolor{currentstroke}{rgb}{0.000000,0.000000,0.000000}%
\pgfsetstrokecolor{currentstroke}%
\pgfsetdash{}{0pt}%
\pgfsys@defobject{currentmarker}{\pgfqpoint{0.000000in}{-0.020833in}}{\pgfqpoint{0.000000in}{0.000000in}}{%
\pgfpathmoveto{\pgfqpoint{0.000000in}{0.000000in}}%
\pgfpathlineto{\pgfqpoint{0.000000in}{-0.020833in}}%
\pgfusepath{stroke,fill}%
}%
\begin{pgfscope}%
\pgfsys@transformshift{4.039299in}{3.899076in}%
\pgfsys@useobject{currentmarker}{}%
\end{pgfscope}%
\end{pgfscope}%
\begin{pgfscope}%
\pgfsetbuttcap%
\pgfsetroundjoin%
\definecolor{currentfill}{rgb}{0.000000,0.000000,0.000000}%
\pgfsetfillcolor{currentfill}%
\pgfsetlinewidth{0.501875pt}%
\definecolor{currentstroke}{rgb}{0.000000,0.000000,0.000000}%
\pgfsetstrokecolor{currentstroke}%
\pgfsetdash{}{0pt}%
\pgfsys@defobject{currentmarker}{\pgfqpoint{0.000000in}{0.000000in}}{\pgfqpoint{0.000000in}{0.020833in}}{%
\pgfpathmoveto{\pgfqpoint{0.000000in}{0.000000in}}%
\pgfpathlineto{\pgfqpoint{0.000000in}{0.020833in}}%
\pgfusepath{stroke,fill}%
}%
\begin{pgfscope}%
\pgfsys@transformshift{4.452209in}{0.405382in}%
\pgfsys@useobject{currentmarker}{}%
\end{pgfscope}%
\end{pgfscope}%
\begin{pgfscope}%
\pgfsetbuttcap%
\pgfsetroundjoin%
\definecolor{currentfill}{rgb}{0.000000,0.000000,0.000000}%
\pgfsetfillcolor{currentfill}%
\pgfsetlinewidth{0.501875pt}%
\definecolor{currentstroke}{rgb}{0.000000,0.000000,0.000000}%
\pgfsetstrokecolor{currentstroke}%
\pgfsetdash{}{0pt}%
\pgfsys@defobject{currentmarker}{\pgfqpoint{0.000000in}{-0.020833in}}{\pgfqpoint{0.000000in}{0.000000in}}{%
\pgfpathmoveto{\pgfqpoint{0.000000in}{0.000000in}}%
\pgfpathlineto{\pgfqpoint{0.000000in}{-0.020833in}}%
\pgfusepath{stroke,fill}%
}%
\begin{pgfscope}%
\pgfsys@transformshift{4.452209in}{3.899076in}%
\pgfsys@useobject{currentmarker}{}%
\end{pgfscope}%
\end{pgfscope}%
\begin{pgfscope}%
\pgfsetbuttcap%
\pgfsetroundjoin%
\definecolor{currentfill}{rgb}{0.000000,0.000000,0.000000}%
\pgfsetfillcolor{currentfill}%
\pgfsetlinewidth{0.501875pt}%
\definecolor{currentstroke}{rgb}{0.000000,0.000000,0.000000}%
\pgfsetstrokecolor{currentstroke}%
\pgfsetdash{}{0pt}%
\pgfsys@defobject{currentmarker}{\pgfqpoint{0.000000in}{0.000000in}}{\pgfqpoint{0.000000in}{0.020833in}}{%
\pgfpathmoveto{\pgfqpoint{0.000000in}{0.000000in}}%
\pgfpathlineto{\pgfqpoint{0.000000in}{0.020833in}}%
\pgfusepath{stroke,fill}%
}%
\begin{pgfscope}%
\pgfsys@transformshift{4.658664in}{0.405382in}%
\pgfsys@useobject{currentmarker}{}%
\end{pgfscope}%
\end{pgfscope}%
\begin{pgfscope}%
\pgfsetbuttcap%
\pgfsetroundjoin%
\definecolor{currentfill}{rgb}{0.000000,0.000000,0.000000}%
\pgfsetfillcolor{currentfill}%
\pgfsetlinewidth{0.501875pt}%
\definecolor{currentstroke}{rgb}{0.000000,0.000000,0.000000}%
\pgfsetstrokecolor{currentstroke}%
\pgfsetdash{}{0pt}%
\pgfsys@defobject{currentmarker}{\pgfqpoint{0.000000in}{-0.020833in}}{\pgfqpoint{0.000000in}{0.000000in}}{%
\pgfpathmoveto{\pgfqpoint{0.000000in}{0.000000in}}%
\pgfpathlineto{\pgfqpoint{0.000000in}{-0.020833in}}%
\pgfusepath{stroke,fill}%
}%
\begin{pgfscope}%
\pgfsys@transformshift{4.658664in}{3.899076in}%
\pgfsys@useobject{currentmarker}{}%
\end{pgfscope}%
\end{pgfscope}%
\begin{pgfscope}%
\pgfsetbuttcap%
\pgfsetroundjoin%
\definecolor{currentfill}{rgb}{0.000000,0.000000,0.000000}%
\pgfsetfillcolor{currentfill}%
\pgfsetlinewidth{0.501875pt}%
\definecolor{currentstroke}{rgb}{0.000000,0.000000,0.000000}%
\pgfsetstrokecolor{currentstroke}%
\pgfsetdash{}{0pt}%
\pgfsys@defobject{currentmarker}{\pgfqpoint{0.000000in}{0.000000in}}{\pgfqpoint{0.000000in}{0.020833in}}{%
\pgfpathmoveto{\pgfqpoint{0.000000in}{0.000000in}}%
\pgfpathlineto{\pgfqpoint{0.000000in}{0.020833in}}%
\pgfusepath{stroke,fill}%
}%
\begin{pgfscope}%
\pgfsys@transformshift{4.865119in}{0.405382in}%
\pgfsys@useobject{currentmarker}{}%
\end{pgfscope}%
\end{pgfscope}%
\begin{pgfscope}%
\pgfsetbuttcap%
\pgfsetroundjoin%
\definecolor{currentfill}{rgb}{0.000000,0.000000,0.000000}%
\pgfsetfillcolor{currentfill}%
\pgfsetlinewidth{0.501875pt}%
\definecolor{currentstroke}{rgb}{0.000000,0.000000,0.000000}%
\pgfsetstrokecolor{currentstroke}%
\pgfsetdash{}{0pt}%
\pgfsys@defobject{currentmarker}{\pgfqpoint{0.000000in}{-0.020833in}}{\pgfqpoint{0.000000in}{0.000000in}}{%
\pgfpathmoveto{\pgfqpoint{0.000000in}{0.000000in}}%
\pgfpathlineto{\pgfqpoint{0.000000in}{-0.020833in}}%
\pgfusepath{stroke,fill}%
}%
\begin{pgfscope}%
\pgfsys@transformshift{4.865119in}{3.899076in}%
\pgfsys@useobject{currentmarker}{}%
\end{pgfscope}%
\end{pgfscope}%
\begin{pgfscope}%
\pgfsetbuttcap%
\pgfsetroundjoin%
\definecolor{currentfill}{rgb}{0.000000,0.000000,0.000000}%
\pgfsetfillcolor{currentfill}%
\pgfsetlinewidth{0.501875pt}%
\definecolor{currentstroke}{rgb}{0.000000,0.000000,0.000000}%
\pgfsetstrokecolor{currentstroke}%
\pgfsetdash{}{0pt}%
\pgfsys@defobject{currentmarker}{\pgfqpoint{0.000000in}{0.000000in}}{\pgfqpoint{0.000000in}{0.020833in}}{%
\pgfpathmoveto{\pgfqpoint{0.000000in}{0.000000in}}%
\pgfpathlineto{\pgfqpoint{0.000000in}{0.020833in}}%
\pgfusepath{stroke,fill}%
}%
\begin{pgfscope}%
\pgfsys@transformshift{5.071574in}{0.405382in}%
\pgfsys@useobject{currentmarker}{}%
\end{pgfscope}%
\end{pgfscope}%
\begin{pgfscope}%
\pgfsetbuttcap%
\pgfsetroundjoin%
\definecolor{currentfill}{rgb}{0.000000,0.000000,0.000000}%
\pgfsetfillcolor{currentfill}%
\pgfsetlinewidth{0.501875pt}%
\definecolor{currentstroke}{rgb}{0.000000,0.000000,0.000000}%
\pgfsetstrokecolor{currentstroke}%
\pgfsetdash{}{0pt}%
\pgfsys@defobject{currentmarker}{\pgfqpoint{0.000000in}{-0.020833in}}{\pgfqpoint{0.000000in}{0.000000in}}{%
\pgfpathmoveto{\pgfqpoint{0.000000in}{0.000000in}}%
\pgfpathlineto{\pgfqpoint{0.000000in}{-0.020833in}}%
\pgfusepath{stroke,fill}%
}%
\begin{pgfscope}%
\pgfsys@transformshift{5.071574in}{3.899076in}%
\pgfsys@useobject{currentmarker}{}%
\end{pgfscope}%
\end{pgfscope}%
\begin{pgfscope}%
\pgfsetbuttcap%
\pgfsetroundjoin%
\definecolor{currentfill}{rgb}{0.000000,0.000000,0.000000}%
\pgfsetfillcolor{currentfill}%
\pgfsetlinewidth{0.501875pt}%
\definecolor{currentstroke}{rgb}{0.000000,0.000000,0.000000}%
\pgfsetstrokecolor{currentstroke}%
\pgfsetdash{}{0pt}%
\pgfsys@defobject{currentmarker}{\pgfqpoint{0.000000in}{0.000000in}}{\pgfqpoint{0.000000in}{0.020833in}}{%
\pgfpathmoveto{\pgfqpoint{0.000000in}{0.000000in}}%
\pgfpathlineto{\pgfqpoint{0.000000in}{0.020833in}}%
\pgfusepath{stroke,fill}%
}%
\begin{pgfscope}%
\pgfsys@transformshift{5.278029in}{0.405382in}%
\pgfsys@useobject{currentmarker}{}%
\end{pgfscope}%
\end{pgfscope}%
\begin{pgfscope}%
\pgfsetbuttcap%
\pgfsetroundjoin%
\definecolor{currentfill}{rgb}{0.000000,0.000000,0.000000}%
\pgfsetfillcolor{currentfill}%
\pgfsetlinewidth{0.501875pt}%
\definecolor{currentstroke}{rgb}{0.000000,0.000000,0.000000}%
\pgfsetstrokecolor{currentstroke}%
\pgfsetdash{}{0pt}%
\pgfsys@defobject{currentmarker}{\pgfqpoint{0.000000in}{-0.020833in}}{\pgfqpoint{0.000000in}{0.000000in}}{%
\pgfpathmoveto{\pgfqpoint{0.000000in}{0.000000in}}%
\pgfpathlineto{\pgfqpoint{0.000000in}{-0.020833in}}%
\pgfusepath{stroke,fill}%
}%
\begin{pgfscope}%
\pgfsys@transformshift{5.278029in}{3.899076in}%
\pgfsys@useobject{currentmarker}{}%
\end{pgfscope}%
\end{pgfscope}%
\begin{pgfscope}%
\pgfsetbuttcap%
\pgfsetroundjoin%
\definecolor{currentfill}{rgb}{0.000000,0.000000,0.000000}%
\pgfsetfillcolor{currentfill}%
\pgfsetlinewidth{0.501875pt}%
\definecolor{currentstroke}{rgb}{0.000000,0.000000,0.000000}%
\pgfsetstrokecolor{currentstroke}%
\pgfsetdash{}{0pt}%
\pgfsys@defobject{currentmarker}{\pgfqpoint{0.000000in}{0.000000in}}{\pgfqpoint{0.000000in}{0.020833in}}{%
\pgfpathmoveto{\pgfqpoint{0.000000in}{0.000000in}}%
\pgfpathlineto{\pgfqpoint{0.000000in}{0.020833in}}%
\pgfusepath{stroke,fill}%
}%
\begin{pgfscope}%
\pgfsys@transformshift{5.690940in}{0.405382in}%
\pgfsys@useobject{currentmarker}{}%
\end{pgfscope}%
\end{pgfscope}%
\begin{pgfscope}%
\pgfsetbuttcap%
\pgfsetroundjoin%
\definecolor{currentfill}{rgb}{0.000000,0.000000,0.000000}%
\pgfsetfillcolor{currentfill}%
\pgfsetlinewidth{0.501875pt}%
\definecolor{currentstroke}{rgb}{0.000000,0.000000,0.000000}%
\pgfsetstrokecolor{currentstroke}%
\pgfsetdash{}{0pt}%
\pgfsys@defobject{currentmarker}{\pgfqpoint{0.000000in}{-0.020833in}}{\pgfqpoint{0.000000in}{0.000000in}}{%
\pgfpathmoveto{\pgfqpoint{0.000000in}{0.000000in}}%
\pgfpathlineto{\pgfqpoint{0.000000in}{-0.020833in}}%
\pgfusepath{stroke,fill}%
}%
\begin{pgfscope}%
\pgfsys@transformshift{5.690940in}{3.899076in}%
\pgfsys@useobject{currentmarker}{}%
\end{pgfscope}%
\end{pgfscope}%
\begin{pgfscope}%
\pgfsetbuttcap%
\pgfsetroundjoin%
\definecolor{currentfill}{rgb}{0.000000,0.000000,0.000000}%
\pgfsetfillcolor{currentfill}%
\pgfsetlinewidth{0.501875pt}%
\definecolor{currentstroke}{rgb}{0.000000,0.000000,0.000000}%
\pgfsetstrokecolor{currentstroke}%
\pgfsetdash{}{0pt}%
\pgfsys@defobject{currentmarker}{\pgfqpoint{0.000000in}{0.000000in}}{\pgfqpoint{0.000000in}{0.020833in}}{%
\pgfpathmoveto{\pgfqpoint{0.000000in}{0.000000in}}%
\pgfpathlineto{\pgfqpoint{0.000000in}{0.020833in}}%
\pgfusepath{stroke,fill}%
}%
\begin{pgfscope}%
\pgfsys@transformshift{5.897395in}{0.405382in}%
\pgfsys@useobject{currentmarker}{}%
\end{pgfscope}%
\end{pgfscope}%
\begin{pgfscope}%
\pgfsetbuttcap%
\pgfsetroundjoin%
\definecolor{currentfill}{rgb}{0.000000,0.000000,0.000000}%
\pgfsetfillcolor{currentfill}%
\pgfsetlinewidth{0.501875pt}%
\definecolor{currentstroke}{rgb}{0.000000,0.000000,0.000000}%
\pgfsetstrokecolor{currentstroke}%
\pgfsetdash{}{0pt}%
\pgfsys@defobject{currentmarker}{\pgfqpoint{0.000000in}{-0.020833in}}{\pgfqpoint{0.000000in}{0.000000in}}{%
\pgfpathmoveto{\pgfqpoint{0.000000in}{0.000000in}}%
\pgfpathlineto{\pgfqpoint{0.000000in}{-0.020833in}}%
\pgfusepath{stroke,fill}%
}%
\begin{pgfscope}%
\pgfsys@transformshift{5.897395in}{3.899076in}%
\pgfsys@useobject{currentmarker}{}%
\end{pgfscope}%
\end{pgfscope}%
\begin{pgfscope}%
\pgfsetbuttcap%
\pgfsetroundjoin%
\definecolor{currentfill}{rgb}{0.000000,0.000000,0.000000}%
\pgfsetfillcolor{currentfill}%
\pgfsetlinewidth{0.501875pt}%
\definecolor{currentstroke}{rgb}{0.000000,0.000000,0.000000}%
\pgfsetstrokecolor{currentstroke}%
\pgfsetdash{}{0pt}%
\pgfsys@defobject{currentmarker}{\pgfqpoint{0.000000in}{0.000000in}}{\pgfqpoint{0.000000in}{0.020833in}}{%
\pgfpathmoveto{\pgfqpoint{0.000000in}{0.000000in}}%
\pgfpathlineto{\pgfqpoint{0.000000in}{0.020833in}}%
\pgfusepath{stroke,fill}%
}%
\begin{pgfscope}%
\pgfsys@transformshift{6.103850in}{0.405382in}%
\pgfsys@useobject{currentmarker}{}%
\end{pgfscope}%
\end{pgfscope}%
\begin{pgfscope}%
\pgfsetbuttcap%
\pgfsetroundjoin%
\definecolor{currentfill}{rgb}{0.000000,0.000000,0.000000}%
\pgfsetfillcolor{currentfill}%
\pgfsetlinewidth{0.501875pt}%
\definecolor{currentstroke}{rgb}{0.000000,0.000000,0.000000}%
\pgfsetstrokecolor{currentstroke}%
\pgfsetdash{}{0pt}%
\pgfsys@defobject{currentmarker}{\pgfqpoint{0.000000in}{-0.020833in}}{\pgfqpoint{0.000000in}{0.000000in}}{%
\pgfpathmoveto{\pgfqpoint{0.000000in}{0.000000in}}%
\pgfpathlineto{\pgfqpoint{0.000000in}{-0.020833in}}%
\pgfusepath{stroke,fill}%
}%
\begin{pgfscope}%
\pgfsys@transformshift{6.103850in}{3.899076in}%
\pgfsys@useobject{currentmarker}{}%
\end{pgfscope}%
\end{pgfscope}%
\begin{pgfscope}%
\pgfsetbuttcap%
\pgfsetroundjoin%
\definecolor{currentfill}{rgb}{0.000000,0.000000,0.000000}%
\pgfsetfillcolor{currentfill}%
\pgfsetlinewidth{1.003750pt}%
\definecolor{currentstroke}{rgb}{0.000000,0.000000,0.000000}%
\pgfsetstrokecolor{currentstroke}%
\pgfsetdash{}{0pt}%
\pgfsys@defobject{currentmarker}{\pgfqpoint{-0.027778in}{-0.000000in}}{\pgfqpoint{0.027778in}{0.000000in}}{%
\pgfpathmoveto{\pgfqpoint{0.027778in}{-0.000000in}}%
\pgfpathlineto{\pgfqpoint{-0.027778in}{0.000000in}}%
\pgfusepath{stroke,fill}%
}%
\begin{pgfscope}%
\pgfsys@transformshift{1.355382in}{0.405382in}%
\pgfsys@useobject{currentmarker}{}%
\end{pgfscope}%
\end{pgfscope}%
\begin{pgfscope}%
\pgfsetbuttcap%
\pgfsetroundjoin%
\definecolor{currentfill}{rgb}{0.000000,0.000000,0.000000}%
\pgfsetfillcolor{currentfill}%
\pgfsetlinewidth{1.003750pt}%
\definecolor{currentstroke}{rgb}{0.000000,0.000000,0.000000}%
\pgfsetstrokecolor{currentstroke}%
\pgfsetdash{}{0pt}%
\pgfsys@defobject{currentmarker}{\pgfqpoint{-0.027778in}{-0.000000in}}{\pgfqpoint{0.027778in}{0.000000in}}{%
\pgfpathmoveto{\pgfqpoint{0.027778in}{-0.000000in}}%
\pgfpathlineto{\pgfqpoint{-0.027778in}{0.000000in}}%
\pgfusepath{stroke,fill}%
}%
\begin{pgfscope}%
\pgfsys@transformshift{6.103850in}{0.405382in}%
\pgfsys@useobject{currentmarker}{}%
\end{pgfscope}%
\end{pgfscope}%
\begin{pgfscope}%
\definecolor{textcolor}{rgb}{0.000000,0.000000,0.000000}%
\pgfsetstrokecolor{textcolor}%
\pgfsetfillcolor{textcolor}%
\pgftext[x=1.278993in,y=0.405382in,right,]{\color{textcolor}{\rmfamily\fontsize{10.000000}{12.000000}\selectfont\catcode`\^=\active\def^{\ifmmode\sp\else\^{}\fi}\catcode`\%=\active\def%{\%}$h(x_1)$}}%
\end{pgfscope}%
\begin{pgfscope}%
\pgfsetbuttcap%
\pgfsetroundjoin%
\definecolor{currentfill}{rgb}{0.000000,0.000000,0.000000}%
\pgfsetfillcolor{currentfill}%
\pgfsetlinewidth{1.003750pt}%
\definecolor{currentstroke}{rgb}{0.000000,0.000000,0.000000}%
\pgfsetstrokecolor{currentstroke}%
\pgfsetdash{}{0pt}%
\pgfsys@defobject{currentmarker}{\pgfqpoint{-0.027778in}{-0.000000in}}{\pgfqpoint{0.027778in}{0.000000in}}{%
\pgfpathmoveto{\pgfqpoint{0.027778in}{-0.000000in}}%
\pgfpathlineto{\pgfqpoint{-0.027778in}{0.000000in}}%
\pgfusepath{stroke,fill}%
}%
\begin{pgfscope}%
\pgfsys@transformshift{1.355382in}{1.403580in}%
\pgfsys@useobject{currentmarker}{}%
\end{pgfscope}%
\end{pgfscope}%
\begin{pgfscope}%
\pgfsetbuttcap%
\pgfsetroundjoin%
\definecolor{currentfill}{rgb}{0.000000,0.000000,0.000000}%
\pgfsetfillcolor{currentfill}%
\pgfsetlinewidth{1.003750pt}%
\definecolor{currentstroke}{rgb}{0.000000,0.000000,0.000000}%
\pgfsetstrokecolor{currentstroke}%
\pgfsetdash{}{0pt}%
\pgfsys@defobject{currentmarker}{\pgfqpoint{-0.027778in}{-0.000000in}}{\pgfqpoint{0.027778in}{0.000000in}}{%
\pgfpathmoveto{\pgfqpoint{0.027778in}{-0.000000in}}%
\pgfpathlineto{\pgfqpoint{-0.027778in}{0.000000in}}%
\pgfusepath{stroke,fill}%
}%
\begin{pgfscope}%
\pgfsys@transformshift{6.103850in}{1.403580in}%
\pgfsys@useobject{currentmarker}{}%
\end{pgfscope}%
\end{pgfscope}%
\begin{pgfscope}%
\definecolor{textcolor}{rgb}{0.000000,0.000000,0.000000}%
\pgfsetstrokecolor{textcolor}%
\pgfsetfillcolor{textcolor}%
\pgftext[x=1.278993in,y=1.403580in,right,]{\color{textcolor}{\rmfamily\fontsize{10.000000}{12.000000}\selectfont\catcode`\^=\active\def^{\ifmmode\sp\else\^{}\fi}\catcode`\%=\active\def%{\%}$h((1-\lambda)x_1 + \lambda x_2)$}}%
\end{pgfscope}%
\begin{pgfscope}%
\pgfsetbuttcap%
\pgfsetroundjoin%
\definecolor{currentfill}{rgb}{0.000000,0.000000,0.000000}%
\pgfsetfillcolor{currentfill}%
\pgfsetlinewidth{1.003750pt}%
\definecolor{currentstroke}{rgb}{0.000000,0.000000,0.000000}%
\pgfsetstrokecolor{currentstroke}%
\pgfsetdash{}{0pt}%
\pgfsys@defobject{currentmarker}{\pgfqpoint{-0.027778in}{-0.000000in}}{\pgfqpoint{0.027778in}{0.000000in}}{%
\pgfpathmoveto{\pgfqpoint{0.027778in}{-0.000000in}}%
\pgfpathlineto{\pgfqpoint{-0.027778in}{0.000000in}}%
\pgfusepath{stroke,fill}%
}%
\begin{pgfscope}%
\pgfsys@transformshift{1.355382in}{1.902679in}%
\pgfsys@useobject{currentmarker}{}%
\end{pgfscope}%
\end{pgfscope}%
\begin{pgfscope}%
\pgfsetbuttcap%
\pgfsetroundjoin%
\definecolor{currentfill}{rgb}{0.000000,0.000000,0.000000}%
\pgfsetfillcolor{currentfill}%
\pgfsetlinewidth{1.003750pt}%
\definecolor{currentstroke}{rgb}{0.000000,0.000000,0.000000}%
\pgfsetstrokecolor{currentstroke}%
\pgfsetdash{}{0pt}%
\pgfsys@defobject{currentmarker}{\pgfqpoint{-0.027778in}{-0.000000in}}{\pgfqpoint{0.027778in}{0.000000in}}{%
\pgfpathmoveto{\pgfqpoint{0.027778in}{-0.000000in}}%
\pgfpathlineto{\pgfqpoint{-0.027778in}{0.000000in}}%
\pgfusepath{stroke,fill}%
}%
\begin{pgfscope}%
\pgfsys@transformshift{6.103850in}{1.902679in}%
\pgfsys@useobject{currentmarker}{}%
\end{pgfscope}%
\end{pgfscope}%
\begin{pgfscope}%
\definecolor{textcolor}{rgb}{0.000000,0.000000,0.000000}%
\pgfsetstrokecolor{textcolor}%
\pgfsetfillcolor{textcolor}%
\pgftext[x=1.278993in,y=1.902679in,right,]{\color{textcolor}{\rmfamily\fontsize{10.000000}{12.000000}\selectfont\catcode`\^=\active\def^{\ifmmode\sp\else\^{}\fi}\catcode`\%=\active\def%{\%}$(1-\lambda)h(x_1) + \lambda h(x_2)$}}%
\end{pgfscope}%
\begin{pgfscope}%
\pgfsetbuttcap%
\pgfsetroundjoin%
\definecolor{currentfill}{rgb}{0.000000,0.000000,0.000000}%
\pgfsetfillcolor{currentfill}%
\pgfsetlinewidth{1.003750pt}%
\definecolor{currentstroke}{rgb}{0.000000,0.000000,0.000000}%
\pgfsetstrokecolor{currentstroke}%
\pgfsetdash{}{0pt}%
\pgfsys@defobject{currentmarker}{\pgfqpoint{-0.027778in}{-0.000000in}}{\pgfqpoint{0.027778in}{0.000000in}}{%
\pgfpathmoveto{\pgfqpoint{0.027778in}{-0.000000in}}%
\pgfpathlineto{\pgfqpoint{-0.027778in}{0.000000in}}%
\pgfusepath{stroke,fill}%
}%
\begin{pgfscope}%
\pgfsys@transformshift{1.355382in}{2.900878in}%
\pgfsys@useobject{currentmarker}{}%
\end{pgfscope}%
\end{pgfscope}%
\begin{pgfscope}%
\pgfsetbuttcap%
\pgfsetroundjoin%
\definecolor{currentfill}{rgb}{0.000000,0.000000,0.000000}%
\pgfsetfillcolor{currentfill}%
\pgfsetlinewidth{1.003750pt}%
\definecolor{currentstroke}{rgb}{0.000000,0.000000,0.000000}%
\pgfsetstrokecolor{currentstroke}%
\pgfsetdash{}{0pt}%
\pgfsys@defobject{currentmarker}{\pgfqpoint{-0.027778in}{-0.000000in}}{\pgfqpoint{0.027778in}{0.000000in}}{%
\pgfpathmoveto{\pgfqpoint{0.027778in}{-0.000000in}}%
\pgfpathlineto{\pgfqpoint{-0.027778in}{0.000000in}}%
\pgfusepath{stroke,fill}%
}%
\begin{pgfscope}%
\pgfsys@transformshift{6.103850in}{2.900878in}%
\pgfsys@useobject{currentmarker}{}%
\end{pgfscope}%
\end{pgfscope}%
\begin{pgfscope}%
\definecolor{textcolor}{rgb}{0.000000,0.000000,0.000000}%
\pgfsetstrokecolor{textcolor}%
\pgfsetfillcolor{textcolor}%
\pgftext[x=1.278993in,y=2.900878in,right,]{\color{textcolor}{\rmfamily\fontsize{10.000000}{12.000000}\selectfont\catcode`\^=\active\def^{\ifmmode\sp\else\^{}\fi}\catcode`\%=\active\def%{\%}$h(x_2)$}}%
\end{pgfscope}%
\begin{pgfscope}%
\pgfsetbuttcap%
\pgfsetroundjoin%
\definecolor{currentfill}{rgb}{0.000000,0.000000,0.000000}%
\pgfsetfillcolor{currentfill}%
\pgfsetlinewidth{0.501875pt}%
\definecolor{currentstroke}{rgb}{0.000000,0.000000,0.000000}%
\pgfsetstrokecolor{currentstroke}%
\pgfsetdash{}{0pt}%
\pgfsys@defobject{currentmarker}{\pgfqpoint{0.000000in}{0.000000in}}{\pgfqpoint{0.020833in}{0.000000in}}{%
\pgfpathmoveto{\pgfqpoint{0.000000in}{0.000000in}}%
\pgfpathlineto{\pgfqpoint{0.020833in}{0.000000in}}%
\pgfusepath{stroke,fill}%
}%
\begin{pgfscope}%
\pgfsys@transformshift{1.355382in}{0.155832in}%
\pgfsys@useobject{currentmarker}{}%
\end{pgfscope}%
\end{pgfscope}%
\begin{pgfscope}%
\pgfsetbuttcap%
\pgfsetroundjoin%
\definecolor{currentfill}{rgb}{0.000000,0.000000,0.000000}%
\pgfsetfillcolor{currentfill}%
\pgfsetlinewidth{0.501875pt}%
\definecolor{currentstroke}{rgb}{0.000000,0.000000,0.000000}%
\pgfsetstrokecolor{currentstroke}%
\pgfsetdash{}{0pt}%
\pgfsys@defobject{currentmarker}{\pgfqpoint{-0.020833in}{0.000000in}}{\pgfqpoint{-0.000000in}{0.000000in}}{%
\pgfpathmoveto{\pgfqpoint{-0.000000in}{0.000000in}}%
\pgfpathlineto{\pgfqpoint{-0.020833in}{0.000000in}}%
\pgfusepath{stroke,fill}%
}%
\begin{pgfscope}%
\pgfsys@transformshift{6.103850in}{0.155832in}%
\pgfsys@useobject{currentmarker}{}%
\end{pgfscope}%
\end{pgfscope}%
\begin{pgfscope}%
\pgfsetbuttcap%
\pgfsetroundjoin%
\definecolor{currentfill}{rgb}{0.000000,0.000000,0.000000}%
\pgfsetfillcolor{currentfill}%
\pgfsetlinewidth{0.501875pt}%
\definecolor{currentstroke}{rgb}{0.000000,0.000000,0.000000}%
\pgfsetstrokecolor{currentstroke}%
\pgfsetdash{}{0pt}%
\pgfsys@defobject{currentmarker}{\pgfqpoint{0.000000in}{0.000000in}}{\pgfqpoint{0.020833in}{0.000000in}}{%
\pgfpathmoveto{\pgfqpoint{0.000000in}{0.000000in}}%
\pgfpathlineto{\pgfqpoint{0.020833in}{0.000000in}}%
\pgfusepath{stroke,fill}%
}%
\begin{pgfscope}%
\pgfsys@transformshift{1.355382in}{0.654931in}%
\pgfsys@useobject{currentmarker}{}%
\end{pgfscope}%
\end{pgfscope}%
\begin{pgfscope}%
\pgfsetbuttcap%
\pgfsetroundjoin%
\definecolor{currentfill}{rgb}{0.000000,0.000000,0.000000}%
\pgfsetfillcolor{currentfill}%
\pgfsetlinewidth{0.501875pt}%
\definecolor{currentstroke}{rgb}{0.000000,0.000000,0.000000}%
\pgfsetstrokecolor{currentstroke}%
\pgfsetdash{}{0pt}%
\pgfsys@defobject{currentmarker}{\pgfqpoint{-0.020833in}{0.000000in}}{\pgfqpoint{-0.000000in}{0.000000in}}{%
\pgfpathmoveto{\pgfqpoint{-0.000000in}{0.000000in}}%
\pgfpathlineto{\pgfqpoint{-0.020833in}{0.000000in}}%
\pgfusepath{stroke,fill}%
}%
\begin{pgfscope}%
\pgfsys@transformshift{6.103850in}{0.654931in}%
\pgfsys@useobject{currentmarker}{}%
\end{pgfscope}%
\end{pgfscope}%
\begin{pgfscope}%
\pgfsetbuttcap%
\pgfsetroundjoin%
\definecolor{currentfill}{rgb}{0.000000,0.000000,0.000000}%
\pgfsetfillcolor{currentfill}%
\pgfsetlinewidth{0.501875pt}%
\definecolor{currentstroke}{rgb}{0.000000,0.000000,0.000000}%
\pgfsetstrokecolor{currentstroke}%
\pgfsetdash{}{0pt}%
\pgfsys@defobject{currentmarker}{\pgfqpoint{0.000000in}{0.000000in}}{\pgfqpoint{0.020833in}{0.000000in}}{%
\pgfpathmoveto{\pgfqpoint{0.000000in}{0.000000in}}%
\pgfpathlineto{\pgfqpoint{0.020833in}{0.000000in}}%
\pgfusepath{stroke,fill}%
}%
\begin{pgfscope}%
\pgfsys@transformshift{1.355382in}{0.904481in}%
\pgfsys@useobject{currentmarker}{}%
\end{pgfscope}%
\end{pgfscope}%
\begin{pgfscope}%
\pgfsetbuttcap%
\pgfsetroundjoin%
\definecolor{currentfill}{rgb}{0.000000,0.000000,0.000000}%
\pgfsetfillcolor{currentfill}%
\pgfsetlinewidth{0.501875pt}%
\definecolor{currentstroke}{rgb}{0.000000,0.000000,0.000000}%
\pgfsetstrokecolor{currentstroke}%
\pgfsetdash{}{0pt}%
\pgfsys@defobject{currentmarker}{\pgfqpoint{-0.020833in}{0.000000in}}{\pgfqpoint{-0.000000in}{0.000000in}}{%
\pgfpathmoveto{\pgfqpoint{-0.000000in}{0.000000in}}%
\pgfpathlineto{\pgfqpoint{-0.020833in}{0.000000in}}%
\pgfusepath{stroke,fill}%
}%
\begin{pgfscope}%
\pgfsys@transformshift{6.103850in}{0.904481in}%
\pgfsys@useobject{currentmarker}{}%
\end{pgfscope}%
\end{pgfscope}%
\begin{pgfscope}%
\pgfsetbuttcap%
\pgfsetroundjoin%
\definecolor{currentfill}{rgb}{0.000000,0.000000,0.000000}%
\pgfsetfillcolor{currentfill}%
\pgfsetlinewidth{0.501875pt}%
\definecolor{currentstroke}{rgb}{0.000000,0.000000,0.000000}%
\pgfsetstrokecolor{currentstroke}%
\pgfsetdash{}{0pt}%
\pgfsys@defobject{currentmarker}{\pgfqpoint{0.000000in}{0.000000in}}{\pgfqpoint{0.020833in}{0.000000in}}{%
\pgfpathmoveto{\pgfqpoint{0.000000in}{0.000000in}}%
\pgfpathlineto{\pgfqpoint{0.020833in}{0.000000in}}%
\pgfusepath{stroke,fill}%
}%
\begin{pgfscope}%
\pgfsys@transformshift{1.355382in}{1.154031in}%
\pgfsys@useobject{currentmarker}{}%
\end{pgfscope}%
\end{pgfscope}%
\begin{pgfscope}%
\pgfsetbuttcap%
\pgfsetroundjoin%
\definecolor{currentfill}{rgb}{0.000000,0.000000,0.000000}%
\pgfsetfillcolor{currentfill}%
\pgfsetlinewidth{0.501875pt}%
\definecolor{currentstroke}{rgb}{0.000000,0.000000,0.000000}%
\pgfsetstrokecolor{currentstroke}%
\pgfsetdash{}{0pt}%
\pgfsys@defobject{currentmarker}{\pgfqpoint{-0.020833in}{0.000000in}}{\pgfqpoint{-0.000000in}{0.000000in}}{%
\pgfpathmoveto{\pgfqpoint{-0.000000in}{0.000000in}}%
\pgfpathlineto{\pgfqpoint{-0.020833in}{0.000000in}}%
\pgfusepath{stroke,fill}%
}%
\begin{pgfscope}%
\pgfsys@transformshift{6.103850in}{1.154031in}%
\pgfsys@useobject{currentmarker}{}%
\end{pgfscope}%
\end{pgfscope}%
\begin{pgfscope}%
\pgfsetbuttcap%
\pgfsetroundjoin%
\definecolor{currentfill}{rgb}{0.000000,0.000000,0.000000}%
\pgfsetfillcolor{currentfill}%
\pgfsetlinewidth{0.501875pt}%
\definecolor{currentstroke}{rgb}{0.000000,0.000000,0.000000}%
\pgfsetstrokecolor{currentstroke}%
\pgfsetdash{}{0pt}%
\pgfsys@defobject{currentmarker}{\pgfqpoint{0.000000in}{0.000000in}}{\pgfqpoint{0.020833in}{0.000000in}}{%
\pgfpathmoveto{\pgfqpoint{0.000000in}{0.000000in}}%
\pgfpathlineto{\pgfqpoint{0.020833in}{0.000000in}}%
\pgfusepath{stroke,fill}%
}%
\begin{pgfscope}%
\pgfsys@transformshift{1.355382in}{1.653130in}%
\pgfsys@useobject{currentmarker}{}%
\end{pgfscope}%
\end{pgfscope}%
\begin{pgfscope}%
\pgfsetbuttcap%
\pgfsetroundjoin%
\definecolor{currentfill}{rgb}{0.000000,0.000000,0.000000}%
\pgfsetfillcolor{currentfill}%
\pgfsetlinewidth{0.501875pt}%
\definecolor{currentstroke}{rgb}{0.000000,0.000000,0.000000}%
\pgfsetstrokecolor{currentstroke}%
\pgfsetdash{}{0pt}%
\pgfsys@defobject{currentmarker}{\pgfqpoint{-0.020833in}{0.000000in}}{\pgfqpoint{-0.000000in}{0.000000in}}{%
\pgfpathmoveto{\pgfqpoint{-0.000000in}{0.000000in}}%
\pgfpathlineto{\pgfqpoint{-0.020833in}{0.000000in}}%
\pgfusepath{stroke,fill}%
}%
\begin{pgfscope}%
\pgfsys@transformshift{6.103850in}{1.653130in}%
\pgfsys@useobject{currentmarker}{}%
\end{pgfscope}%
\end{pgfscope}%
\begin{pgfscope}%
\pgfsetbuttcap%
\pgfsetroundjoin%
\definecolor{currentfill}{rgb}{0.000000,0.000000,0.000000}%
\pgfsetfillcolor{currentfill}%
\pgfsetlinewidth{0.501875pt}%
\definecolor{currentstroke}{rgb}{0.000000,0.000000,0.000000}%
\pgfsetstrokecolor{currentstroke}%
\pgfsetdash{}{0pt}%
\pgfsys@defobject{currentmarker}{\pgfqpoint{0.000000in}{0.000000in}}{\pgfqpoint{0.020833in}{0.000000in}}{%
\pgfpathmoveto{\pgfqpoint{0.000000in}{0.000000in}}%
\pgfpathlineto{\pgfqpoint{0.020833in}{0.000000in}}%
\pgfusepath{stroke,fill}%
}%
\begin{pgfscope}%
\pgfsys@transformshift{1.355382in}{2.152229in}%
\pgfsys@useobject{currentmarker}{}%
\end{pgfscope}%
\end{pgfscope}%
\begin{pgfscope}%
\pgfsetbuttcap%
\pgfsetroundjoin%
\definecolor{currentfill}{rgb}{0.000000,0.000000,0.000000}%
\pgfsetfillcolor{currentfill}%
\pgfsetlinewidth{0.501875pt}%
\definecolor{currentstroke}{rgb}{0.000000,0.000000,0.000000}%
\pgfsetstrokecolor{currentstroke}%
\pgfsetdash{}{0pt}%
\pgfsys@defobject{currentmarker}{\pgfqpoint{-0.020833in}{0.000000in}}{\pgfqpoint{-0.000000in}{0.000000in}}{%
\pgfpathmoveto{\pgfqpoint{-0.000000in}{0.000000in}}%
\pgfpathlineto{\pgfqpoint{-0.020833in}{0.000000in}}%
\pgfusepath{stroke,fill}%
}%
\begin{pgfscope}%
\pgfsys@transformshift{6.103850in}{2.152229in}%
\pgfsys@useobject{currentmarker}{}%
\end{pgfscope}%
\end{pgfscope}%
\begin{pgfscope}%
\pgfsetbuttcap%
\pgfsetroundjoin%
\definecolor{currentfill}{rgb}{0.000000,0.000000,0.000000}%
\pgfsetfillcolor{currentfill}%
\pgfsetlinewidth{0.501875pt}%
\definecolor{currentstroke}{rgb}{0.000000,0.000000,0.000000}%
\pgfsetstrokecolor{currentstroke}%
\pgfsetdash{}{0pt}%
\pgfsys@defobject{currentmarker}{\pgfqpoint{0.000000in}{0.000000in}}{\pgfqpoint{0.020833in}{0.000000in}}{%
\pgfpathmoveto{\pgfqpoint{0.000000in}{0.000000in}}%
\pgfpathlineto{\pgfqpoint{0.020833in}{0.000000in}}%
\pgfusepath{stroke,fill}%
}%
\begin{pgfscope}%
\pgfsys@transformshift{1.355382in}{2.401779in}%
\pgfsys@useobject{currentmarker}{}%
\end{pgfscope}%
\end{pgfscope}%
\begin{pgfscope}%
\pgfsetbuttcap%
\pgfsetroundjoin%
\definecolor{currentfill}{rgb}{0.000000,0.000000,0.000000}%
\pgfsetfillcolor{currentfill}%
\pgfsetlinewidth{0.501875pt}%
\definecolor{currentstroke}{rgb}{0.000000,0.000000,0.000000}%
\pgfsetstrokecolor{currentstroke}%
\pgfsetdash{}{0pt}%
\pgfsys@defobject{currentmarker}{\pgfqpoint{-0.020833in}{0.000000in}}{\pgfqpoint{-0.000000in}{0.000000in}}{%
\pgfpathmoveto{\pgfqpoint{-0.000000in}{0.000000in}}%
\pgfpathlineto{\pgfqpoint{-0.020833in}{0.000000in}}%
\pgfusepath{stroke,fill}%
}%
\begin{pgfscope}%
\pgfsys@transformshift{6.103850in}{2.401779in}%
\pgfsys@useobject{currentmarker}{}%
\end{pgfscope}%
\end{pgfscope}%
\begin{pgfscope}%
\pgfsetbuttcap%
\pgfsetroundjoin%
\definecolor{currentfill}{rgb}{0.000000,0.000000,0.000000}%
\pgfsetfillcolor{currentfill}%
\pgfsetlinewidth{0.501875pt}%
\definecolor{currentstroke}{rgb}{0.000000,0.000000,0.000000}%
\pgfsetstrokecolor{currentstroke}%
\pgfsetdash{}{0pt}%
\pgfsys@defobject{currentmarker}{\pgfqpoint{0.000000in}{0.000000in}}{\pgfqpoint{0.020833in}{0.000000in}}{%
\pgfpathmoveto{\pgfqpoint{0.000000in}{0.000000in}}%
\pgfpathlineto{\pgfqpoint{0.020833in}{0.000000in}}%
\pgfusepath{stroke,fill}%
}%
\begin{pgfscope}%
\pgfsys@transformshift{1.355382in}{2.651328in}%
\pgfsys@useobject{currentmarker}{}%
\end{pgfscope}%
\end{pgfscope}%
\begin{pgfscope}%
\pgfsetbuttcap%
\pgfsetroundjoin%
\definecolor{currentfill}{rgb}{0.000000,0.000000,0.000000}%
\pgfsetfillcolor{currentfill}%
\pgfsetlinewidth{0.501875pt}%
\definecolor{currentstroke}{rgb}{0.000000,0.000000,0.000000}%
\pgfsetstrokecolor{currentstroke}%
\pgfsetdash{}{0pt}%
\pgfsys@defobject{currentmarker}{\pgfqpoint{-0.020833in}{0.000000in}}{\pgfqpoint{-0.000000in}{0.000000in}}{%
\pgfpathmoveto{\pgfqpoint{-0.000000in}{0.000000in}}%
\pgfpathlineto{\pgfqpoint{-0.020833in}{0.000000in}}%
\pgfusepath{stroke,fill}%
}%
\begin{pgfscope}%
\pgfsys@transformshift{6.103850in}{2.651328in}%
\pgfsys@useobject{currentmarker}{}%
\end{pgfscope}%
\end{pgfscope}%
\begin{pgfscope}%
\pgfsetbuttcap%
\pgfsetroundjoin%
\definecolor{currentfill}{rgb}{0.000000,0.000000,0.000000}%
\pgfsetfillcolor{currentfill}%
\pgfsetlinewidth{0.501875pt}%
\definecolor{currentstroke}{rgb}{0.000000,0.000000,0.000000}%
\pgfsetstrokecolor{currentstroke}%
\pgfsetdash{}{0pt}%
\pgfsys@defobject{currentmarker}{\pgfqpoint{0.000000in}{0.000000in}}{\pgfqpoint{0.020833in}{0.000000in}}{%
\pgfpathmoveto{\pgfqpoint{0.000000in}{0.000000in}}%
\pgfpathlineto{\pgfqpoint{0.020833in}{0.000000in}}%
\pgfusepath{stroke,fill}%
}%
\begin{pgfscope}%
\pgfsys@transformshift{1.355382in}{3.150427in}%
\pgfsys@useobject{currentmarker}{}%
\end{pgfscope}%
\end{pgfscope}%
\begin{pgfscope}%
\pgfsetbuttcap%
\pgfsetroundjoin%
\definecolor{currentfill}{rgb}{0.000000,0.000000,0.000000}%
\pgfsetfillcolor{currentfill}%
\pgfsetlinewidth{0.501875pt}%
\definecolor{currentstroke}{rgb}{0.000000,0.000000,0.000000}%
\pgfsetstrokecolor{currentstroke}%
\pgfsetdash{}{0pt}%
\pgfsys@defobject{currentmarker}{\pgfqpoint{-0.020833in}{0.000000in}}{\pgfqpoint{-0.000000in}{0.000000in}}{%
\pgfpathmoveto{\pgfqpoint{-0.000000in}{0.000000in}}%
\pgfpathlineto{\pgfqpoint{-0.020833in}{0.000000in}}%
\pgfusepath{stroke,fill}%
}%
\begin{pgfscope}%
\pgfsys@transformshift{6.103850in}{3.150427in}%
\pgfsys@useobject{currentmarker}{}%
\end{pgfscope}%
\end{pgfscope}%
\begin{pgfscope}%
\pgfsetbuttcap%
\pgfsetroundjoin%
\definecolor{currentfill}{rgb}{0.000000,0.000000,0.000000}%
\pgfsetfillcolor{currentfill}%
\pgfsetlinewidth{0.501875pt}%
\definecolor{currentstroke}{rgb}{0.000000,0.000000,0.000000}%
\pgfsetstrokecolor{currentstroke}%
\pgfsetdash{}{0pt}%
\pgfsys@defobject{currentmarker}{\pgfqpoint{0.000000in}{0.000000in}}{\pgfqpoint{0.020833in}{0.000000in}}{%
\pgfpathmoveto{\pgfqpoint{0.000000in}{0.000000in}}%
\pgfpathlineto{\pgfqpoint{0.020833in}{0.000000in}}%
\pgfusepath{stroke,fill}%
}%
\begin{pgfscope}%
\pgfsys@transformshift{1.355382in}{3.399977in}%
\pgfsys@useobject{currentmarker}{}%
\end{pgfscope}%
\end{pgfscope}%
\begin{pgfscope}%
\pgfsetbuttcap%
\pgfsetroundjoin%
\definecolor{currentfill}{rgb}{0.000000,0.000000,0.000000}%
\pgfsetfillcolor{currentfill}%
\pgfsetlinewidth{0.501875pt}%
\definecolor{currentstroke}{rgb}{0.000000,0.000000,0.000000}%
\pgfsetstrokecolor{currentstroke}%
\pgfsetdash{}{0pt}%
\pgfsys@defobject{currentmarker}{\pgfqpoint{-0.020833in}{0.000000in}}{\pgfqpoint{-0.000000in}{0.000000in}}{%
\pgfpathmoveto{\pgfqpoint{-0.000000in}{0.000000in}}%
\pgfpathlineto{\pgfqpoint{-0.020833in}{0.000000in}}%
\pgfusepath{stroke,fill}%
}%
\begin{pgfscope}%
\pgfsys@transformshift{6.103850in}{3.399977in}%
\pgfsys@useobject{currentmarker}{}%
\end{pgfscope}%
\end{pgfscope}%
\begin{pgfscope}%
\pgfsetbuttcap%
\pgfsetroundjoin%
\definecolor{currentfill}{rgb}{0.000000,0.000000,0.000000}%
\pgfsetfillcolor{currentfill}%
\pgfsetlinewidth{0.501875pt}%
\definecolor{currentstroke}{rgb}{0.000000,0.000000,0.000000}%
\pgfsetstrokecolor{currentstroke}%
\pgfsetdash{}{0pt}%
\pgfsys@defobject{currentmarker}{\pgfqpoint{0.000000in}{0.000000in}}{\pgfqpoint{0.020833in}{0.000000in}}{%
\pgfpathmoveto{\pgfqpoint{0.000000in}{0.000000in}}%
\pgfpathlineto{\pgfqpoint{0.020833in}{0.000000in}}%
\pgfusepath{stroke,fill}%
}%
\begin{pgfscope}%
\pgfsys@transformshift{1.355382in}{3.649527in}%
\pgfsys@useobject{currentmarker}{}%
\end{pgfscope}%
\end{pgfscope}%
\begin{pgfscope}%
\pgfsetbuttcap%
\pgfsetroundjoin%
\definecolor{currentfill}{rgb}{0.000000,0.000000,0.000000}%
\pgfsetfillcolor{currentfill}%
\pgfsetlinewidth{0.501875pt}%
\definecolor{currentstroke}{rgb}{0.000000,0.000000,0.000000}%
\pgfsetstrokecolor{currentstroke}%
\pgfsetdash{}{0pt}%
\pgfsys@defobject{currentmarker}{\pgfqpoint{-0.020833in}{0.000000in}}{\pgfqpoint{-0.000000in}{0.000000in}}{%
\pgfpathmoveto{\pgfqpoint{-0.000000in}{0.000000in}}%
\pgfpathlineto{\pgfqpoint{-0.020833in}{0.000000in}}%
\pgfusepath{stroke,fill}%
}%
\begin{pgfscope}%
\pgfsys@transformshift{6.103850in}{3.649527in}%
\pgfsys@useobject{currentmarker}{}%
\end{pgfscope}%
\end{pgfscope}%
\begin{pgfscope}%
\pgfsetbuttcap%
\pgfsetroundjoin%
\definecolor{currentfill}{rgb}{0.000000,0.000000,0.000000}%
\pgfsetfillcolor{currentfill}%
\pgfsetlinewidth{0.501875pt}%
\definecolor{currentstroke}{rgb}{0.000000,0.000000,0.000000}%
\pgfsetstrokecolor{currentstroke}%
\pgfsetdash{}{0pt}%
\pgfsys@defobject{currentmarker}{\pgfqpoint{0.000000in}{0.000000in}}{\pgfqpoint{0.020833in}{0.000000in}}{%
\pgfpathmoveto{\pgfqpoint{0.000000in}{0.000000in}}%
\pgfpathlineto{\pgfqpoint{0.020833in}{0.000000in}}%
\pgfusepath{stroke,fill}%
}%
\begin{pgfscope}%
\pgfsys@transformshift{1.355382in}{3.899076in}%
\pgfsys@useobject{currentmarker}{}%
\end{pgfscope}%
\end{pgfscope}%
\begin{pgfscope}%
\pgfsetbuttcap%
\pgfsetroundjoin%
\definecolor{currentfill}{rgb}{0.000000,0.000000,0.000000}%
\pgfsetfillcolor{currentfill}%
\pgfsetlinewidth{0.501875pt}%
\definecolor{currentstroke}{rgb}{0.000000,0.000000,0.000000}%
\pgfsetstrokecolor{currentstroke}%
\pgfsetdash{}{0pt}%
\pgfsys@defobject{currentmarker}{\pgfqpoint{-0.020833in}{0.000000in}}{\pgfqpoint{-0.000000in}{0.000000in}}{%
\pgfpathmoveto{\pgfqpoint{-0.000000in}{0.000000in}}%
\pgfpathlineto{\pgfqpoint{-0.020833in}{0.000000in}}%
\pgfusepath{stroke,fill}%
}%
\begin{pgfscope}%
\pgfsys@transformshift{6.103850in}{3.899076in}%
\pgfsys@useobject{currentmarker}{}%
\end{pgfscope}%
\end{pgfscope}%
\begin{pgfscope}%
\pgfpathrectangle{\pgfqpoint{1.355382in}{0.155832in}}{\pgfqpoint{4.748468in}{3.743244in}}%
\pgfusepath{clip}%
\pgfsetbuttcap%
\pgfsetroundjoin%
\pgfsetlinewidth{0.803000pt}%
\definecolor{currentstroke}{rgb}{0.000000,0.000000,0.000000}%
\pgfsetstrokecolor{currentstroke}%
\pgfsetdash{{0.800000pt}{1.320000pt}}{0.000000pt}%
\pgfpathmoveto{\pgfqpoint{5.484484in}{0.405382in}}%
\pgfpathlineto{\pgfqpoint{5.484484in}{2.900878in}}%
\pgfusepath{stroke}%
\end{pgfscope}%
\begin{pgfscope}%
\pgfpathrectangle{\pgfqpoint{1.355382in}{0.155832in}}{\pgfqpoint{4.748468in}{3.743244in}}%
\pgfusepath{clip}%
\pgfsetbuttcap%
\pgfsetroundjoin%
\pgfsetlinewidth{0.803000pt}%
\definecolor{currentstroke}{rgb}{0.000000,0.000000,0.000000}%
\pgfsetstrokecolor{currentstroke}%
\pgfsetdash{{0.800000pt}{1.320000pt}}{0.000000pt}%
\pgfpathmoveto{\pgfqpoint{1.355382in}{2.900878in}}%
\pgfpathlineto{\pgfqpoint{5.484484in}{2.900878in}}%
\pgfusepath{stroke}%
\end{pgfscope}%
\begin{pgfscope}%
\pgfpathrectangle{\pgfqpoint{1.355382in}{0.155832in}}{\pgfqpoint{4.748468in}{3.743244in}}%
\pgfusepath{clip}%
\pgfsetbuttcap%
\pgfsetroundjoin%
\pgfsetlinewidth{0.803000pt}%
\definecolor{currentstroke}{rgb}{0.000000,0.000000,0.000000}%
\pgfsetstrokecolor{currentstroke}%
\pgfsetdash{{0.800000pt}{1.320000pt}}{0.000000pt}%
\pgfpathmoveto{\pgfqpoint{4.245754in}{0.405382in}}%
\pgfpathlineto{\pgfqpoint{4.245754in}{1.902679in}}%
\pgfusepath{stroke}%
\end{pgfscope}%
\begin{pgfscope}%
\pgfpathrectangle{\pgfqpoint{1.355382in}{0.155832in}}{\pgfqpoint{4.748468in}{3.743244in}}%
\pgfusepath{clip}%
\pgfsetbuttcap%
\pgfsetroundjoin%
\pgfsetlinewidth{0.803000pt}%
\definecolor{currentstroke}{rgb}{0.000000,0.000000,0.000000}%
\pgfsetstrokecolor{currentstroke}%
\pgfsetdash{{0.800000pt}{1.320000pt}}{0.000000pt}%
\pgfpathmoveto{\pgfqpoint{1.355382in}{1.403580in}}%
\pgfpathlineto{\pgfqpoint{4.245754in}{1.403580in}}%
\pgfusepath{stroke}%
\end{pgfscope}%
\begin{pgfscope}%
\pgfpathrectangle{\pgfqpoint{1.355382in}{0.155832in}}{\pgfqpoint{4.748468in}{3.743244in}}%
\pgfusepath{clip}%
\pgfsetbuttcap%
\pgfsetroundjoin%
\pgfsetlinewidth{0.803000pt}%
\definecolor{currentstroke}{rgb}{0.000000,0.000000,0.000000}%
\pgfsetstrokecolor{currentstroke}%
\pgfsetdash{{0.800000pt}{1.320000pt}}{0.000000pt}%
\pgfpathmoveto{\pgfqpoint{1.355382in}{1.902679in}}%
\pgfpathlineto{\pgfqpoint{4.245754in}{1.902679in}}%
\pgfusepath{stroke}%
\end{pgfscope}%
\begin{pgfscope}%
\pgfsetbuttcap%
\pgfsetmiterjoin%
\definecolor{currentfill}{rgb}{0.000000,0.000000,0.000000}%
\pgfsetfillcolor{currentfill}%
\pgfsetlinewidth{1.003750pt}%
\definecolor{currentstroke}{rgb}{0.000000,0.000000,0.000000}%
\pgfsetstrokecolor{currentstroke}%
\pgfsetdash{}{0pt}%
\pgfsys@defobject{currentmarker}{\pgfqpoint{-0.041667in}{-0.041667in}}{\pgfqpoint{0.041667in}{0.041667in}}{%
\pgfpathmoveto{\pgfqpoint{0.041667in}{-0.000000in}}%
\pgfpathlineto{\pgfqpoint{-0.041667in}{0.041667in}}%
\pgfpathlineto{\pgfqpoint{-0.041667in}{-0.041667in}}%
\pgfpathlineto{\pgfqpoint{0.041667in}{-0.000000in}}%
\pgfpathclose%
\pgfusepath{stroke,fill}%
}%
\begin{pgfscope}%
\pgfsys@transformshift{6.103850in}{0.405382in}%
\pgfsys@useobject{currentmarker}{}%
\end{pgfscope}%
\end{pgfscope}%
\begin{pgfscope}%
\pgfsetbuttcap%
\pgfsetmiterjoin%
\definecolor{currentfill}{rgb}{0.000000,0.000000,0.000000}%
\pgfsetfillcolor{currentfill}%
\pgfsetlinewidth{1.003750pt}%
\definecolor{currentstroke}{rgb}{0.000000,0.000000,0.000000}%
\pgfsetstrokecolor{currentstroke}%
\pgfsetdash{}{0pt}%
\pgfsys@defobject{currentmarker}{\pgfqpoint{-0.041667in}{-0.041667in}}{\pgfqpoint{0.041667in}{0.041667in}}{%
\pgfpathmoveto{\pgfqpoint{0.000000in}{0.041667in}}%
\pgfpathlineto{\pgfqpoint{-0.041667in}{-0.041667in}}%
\pgfpathlineto{\pgfqpoint{0.041667in}{-0.041667in}}%
\pgfpathlineto{\pgfqpoint{0.000000in}{0.041667in}}%
\pgfpathclose%
\pgfusepath{stroke,fill}%
}%
\begin{pgfscope}%
\pgfsys@transformshift{1.355382in}{3.899076in}%
\pgfsys@useobject{currentmarker}{}%
\end{pgfscope}%
\end{pgfscope}%
\begin{pgfscope}%
\pgfsetrectcap%
\pgfsetmiterjoin%
\pgfsetlinewidth{0.501875pt}%
\definecolor{currentstroke}{rgb}{0.000000,0.000000,0.000000}%
\pgfsetstrokecolor{currentstroke}%
\pgfsetdash{}{0pt}%
\pgfpathmoveto{\pgfqpoint{1.355382in}{0.155832in}}%
\pgfpathlineto{\pgfqpoint{1.355382in}{3.899076in}}%
\pgfusepath{stroke}%
\end{pgfscope}%
\begin{pgfscope}%
\pgfsetrectcap%
\pgfsetmiterjoin%
\pgfsetlinewidth{0.501875pt}%
\definecolor{currentstroke}{rgb}{0.000000,0.000000,0.000000}%
\pgfsetstrokecolor{currentstroke}%
\pgfsetdash{}{0pt}%
\pgfpathmoveto{\pgfqpoint{1.355382in}{0.405382in}}%
\pgfpathlineto{\pgfqpoint{6.103850in}{0.405382in}}%
\pgfusepath{stroke}%
\end{pgfscope}%
\begin{pgfscope}%
\definecolor{textcolor}{rgb}{0.000000,0.000000,0.000000}%
\pgfsetstrokecolor{textcolor}%
\pgfsetfillcolor{textcolor}%
\pgftext[x=6.103850in,y=0.280607in,right,top]{\color{textcolor}{\rmfamily\fontsize{10.000000}{12.000000}\selectfont\catcode`\^=\active\def^{\ifmmode\sp\else\^{}\fi}\catcode`\%=\active\def%{\%}$x$}}%
\end{pgfscope}%
\begin{pgfscope}%
\definecolor{textcolor}{rgb}{0.000000,0.000000,0.000000}%
\pgfsetstrokecolor{textcolor}%
\pgfsetfillcolor{textcolor}%
\pgftext[x=1.252155in,y=3.899076in,right,top]{\color{textcolor}{\rmfamily\fontsize{10.000000}{12.000000}\selectfont\catcode`\^=\active\def^{\ifmmode\sp\else\^{}\fi}\catcode`\%=\active\def%{\%}$y$}}%
\end{pgfscope}%
\begin{pgfscope}%
\definecolor{textcolor}{rgb}{0.000000,0.000000,0.000000}%
\pgfsetstrokecolor{textcolor}%
\pgfsetfillcolor{textcolor}%
\pgftext[x=6.103850in,y=3.649527in,right,bottom]{\color{textcolor}{\rmfamily\fontsize{10.000000}{12.000000}\selectfont\catcode`\^=\active\def^{\ifmmode\sp\else\^{}\fi}\catcode`\%=\active\def%{\%}$y = h(x)$}}%
\end{pgfscope}%
\begin{pgfscope}%
\pgfpathrectangle{\pgfqpoint{1.355382in}{0.155832in}}{\pgfqpoint{4.748468in}{3.743244in}}%
\pgfusepath{clip}%
\pgfsetbuttcap%
\pgfsetroundjoin%
\pgfsetlinewidth{1.003750pt}%
\definecolor{currentstroke}{rgb}{0.000000,0.000000,0.000000}%
\pgfsetstrokecolor{currentstroke}%
\pgfsetdash{{3.700000pt}{1.600000pt}}{0.000000pt}%
\pgfpathmoveto{\pgfqpoint{2.387658in}{0.405382in}}%
\pgfpathlineto{\pgfqpoint{5.484484in}{2.900878in}}%
\pgfusepath{stroke}%
\end{pgfscope}%
\begin{pgfscope}%
\pgfpathrectangle{\pgfqpoint{1.355382in}{0.155832in}}{\pgfqpoint{4.748468in}{3.743244in}}%
\pgfusepath{clip}%
\pgfsetrectcap%
\pgfsetroundjoin%
\pgfsetlinewidth{2.007500pt}%
\definecolor{currentstroke}{rgb}{0.000000,0.000000,0.000000}%
\pgfsetstrokecolor{currentstroke}%
\pgfsetdash{}{0pt}%
\pgfpathmoveto{\pgfqpoint{1.355382in}{0.405382in}}%
\pgfpathlineto{\pgfqpoint{3.418280in}{0.405382in}}%
\pgfpathlineto{\pgfqpoint{3.423033in}{0.409129in}}%
\pgfpathlineto{\pgfqpoint{6.103850in}{3.649527in}}%
\pgfpathlineto{\pgfqpoint{6.103850in}{3.649527in}}%
\pgfusepath{stroke}%
\end{pgfscope}%
\end{pgfpicture}%
\makeatother%
\endgroup%

	\caption{The convex function \(h(x) = (x-K)^+\).}\label{fig:ConvexFunctionPayoff}
\end{figure}

\begin{thm}{Price of American Derivative}{PriceOfAmericanDerivative}
	Let \(h(x)\) be a nonnegative, convex function of \(x \ge 0\) satisfying \(h(0) = 0\). Then the price of the American derivative security expiring at time \(T\) and having intrinsic value \(h(S(t))\), \(0 \le t \le T\), is the same as the price of the European derivative security paying \(h(S(T))\) at expiration \(T\).
\end{thm}

\begin{proof}
	Replacing \(t\) by \(T\) in \Cref{eq:PreDiscountedSubmartingaleInequality}, we obtain
	\[
		\tilde{\E}\left[e^{-r(T-u)}h(S(T)) \gvn \cF(u)\right] \ge h(S(u)), \quad 0 \le u \le T.
	\]
	In other words, the European derivative security price always dominates the intrinsic value of the American derivative security. This shows that the option to exercise early is worthless, and the price of the American derivative security agrees with the price of the European security.
\end{proof}

\begin{cor}{American Call Price}{AmericanCallPrice}
	The price of an American call on an asset not paying a dividend is the same as the price of the European call on the same asset with the same expiration.
\end{cor}

\begin{proof}
	Take \(h(x) = (x - K)^+\) in \Cref{thm:PriceOfAmericanDerivative}.
\end{proof}

The idea behind \Cref{cor:AmericanCallPrice} is that the discounted process \(e^{-rt}(S(t) - K)^+\) is a submartingale under \(\tilde{\Prob}\) and hence tends to rise. Therefore, it is optimal to wait until expiration before deciding whether to exercise.

There are two factors that contribute to the submartingale property for \(e^{-rt}(S(t)-K)^+\). One is the discounting of the strike. In fact, \(e^{-rt}(S(t)-K)\) (without the \(^+\)) is a submartingale because \(e^{-rt}S(t)\) is a martingale under the risk-neutral measure \(\tilde{\Prob}\) and \(-e^{-rt}K\) increases as \(t\) increases (throughout this chapter, we assume a strictly positive interest rate \(r\)). When we reinstate the \(^+\), we are taking a convex function of a submartingale and, because of Jensen's inequality, this reinforces the upward trend.

The previous argument does not apply to the American put, whose discounted intrinsic value \(e^{-rt}(K - S(t))\) (without the \(^+\)) is a supermartingale (\(e^{-rt}K\) falls and \(-e^{-rt}S(t)\) is a martingale). Jensen's inequality creates an upward trend that competes with this supermartingale property, and the analysis becomes complicated.

If the underlying asset pays a dividend, the case considered in the next subsection, the argument above no longer applies to the American call. In this case, \(e^{-rt}S(t)\) is a supermartingale and tends to fall because of the dividend outflow.

\subsubsection{Underlying Asset Pays Dividends}

In this subsection, we consider an American call on an asset whose price process is a geometric Brownian motion governed by (8.5.1) between dividend payment dates.

We assume there are times \(0 < t_1 < t_2 < \cdots < t_n < T\), and at each time \(t_j\) the dividend paid is \(a_j S(t_j-)\), where \(S(t_j-)\) denotes the asset price just prior to the dividend payment. The asset price \(S(t_j)\) after the dividend payment is the asset price before the dividend payment less the dividend payment:
\begin{equation}
	S(t_j) = S(t_j-) - a_j S(t_j-) = (1 - a_j)S(t_j-). \label{eq:DividendAdjustment}
\end{equation}

We assume that each \(a_j, j=1,\dots,n\), is a number between 0 and 1. We set \(t_0 = 0\), but this is not a dividend payment date. We also assume that \(T\) is not a dividend payment date, although it is not difficult to modify the analysis given below to handle the case when \(T\) is a dividend payment date.

We shall see that it is not optimal to exercise an American call on this asset except possibly immediately before a dividend payment. The price of the call will be seen to satisfy the Black-Scholes-Merton partial differential equation between dividend payment dates. At dividend payment dates, the price of the call is the maximum of the call's intrinsic value and the price of the call after the dividend is paid and the stock price is reduced by the amount of the payment. These observations lead to a recursive algorithm for determining the price, and that is developed in this subsection.

For \(t_j \le t < t_{j+1}\), we have
\begin{equation}
	S(t) = S(t_j) \exp \set[\bigg]{ \sigma (\widetilde{W}(t) - \widetilde{W}(t_j)) + \left( r - \frac{1}{2}\sigma^2 \right) (t - t_j) }, \label{eq:AssetPriceDynamics}
\end{equation}
which implies
\begin{equation}
	S(t_{j+1}-) = S(t_j) \exp \set[\bigg]{ \sigma (\widetilde{W}(t_{j+1}) - \widetilde{W}(t_j)) + \left( r - \frac{1}{2}\sigma^2 \right) (t_{j+1} - t_j) } \label{eq:AssetPricePreDividend}
\end{equation}
and
\begin{equation}
	S(t_{j+1}) = (1 - a_{j+1})S(t_j) \exp \set[\bigg]{ \sigma (\widetilde{W}(t_{j+1}) - \widetilde{W}(t_j)) + \left( r - \frac{1}{2}\sigma^2 \right) (t_{j+1} - t_j) }. \label{eq:AssetPricePostDividend}
\end{equation}
We also have
\begin{equation}
	S(T) = S(t_n) \exp \set[\bigg]{ \sigma (\widetilde{W}(T) - \widetilde{W}(t_n)) + \left( r - \frac{1}{2}\sigma^2 \right) (T - t_n) }. \label{eq:TerminalAssetPrice}
\end{equation}

We consider an American call expiring at time \(T\) with strike price \(K\). For \(t_n \le t \le T\), the discounted asset price \(\e^{-rt}S(t)\) is a martingale under \(\widetilde{\Prob}\), and \Cref{lem:DiscountedIntrinsicValueSubmartingale} can be invoked to show that \(\e^{-rt}(S(t) - K)^+\) is a submartingale. Therefore,
\begin{equation}
	\widetilde{\E} \left[ \e^{-rT} (S(T) - K)^+ \gvn \cF(t) \right] \ge \e^{-rt} (S(t) - K)^+, \quad t_n \le t \le T. \label{eq:SubmartingaleProperty}
\end{equation}
This shows that, for all \(t \in [t_n, T]\), the price of the European call at time \(t\),
\begin{equation}
	c_n(t, S(t)) = \widetilde{\E} \left[ \e^{-r(T-t)} (S(T) - K)^+ \gvn \cF(t) \right], \label{eq:EuropeanCallValue}
\end{equation}

is greater than the intrinsic value of the American call, \((S(t) - K)^+\). Consequently, the early exercise feature of the American call is worthless, and the prices at time \(t\) of the European and American calls agree for \(t_n \le t \le T\). This price is given by the Black-Scholes-Merton formula
\begin{equation}
	c_n(t, x) = x N(d_+(T - t, x)) - K \e^{-r(T-t)} N(d_-(T - t, x)), \label{eq:BSMFormula}
\end{equation}
where
\[
	d_{\pm}(\tau, x) = \frac{1}{\sigma\sqrt{\tau}} \left[ \log \frac{x}{K} + \left( r \pm \frac{1}{2}\sigma^2 \right) \tau \right].
\]
Although one cannot simply substitute \(x = 0\) into (8.5.14), we have \(c(t, 0) = 0\); see equation (4.5.17) and Exercise 4.9. Formula (8.5.14) can be determined by computing the conditional expectation in (8.5.13) under the condition \(S(t) = x\). In the case \(t = t_n\), using (8.5.12), this leads to
\begin{equation}
	c_n(t_n, x) = \widetilde{\E} \left[ \e^{-r(T-t_n)} \left( x \exp \set[\bigg]{ \sigma (\widetilde{W}(T) - \widetilde{W}(t_n)) + \left( r - \frac{1}{2}\sigma^2 \right) (T - t_n) } - K \right)^+ \right]. \label{eq:ConditionalExpectationTn}
\end{equation}
The function \(c_n(t, x)\) also satisfies the Black-Scholes-Merton differential equation
\begin{equation}
	\frac{\partial}{\partial t} c_n(t, x) + rx \frac{\partial}{\partial x} c_n(t, x) + \frac{1}{2}\sigma^2 x^2 \frac{\partial^2}{\partial x^2} c_n(t, x) = r c_n(t, x), \quad t_n \le t < T, \ x \ge 0, \label{eq:BSMPDE}
\end{equation}
and the terminal condition
\begin{equation}
	c_n(t, x) = (x - K)^+, \ x \ge 0. \label{eq:TerminalCondition}
\end{equation}
The function \(c_n(t_n, x)\) is convex in \(x\). This is well-known, but we establish it here anyway to demonstrate a method we need later. To show convexity in \(x\), we show that, whenever \(0 \le x_1 \le x_2\) and \(0 \le \lambda \le 1\), we have
\begin{equation}
	c_n(t_n, (1 - \lambda)x_1 + \lambda x_2) \le (1 - \lambda)c_n(t_n, x_1) + \lambda c_n(t_n, x_2). \label{eq:ConvexityInequality}
\end{equation}
We begin with the observation that, for any number \(\alpha\), the function \((\alpha x - K)^+\) is convex in \(x\), and therefore
\[
	\left( x \exp \set[\bigg]{ \sigma (\widetilde{W}(T) - \widetilde{W}(t_n)) + \left( r - \frac{1}{2}\sigma^2 \right) (T - t_n) } - K \right)^+
\]
is convex in \(x\). It follows that
\begin{align}
	 & c_n(t_n, (1 - \lambda)x_1 + \lambda x_2) \nonumber                                                                                                                                                           \\
	 & = \widetilde{\E} \Bigg[ \e^{-r(T-t_n)} \bigg( ((1 - \lambda)x_1 + \lambda x_2) \exp\Bigg\{ \sigma (\widetilde{W}(T) - \widetilde{W}(t_n)) \nonumber                                                          \\
	 & \qquad + \left( r - \frac{1}{2}\sigma^2 \right) \Bigg\} - K \bigg)^+ \Bigg] \nonumber                                                                                                                        \\
	 & \le (1 - \lambda)\widetilde{\E} \left[ \e^{-r(T-t_n)} \left( x_1 \exp \set[\bigg]{ \sigma (\widetilde{W}(T) - \widetilde{W}(t_n)) + \left( r - \frac{1}{2}\sigma^2 \right) } - K \right)^+ \right] \nonumber \\
	 & \qquad + \lambda\widetilde{\E} \left[ \e^{-r(T-t_n)} \left( x_2 \exp \set[\bigg]{ \sigma (\widetilde{W}(T) - \widetilde{W}(t_n)) + \left( r - \frac{1}{2}\sigma^2 \right) } - K \right)^+ \right] \nonumber  \\
	 & = (1 - \lambda)c_n(t_n, x_1) + \lambda c_n(t_n, x_2). \label{eq:ConvexityProof}
\end{align}

This proves \Cref{eq:ConvexityInequality}.

At time \(t_n\), immediately before the dividend payment, the owner of the American call has two choices. She can exercise the option and receive \(S(t_n-) - K\), or she can decline to exercise, permit the dividend to be paid (not to her) and the asset price to fall to \(S(t_n) = (1 - a_n)S(t_n-)\), and have an option valued at \(c_n(t_n, (1 - a_n)S(t_n-))\). The optimal decision is to exercise if \(S(t_n-) - K > c_n(t_n, (1 - a_n)S(t_n-))\) and to decline to exercise if \(S(t_n-) - K < c_n(t_n, (1 - a_n)S(t_n-))\). If \(S(t_n-) - K = c_n(t_n, (1 - a_n)S(t_n-))\), it does not matter whether she exercises or declines to exercise. Therefore, the call value at time \(t_n\) immediately before the dividend is paid is \(h_n(S(t_n-))\), where
\begin{equation}
	h_n(x) = \max \set{ x - K, c_n(t_n, (1 - a_n)x) }, \ x \ge 0. \label{eq:ValueFunctionHn}
\end{equation}

We show that \(h_n(x)\) satisfies the assumptions of \Cref{lem:DiscountedIntrinsicValueSubmartingale}. It is clear that \(h_n(x) \ge 0\) for all \(x \ge 0\) because \(c_n(t_n, (1 - a_n)x) \ge 0\) for all \(x \ge 0\). It is also clear that \(h_n(0) = 0\) because \(c_n(t_n, (1 - a_n)0) = 0\). To establish the convexity of \(h_n(x)\), we recall from \Cref{eq:ConvexityInequality} that \(c_n(t_n, x)\) is convex in \(x\). For \(0 \le x_1 \le x_2\) and \(0 \le \lambda \le 1\), we replace \(x_1\) in \Cref{eq:ConvexityInequality} by \((1 - a_n)x_1\) and replace \(x_2\) by \((1 - a_n)x_2\) to obtain
\[
	c_n(t_n, (1 - a_n)((1 - \lambda)x_1 + \lambda x_2)) \le (1 - \lambda)c_n(t_n, (1 - a_n)x_1) + \lambda c_n(t_n, (1 - a_n)x_2).
\]
This shows that \(c_n(t, (1 - a_n)x)\) is a convex function of \(x\). The maximum of two convex functions is convex, and therefore \(h_n(x)\) defined by \Cref{eq:ValueFunctionHn} is convex.

Starting from time \(t\), where \(t_{n-1} \le t < t_n\), the owner of the American call can exercise at any time \(u \in [t, t_n)\), and if she does, she receives \(S(u) - K\). If she does not exercise prior to \(t_n\), then at time \(t_n\), immediately before the dividend payment, she owns a call whose value we have just determined to be \(h_n(S(t_n-))\). Therefore, for \(t_{n-1} \le t < t_n\), the American call expiring at time \(T\) has the same price as the American call expiring immediately before the dividend payment at date \(t_n\) and paying \(h_n(S(t_n-))\) upon expiration.

Because the underlying asset evolves as a geometric Brownian motion after the dividend is paid at time \(t_{n-1}\) until the dividend is paid at time \(t_n\), \Cref{lem:DiscountedIntrinsicValueSubmartingale} implies that \(e^{-rt} h_n (S(t))\) is a submartingale for \(t_{n-1} \leq t < t_n\).

In particular,
\[
	\tilde{\E} \left[ e^{-r(u-t)}h_n(S(u)) \mid \cF(t) \right] \geq h_n(S(t)), \quad t_{n-1} \leq t \leq u < t_n,
\]
and letting \(u \uparrow t_n\), we obtain
\begin{equation}
	\tilde{\E} \left[ e^{-r(t_n-t)}h_n(S(t_n-)) \mid \cF(t) \right] \geq h_n(S(t)). \label{eq:SubmartingaleLimitCondition}
\end{equation}

By the definition of \(h_n(x)\),
\begin{equation}
	h_n(S(t)) \geq S(t) - K. \label{eq:HnIntrinsicBound}
\end{equation}

This shows that the value of the European call expiring at time \(t_n\) immediately before the dividend is paid and paying \(h_n(S(t_n-))\) upon expiration, which is the left-hand side of~\eqref{eq:SubmartingaleLimitCondition}, is greater than or equal to the intrinsic value of the American call, which is the right-hand side of~\eqref{eq:HnIntrinsicBound}. Therefore, the option to exercise the American call before time \(t_n\) is worthless, and the American call value is the same as the value of the European call just described.

Because \(S(t)\) is a Markov process, there is some function \(c_{n-1}(t, x)\) such that the left-hand side of~\eqref{eq:SubmartingaleLimitCondition}, the European call value, is
\begin{equation}
	c_{n-1}(t, S(t)) = \tilde{\E} \left[ e^{-r(t_n-t)}h_n(S(t_n-)) \mid \cF(t) \right]. \label{eq:EuropeanCallFunction}
\end{equation}

The function \(c_{n-1}(t, x)\) can be determined by computing the conditional expectation in~\eqref{eq:EuropeanCallFunction} under the condition \(S(t) = x\). In the case \(t = t_{n-1}\), using \Cref{eq:AssetPricePreDividend}, this leads to
\begin{align}
	c_{n-1}(t_{n-1}, x)
	 & = \tilde{\E} \bigg[ e^{-r(t_n-t_{n-1})}
		\times h_n \left( x \exp \set[\bigg]{ \sigma (\tilde{W}(t_n) - \tilde{W}(t_{n-1})) + \left( r - \frac{1}{2}\sigma^2 \right) (t_n - t_{n-1}) } \right) \bigg]. \label{eq:RecursiveExpectationFormula}
\end{align}

The function \(c_{n-1}(t, x)\) also satisfies the Black-Scholes-Merton differential equation
\begin{equation}
	\frac{\partial}{\partial t}c_{n-1}(t, x) + rx \frac{\partial}{\partial x}c_{n-1}(t, x) + \frac{1}{2}\sigma^2 x^2 \frac{\partial^2}{\partial x^2}c_{n-1}(t, x) = r c_{n-1}(t, x), \quad t_{n-1} \leq t < t_n, \quad x \geq 0. \label{eq:BlackScholesMertonPDE}
\end{equation}

and the terminal condition

\begin{equation}
	c_{n-1}(t_n, x) = h_n(x), \quad x \ge 0.
	\label{eq:CNMinusOneTerminalAtTn}
\end{equation}

We repeat this process, defining
\[
	h_{n-1}(x) = \max \set{ x - K, c_{n-1}(t_{n-1}, (1 - a_{n-1})x) }, \quad x \ge 0.
\]
We can show as above that \(h_{n-1}(x)\) satisfies the hypotheses of Lemma 8.5.1, and we continue.

In conclusion, we obtain an algorithm for the American call price on an asset paying dividends at the dates \(t_1, t_2, \dots, t_n\). Solve recursively for \(j = n, n-1, \dots, 0\), the partial differential equation
\begin{equation}
	\begin{split}
		\frac{\partial}{\partial t}c_{j-1}(t, x)
		+ rx \frac{\partial}{\partial x}c_{j-1}(t, x)
		+ \frac{1}{2}\sigma^2 x^2 \frac{\partial^2}{\partial x^2}c_{j-1}(t, x)
		= r c_{j-1}(t, x), \quad t_{j-1} \le t < t_j, \quad x \ge 0,
	\end{split}
	\label{eq:AmericanCallRecursionPDE}
\end{equation}
with the terminal condition
\begin{equation}
	c_{j-1}(t_j, x) = h_j(x), \quad x \ge 0.
	\label{eq:AmericanCallRecursionTerminal}
\end{equation}

The functions \(c_n(t, x)\) and \(h_n(x)\) needed to get started are given by \Crefrange{eq:BSMFormula}{eq:ValueFunctionHn}, and the function \(h_{j-1}(x)\) needed for the next step is given by
\begin{equation}
	h_{j-1}(x) = \max \set{ x - K, c_{j-1}(t_{j-1}, (1 - a_{j-1})x) }, \quad x \ge 0.
	\label{eq:HJMinusOneDefinition}
\end{equation}

For \(t_{j-1} \le t < t_j\), if \(S(t) = x\), then \(c_{j-1}(t, x)\) is the American call price. Within each interval \([t_{j-1}, t_j)\), the American call price is actually the price of a European call expiring at \(t_j\). The optimal exercise time is immediately prior to the dividend payment at the smallest time \(t_j\) for which \(S(t_j-) - K\) exceeds \(c_j(t_j, (1 - a_j)S(t_j-))\). If there is no \(t_j\) for which this condition is satisfied, then optimal exercise takes place at time \(T\) if \(S(T) > K\), and otherwise the option should be allowed to expire unexercised.

\subsection{Approximate Continuation Values}
Regression-based methods posit an expression for the continuation value of the form
\begin{equation} \label{eq:ContinuationValueExpectation}
	\E[V_{i+1}(X_{i+1}) \mid X_i = x] = \sum_{r=1}^M \beta_{ir} \psi_r(x),
\end{equation}
for some basis functions \(\psi_r : \R^d \to \R\) and constants \(\beta_{ir}, r = 1, \dots, M\). Using the notation \(C_i\) for the continuation value at time \(i\), we may equivalently write \Cref{eq:ContinuationValueExpectation} as
\begin{equation} \label{eq:ContinuationValueVectorForm}
	C_i(x) = \beta_i^\top \psi(x),
\end{equation}
with
\begin{equation*}
	\beta_i^\top = (\beta_{i1}, \dots, \beta_{iM}), \quad \psi(x) = (\psi_1(x), \dots, \psi_M(x))^\top.
\end{equation*}
One could use different basis functions at different exercise dates, but to simplify notation, we suppress any dependence of \(\psi\) on \(i\).

Assuming a relation of the form \Cref{eq:ContinuationValueExpectation} holds, the vector \(\beta_i\) is given by
\begin{equation} \label{eq:BetaVectorDefinition}
	\beta_i = (\E[\psi(X_i)\psi(X_i)^\top])^{-1} \E[\psi(X_i)V_{i+1}(X_{i+1})] \equiv B_\psi^{-1} B_{\psi V}.
\end{equation}
Here, \(B_\psi\) is the indicated \(M \times M\) matrix (assumed nonsingular) and \(B_{\psi V}\) the indicated vector of length \(M\). The variables \((X_i, X_{i+1})\) inside the expectations have a joint distribution corresponding to the state of the underlying Markov chain at dates \(i\) and \(i+1\).

The coefficients \(\beta_{ir}\) could be estimated from observations of pairs \((X_{ij}, V_{i+1}(X_{i+1,j}))\), \(j = 1, \dots, b\), each consisting of the state at time \(i\) and the corresponding option value at time \(i+1\). Consider, in particular, independent paths \((X_{1j}, \dots, X_{mj})\), \(j = 1, \dots, b\), and suppose for a moment that the values \(V_{i+1}(X_{i+1,j})\) are known. The least-squares estimate of \(\beta_i\) is then given by
\begin{equation*}
	\hat{\beta}_i = \hat{B}_\psi^{-1} \hat{B}_{\psi V},
\end{equation*}
where \(\hat{B}_\psi\) and \(\hat{B}_{\psi V}\) are the sample counterparts of \(B_\psi\) and \(B_{\psi V}\). More explicitly, \(\hat{B}_\psi\) is the \(M \times M\) matrix with \(qr\) entry
\begin{equation*}
	\frac{1}{b} \sum_{j=1}^b \psi_q(X_{ij}) \psi_r(X_{ij})
\end{equation*}
and \(\hat{B}_{\psi V}\) is the \(M\)-vector with \(r\)th entry
\begin{equation} \label{eq:SampleCounterpartVector}
	\frac{1}{b} \sum_{k=1}^b \psi_r(X_{ik}) V_{i+1}(X_{i+1, k}).
\end{equation}
All of these quantities can be calculated from function values at pairs of consecutive nodes \((X_{ij}, X_{i+1, j})\), \(j = 1, \dots, b\). In practice, \(V_{i+1}\) is unknown and must be replaced by estimated values \(\hat{V}_{i+1}\) at downstream nodes. The estimate \(\hat{\beta}_i\) then defines an estimate
\begin{equation} \label{eq:EstimatedContinuationValue}
	\hat{C}_i(x) = \hat{\beta}_i^\top \psi(x),
\end{equation}
of the continuation value at an arbitrary point \(x\) in the state space \(\R^d\). This, in turn, defines a procedure for estimating the option value.

\begin{algorithm}[H]
	\caption{Regression-Based Pricing Algorithm}\label{alg:GeneralRegressionPricing}
	\begin{algorithmic}[1]
		\State{} Simulate \(b\) independent paths \(\{X_{1j}, \dots, X_{mj}\}, j = 1, \dots, b\), of the Markov chain.
		\State{} At terminal nodes, set \(\hat{V}_{mj} = h_m(X_{mj})\) for \(j = 1, \dots, b\).
		\For{\(i = m-1\) \textbf{down to} \(1\)}
		\State{} Given estimated values \(\hat{V}_{i+1, j}, j = 1, \dots, b\), use regression to calculate \(\hat{\beta}_i = \hat{B}_\psi^{-1} \hat{B}_{\psi V}\).
		\State{} Compute estimated continuation values \(\hat{C}_i(X_{ij})\) using \Cref{eq:EstimatedContinuationValue}.
		\For{\(j = 1, \dots, b\)}
		\State{} Set value based on immediate exercise versus estimated continuation:
		\begin{equation} \label{eq:GeneralRegressionUpdate}
			\hat{V}_{ij} = \max \left\{ h_i(X_{ij}), \hat{C}_i(X_{ij}) \right\}.
		\end{equation}
		\EndFor{}
		\EndFor{}
		\State{} Set \(\hat{V}_0 = (\hat{V}_{11} + \cdots + \hat{V}_{1b}) / b\).
	\end{algorithmic}
\end{algorithm}

\begin{defn}{American Contingent Claim}{AmericanContingentClaim}
	An \textit{American contingent claim} \(X^a = (X, \mathcal{T}_{[0,T]})\) expiring at time \(T\) consists of a sequence of payoffs \((X_t)_{t=0}^T\), where \(X_t\) is an \(\cF_t\)-measurable random variable for \(t = 0, 1, \dots, T\), and of the set \(\mathcal{T}_{[0,T]}\) of admissible exercise policies.
\end{defn}

\begin{prop}{Arbitrage Price of American Put}{ArbitragePriceAmericanPut}
	The arbitrage price \(P_t^a\) of an American put option equals
	\begin{equation} \label{eq:AmericanPutArbitragePrice}
		P_t^a = \max_{\tau \in \mathcal{T}_{[t,T]}} \E_{\Pm^*} \left( \hat{r}^{-(\tau-t)} (K - S_\tau)^+ \mid \cF_t \right), \quad \forall t \le T.
	\end{equation}
	Moreover, for any \(t \le T\) the stopping time \(\tau_t^*\) that realizes the maximum in \Cref{eq:AmericanPutArbitragePrice} is given by the expression (by convention \(\min \emptyset = T\))
	\begin{equation} \label{eq:OptimalStoppingTime}
		\tau_t^* = \min \left\{ u \in \{t, t + 1, \dots, T\} \mid (K - S_u)^+ \ge P_u^a \right\}.
	\end{equation}
\end{prop}

\begin{prop}{Arbitrage Price of American Claim in CRR}{ArbitragePriceAmericanClaimCRR}
	For every \(t \le T\), the arbitrage price \(\pi(X^a)\) of an arbitrary American claim \(X^a\) in the CRR model equals
	\begin{equation} \label{eq:GeneralAmericanArbitragePrice}
		\pi_t(X^a) = \max_{\tau \in \mathcal{T}_{[t,T]}} \E_{\Pm^*} \left( \hat{r}^{-(\tau-t)} X_\tau \mid \cF_t \right).
	\end{equation}
	The price process \(\pi(X^a)\) can be determined using the following recurrence relation, for every \(t \le T - 1\),
	\begin{equation} \label{eq:SnellEnvelopeRecurrence}
		\pi_t(X^a) = \max \left\{ X_t, \E_{\Pm^*} \left( \hat{r}^{-1} \pi_{t+1}(X^a) \mid \cF_t \right) \right\}
	\end{equation}
	subject to the terminal condition \(\pi_T(X^a) = X_T\). In the case of a path-independent American claim \(X^a\) with the reward function \(h\) we have that, for every \(t \le T - 1\),
	\begin{equation} \label{eq:PathIndependentRecurrence}
		\pi_t(X^a) = \max \left\{ h(S_t, t), \hat{r}^{-1} \left( p_* \pi_{t+1}^u(X^a) + (1 - p_*) \pi_{t+1}^d(X^a) \right) \right\},
	\end{equation}
	where, for a generic stock price value \(S_t\) at time \(t\), we write \(\pi_{t+1}^u(X^a)\) and \(\pi_{t+1}^d(X^a)\) to denote the values of the price process \(\pi(X^a)\) at time \(t + 1\) in the nodes that correspond to the upward and downward movements of the stock price during the time-period \([t, t + 1]\), that is, for the values \(uS_t\) and \(dS_t\) of the stock price at time \(t + 1\), respectively.
\end{prop}
